\documentclass[12pt]{article}
\usepackage[margin=0.9in]{geometry}
\geometry{letterpaper}
\usepackage[T1]{fontenc} 
\usepackage[utf8]{inputenc} 
\usepackage[french]{babel}
\usepackage{graphicx} % Support for including images
\usepackage[urlcolor = blue]{hyperref} % Support for hyperlinks
%\hypersetup{hidelinks = False}
\usepackage{amsmath}
\usepackage{amsfonts}
\usepackage{siunitx}
\usepackage{xcolor}
%\sisetup{
%    round-mode = figures,
%    round-precision    = 3           ,
%    scientific-notation = engineering
%    }
\usepackage{float}
\usepackage{cancel}
\usepackage{tikz}
\usepackage{booktabs}
\usepackage{comment} 
\usepackage{minted}
\usepackage{tabularx}
\usepackage{placeins}
\renewcommand{\thesection}{\arabic{section}}
\newcommand\pder[2][]{\ensuremath{\frac{\partial#1}{\partial#2}}} 
\newcommand\pdder[2][]{\ensuremath{\frac{\partial^2 #1}{\partial#2^2}}} 
\usepackage{bibentry}
\nobibliography*
%------------------------------------------------------------------
% TITLE
%------------------------------------------------------------------
\title{

\centerline{
\begin{tabular}{lcr}
    \includegraphics[width=0.4\textwidth]{images/poly}
 &  & \includegraphics[width=0.5\textwidth]{images/McGillLogo.png}
\end{tabular}
}
\vspace{0.2 cm}
Laboratory Report\\
\vspace{0.2cm}
MECH 656\\
\vspace{0.2cm}
Fundamentals of Turbulent Flows 
  }

\author{
    \begin{minipage}{0.6\textwidth}
        \begin{center}
            Gaétan \textsc{Raynaud} \\
            2022063 - \texttt{gaetan.raynaud@polymtl.ca} \\
            260964987 - \texttt{gaetan.raynaud@mail.mcgill.ca}
        \end{center}
    \end{minipage}
}    

\date{\today}


%------------------------------------------------------------------
% DOCUMENT START HERE
%------------------------------------------------------------------
\begin{document}
\maketitle
\tableofcontents
\newpage

\section{Introduction}

\section{Theoretical Background}

\section{Experimental equipment and procedure}

\subsection{Anemometer calibration}
\begin{figure}
    \centering
    %% Creator: Matplotlib, PGF backend
%%
%% To include the figure in your LaTeX document, write
%%   \input{<filename>.pgf}
%%
%% Make sure the required packages are loaded in your preamble
%%   \usepackage{pgf}
%%
%% Figures using additional raster images can only be included by \input if
%% they are in the same directory as the main LaTeX file. For loading figures
%% from other directories you can use the `import` package
%%   \usepackage{import}
%% and then include the figures with
%%   \import{<path to file>}{<filename>.pgf}
%%
%% Matplotlib used the following preamble
%%
\begingroup%
\makeatletter%
\begin{pgfpicture}%
\pgfpathrectangle{\pgfpointorigin}{\pgfqpoint{6.400000in}{4.800000in}}%
\pgfusepath{use as bounding box, clip}%
\begin{pgfscope}%
\pgfsetbuttcap%
\pgfsetmiterjoin%
\definecolor{currentfill}{rgb}{1.000000,1.000000,1.000000}%
\pgfsetfillcolor{currentfill}%
\pgfsetlinewidth{0.000000pt}%
\definecolor{currentstroke}{rgb}{1.000000,1.000000,1.000000}%
\pgfsetstrokecolor{currentstroke}%
\pgfsetdash{}{0pt}%
\pgfpathmoveto{\pgfqpoint{0.000000in}{0.000000in}}%
\pgfpathlineto{\pgfqpoint{6.400000in}{0.000000in}}%
\pgfpathlineto{\pgfqpoint{6.400000in}{4.800000in}}%
\pgfpathlineto{\pgfqpoint{0.000000in}{4.800000in}}%
\pgfpathclose%
\pgfusepath{fill}%
\end{pgfscope}%
\begin{pgfscope}%
\pgfsetbuttcap%
\pgfsetmiterjoin%
\definecolor{currentfill}{rgb}{1.000000,1.000000,1.000000}%
\pgfsetfillcolor{currentfill}%
\pgfsetlinewidth{0.000000pt}%
\definecolor{currentstroke}{rgb}{0.000000,0.000000,0.000000}%
\pgfsetstrokecolor{currentstroke}%
\pgfsetstrokeopacity{0.000000}%
\pgfsetdash{}{0pt}%
\pgfpathmoveto{\pgfqpoint{0.800000in}{0.528000in}}%
\pgfpathlineto{\pgfqpoint{5.760000in}{0.528000in}}%
\pgfpathlineto{\pgfqpoint{5.760000in}{4.224000in}}%
\pgfpathlineto{\pgfqpoint{0.800000in}{4.224000in}}%
\pgfpathclose%
\pgfusepath{fill}%
\end{pgfscope}%
\begin{pgfscope}%
\pgfsetbuttcap%
\pgfsetroundjoin%
\definecolor{currentfill}{rgb}{0.000000,0.000000,0.000000}%
\pgfsetfillcolor{currentfill}%
\pgfsetlinewidth{0.803000pt}%
\definecolor{currentstroke}{rgb}{0.000000,0.000000,0.000000}%
\pgfsetstrokecolor{currentstroke}%
\pgfsetdash{}{0pt}%
\pgfsys@defobject{currentmarker}{\pgfqpoint{0.000000in}{-0.048611in}}{\pgfqpoint{0.000000in}{0.000000in}}{%
\pgfpathmoveto{\pgfqpoint{0.000000in}{0.000000in}}%
\pgfpathlineto{\pgfqpoint{0.000000in}{-0.048611in}}%
\pgfusepath{stroke,fill}%
}%
\begin{pgfscope}%
\pgfsys@transformshift{1.168327in}{0.528000in}%
\pgfsys@useobject{currentmarker}{}%
\end{pgfscope}%
\end{pgfscope}%
\begin{pgfscope}%
\definecolor{textcolor}{rgb}{0.000000,0.000000,0.000000}%
\pgfsetstrokecolor{textcolor}%
\pgfsetfillcolor{textcolor}%
\pgftext[x=1.168327in,y=0.430778in,,top]{\color{textcolor}\rmfamily\fontsize{10.000000}{12.000000}\selectfont \(\displaystyle 2.5\)}%
\end{pgfscope}%
\begin{pgfscope}%
\pgfsetbuttcap%
\pgfsetroundjoin%
\definecolor{currentfill}{rgb}{0.000000,0.000000,0.000000}%
\pgfsetfillcolor{currentfill}%
\pgfsetlinewidth{0.803000pt}%
\definecolor{currentstroke}{rgb}{0.000000,0.000000,0.000000}%
\pgfsetstrokecolor{currentstroke}%
\pgfsetdash{}{0pt}%
\pgfsys@defobject{currentmarker}{\pgfqpoint{0.000000in}{-0.048611in}}{\pgfqpoint{0.000000in}{0.000000in}}{%
\pgfpathmoveto{\pgfqpoint{0.000000in}{0.000000in}}%
\pgfpathlineto{\pgfqpoint{0.000000in}{-0.048611in}}%
\pgfusepath{stroke,fill}%
}%
\begin{pgfscope}%
\pgfsys@transformshift{1.766500in}{0.528000in}%
\pgfsys@useobject{currentmarker}{}%
\end{pgfscope}%
\end{pgfscope}%
\begin{pgfscope}%
\definecolor{textcolor}{rgb}{0.000000,0.000000,0.000000}%
\pgfsetstrokecolor{textcolor}%
\pgfsetfillcolor{textcolor}%
\pgftext[x=1.766500in,y=0.430778in,,top]{\color{textcolor}\rmfamily\fontsize{10.000000}{12.000000}\selectfont \(\displaystyle 5.0\)}%
\end{pgfscope}%
\begin{pgfscope}%
\pgfsetbuttcap%
\pgfsetroundjoin%
\definecolor{currentfill}{rgb}{0.000000,0.000000,0.000000}%
\pgfsetfillcolor{currentfill}%
\pgfsetlinewidth{0.803000pt}%
\definecolor{currentstroke}{rgb}{0.000000,0.000000,0.000000}%
\pgfsetstrokecolor{currentstroke}%
\pgfsetdash{}{0pt}%
\pgfsys@defobject{currentmarker}{\pgfqpoint{0.000000in}{-0.048611in}}{\pgfqpoint{0.000000in}{0.000000in}}{%
\pgfpathmoveto{\pgfqpoint{0.000000in}{0.000000in}}%
\pgfpathlineto{\pgfqpoint{0.000000in}{-0.048611in}}%
\pgfusepath{stroke,fill}%
}%
\begin{pgfscope}%
\pgfsys@transformshift{2.364673in}{0.528000in}%
\pgfsys@useobject{currentmarker}{}%
\end{pgfscope}%
\end{pgfscope}%
\begin{pgfscope}%
\definecolor{textcolor}{rgb}{0.000000,0.000000,0.000000}%
\pgfsetstrokecolor{textcolor}%
\pgfsetfillcolor{textcolor}%
\pgftext[x=2.364673in,y=0.430778in,,top]{\color{textcolor}\rmfamily\fontsize{10.000000}{12.000000}\selectfont \(\displaystyle 7.5\)}%
\end{pgfscope}%
\begin{pgfscope}%
\pgfsetbuttcap%
\pgfsetroundjoin%
\definecolor{currentfill}{rgb}{0.000000,0.000000,0.000000}%
\pgfsetfillcolor{currentfill}%
\pgfsetlinewidth{0.803000pt}%
\definecolor{currentstroke}{rgb}{0.000000,0.000000,0.000000}%
\pgfsetstrokecolor{currentstroke}%
\pgfsetdash{}{0pt}%
\pgfsys@defobject{currentmarker}{\pgfqpoint{0.000000in}{-0.048611in}}{\pgfqpoint{0.000000in}{0.000000in}}{%
\pgfpathmoveto{\pgfqpoint{0.000000in}{0.000000in}}%
\pgfpathlineto{\pgfqpoint{0.000000in}{-0.048611in}}%
\pgfusepath{stroke,fill}%
}%
\begin{pgfscope}%
\pgfsys@transformshift{2.962846in}{0.528000in}%
\pgfsys@useobject{currentmarker}{}%
\end{pgfscope}%
\end{pgfscope}%
\begin{pgfscope}%
\definecolor{textcolor}{rgb}{0.000000,0.000000,0.000000}%
\pgfsetstrokecolor{textcolor}%
\pgfsetfillcolor{textcolor}%
\pgftext[x=2.962846in,y=0.430778in,,top]{\color{textcolor}\rmfamily\fontsize{10.000000}{12.000000}\selectfont \(\displaystyle 10.0\)}%
\end{pgfscope}%
\begin{pgfscope}%
\pgfsetbuttcap%
\pgfsetroundjoin%
\definecolor{currentfill}{rgb}{0.000000,0.000000,0.000000}%
\pgfsetfillcolor{currentfill}%
\pgfsetlinewidth{0.803000pt}%
\definecolor{currentstroke}{rgb}{0.000000,0.000000,0.000000}%
\pgfsetstrokecolor{currentstroke}%
\pgfsetdash{}{0pt}%
\pgfsys@defobject{currentmarker}{\pgfqpoint{0.000000in}{-0.048611in}}{\pgfqpoint{0.000000in}{0.000000in}}{%
\pgfpathmoveto{\pgfqpoint{0.000000in}{0.000000in}}%
\pgfpathlineto{\pgfqpoint{0.000000in}{-0.048611in}}%
\pgfusepath{stroke,fill}%
}%
\begin{pgfscope}%
\pgfsys@transformshift{3.561019in}{0.528000in}%
\pgfsys@useobject{currentmarker}{}%
\end{pgfscope}%
\end{pgfscope}%
\begin{pgfscope}%
\definecolor{textcolor}{rgb}{0.000000,0.000000,0.000000}%
\pgfsetstrokecolor{textcolor}%
\pgfsetfillcolor{textcolor}%
\pgftext[x=3.561019in,y=0.430778in,,top]{\color{textcolor}\rmfamily\fontsize{10.000000}{12.000000}\selectfont \(\displaystyle 12.5\)}%
\end{pgfscope}%
\begin{pgfscope}%
\pgfsetbuttcap%
\pgfsetroundjoin%
\definecolor{currentfill}{rgb}{0.000000,0.000000,0.000000}%
\pgfsetfillcolor{currentfill}%
\pgfsetlinewidth{0.803000pt}%
\definecolor{currentstroke}{rgb}{0.000000,0.000000,0.000000}%
\pgfsetstrokecolor{currentstroke}%
\pgfsetdash{}{0pt}%
\pgfsys@defobject{currentmarker}{\pgfqpoint{0.000000in}{-0.048611in}}{\pgfqpoint{0.000000in}{0.000000in}}{%
\pgfpathmoveto{\pgfqpoint{0.000000in}{0.000000in}}%
\pgfpathlineto{\pgfqpoint{0.000000in}{-0.048611in}}%
\pgfusepath{stroke,fill}%
}%
\begin{pgfscope}%
\pgfsys@transformshift{4.159192in}{0.528000in}%
\pgfsys@useobject{currentmarker}{}%
\end{pgfscope}%
\end{pgfscope}%
\begin{pgfscope}%
\definecolor{textcolor}{rgb}{0.000000,0.000000,0.000000}%
\pgfsetstrokecolor{textcolor}%
\pgfsetfillcolor{textcolor}%
\pgftext[x=4.159192in,y=0.430778in,,top]{\color{textcolor}\rmfamily\fontsize{10.000000}{12.000000}\selectfont \(\displaystyle 15.0\)}%
\end{pgfscope}%
\begin{pgfscope}%
\pgfsetbuttcap%
\pgfsetroundjoin%
\definecolor{currentfill}{rgb}{0.000000,0.000000,0.000000}%
\pgfsetfillcolor{currentfill}%
\pgfsetlinewidth{0.803000pt}%
\definecolor{currentstroke}{rgb}{0.000000,0.000000,0.000000}%
\pgfsetstrokecolor{currentstroke}%
\pgfsetdash{}{0pt}%
\pgfsys@defobject{currentmarker}{\pgfqpoint{0.000000in}{-0.048611in}}{\pgfqpoint{0.000000in}{0.000000in}}{%
\pgfpathmoveto{\pgfqpoint{0.000000in}{0.000000in}}%
\pgfpathlineto{\pgfqpoint{0.000000in}{-0.048611in}}%
\pgfusepath{stroke,fill}%
}%
\begin{pgfscope}%
\pgfsys@transformshift{4.757365in}{0.528000in}%
\pgfsys@useobject{currentmarker}{}%
\end{pgfscope}%
\end{pgfscope}%
\begin{pgfscope}%
\definecolor{textcolor}{rgb}{0.000000,0.000000,0.000000}%
\pgfsetstrokecolor{textcolor}%
\pgfsetfillcolor{textcolor}%
\pgftext[x=4.757365in,y=0.430778in,,top]{\color{textcolor}\rmfamily\fontsize{10.000000}{12.000000}\selectfont \(\displaystyle 17.5\)}%
\end{pgfscope}%
\begin{pgfscope}%
\pgfsetbuttcap%
\pgfsetroundjoin%
\definecolor{currentfill}{rgb}{0.000000,0.000000,0.000000}%
\pgfsetfillcolor{currentfill}%
\pgfsetlinewidth{0.803000pt}%
\definecolor{currentstroke}{rgb}{0.000000,0.000000,0.000000}%
\pgfsetstrokecolor{currentstroke}%
\pgfsetdash{}{0pt}%
\pgfsys@defobject{currentmarker}{\pgfqpoint{0.000000in}{-0.048611in}}{\pgfqpoint{0.000000in}{0.000000in}}{%
\pgfpathmoveto{\pgfqpoint{0.000000in}{0.000000in}}%
\pgfpathlineto{\pgfqpoint{0.000000in}{-0.048611in}}%
\pgfusepath{stroke,fill}%
}%
\begin{pgfscope}%
\pgfsys@transformshift{5.355539in}{0.528000in}%
\pgfsys@useobject{currentmarker}{}%
\end{pgfscope}%
\end{pgfscope}%
\begin{pgfscope}%
\definecolor{textcolor}{rgb}{0.000000,0.000000,0.000000}%
\pgfsetstrokecolor{textcolor}%
\pgfsetfillcolor{textcolor}%
\pgftext[x=5.355539in,y=0.430778in,,top]{\color{textcolor}\rmfamily\fontsize{10.000000}{12.000000}\selectfont \(\displaystyle 20.0\)}%
\end{pgfscope}%
\begin{pgfscope}%
\definecolor{textcolor}{rgb}{0.000000,0.000000,0.000000}%
\pgfsetstrokecolor{textcolor}%
\pgfsetfillcolor{textcolor}%
\pgftext[x=3.280000in,y=0.251766in,,top]{\color{textcolor}\rmfamily\fontsize{10.000000}{12.000000}\selectfont Velocity \(\displaystyle U\)}%
\end{pgfscope}%
\begin{pgfscope}%
\pgfsetbuttcap%
\pgfsetroundjoin%
\definecolor{currentfill}{rgb}{0.000000,0.000000,0.000000}%
\pgfsetfillcolor{currentfill}%
\pgfsetlinewidth{0.803000pt}%
\definecolor{currentstroke}{rgb}{0.000000,0.000000,0.000000}%
\pgfsetstrokecolor{currentstroke}%
\pgfsetdash{}{0pt}%
\pgfsys@defobject{currentmarker}{\pgfqpoint{-0.048611in}{0.000000in}}{\pgfqpoint{0.000000in}{0.000000in}}{%
\pgfpathmoveto{\pgfqpoint{0.000000in}{0.000000in}}%
\pgfpathlineto{\pgfqpoint{-0.048611in}{0.000000in}}%
\pgfusepath{stroke,fill}%
}%
\begin{pgfscope}%
\pgfsys@transformshift{0.800000in}{1.014674in}%
\pgfsys@useobject{currentmarker}{}%
\end{pgfscope}%
\end{pgfscope}%
\begin{pgfscope}%
\definecolor{textcolor}{rgb}{0.000000,0.000000,0.000000}%
\pgfsetstrokecolor{textcolor}%
\pgfsetfillcolor{textcolor}%
\pgftext[x=0.525308in,y=0.966448in,left,base]{\color{textcolor}\rmfamily\fontsize{10.000000}{12.000000}\selectfont \(\displaystyle 1.2\)}%
\end{pgfscope}%
\begin{pgfscope}%
\pgfsetbuttcap%
\pgfsetroundjoin%
\definecolor{currentfill}{rgb}{0.000000,0.000000,0.000000}%
\pgfsetfillcolor{currentfill}%
\pgfsetlinewidth{0.803000pt}%
\definecolor{currentstroke}{rgb}{0.000000,0.000000,0.000000}%
\pgfsetstrokecolor{currentstroke}%
\pgfsetdash{}{0pt}%
\pgfsys@defobject{currentmarker}{\pgfqpoint{-0.048611in}{0.000000in}}{\pgfqpoint{0.000000in}{0.000000in}}{%
\pgfpathmoveto{\pgfqpoint{0.000000in}{0.000000in}}%
\pgfpathlineto{\pgfqpoint{-0.048611in}{0.000000in}}%
\pgfusepath{stroke,fill}%
}%
\begin{pgfscope}%
\pgfsys@transformshift{0.800000in}{1.579440in}%
\pgfsys@useobject{currentmarker}{}%
\end{pgfscope}%
\end{pgfscope}%
\begin{pgfscope}%
\definecolor{textcolor}{rgb}{0.000000,0.000000,0.000000}%
\pgfsetstrokecolor{textcolor}%
\pgfsetfillcolor{textcolor}%
\pgftext[x=0.525308in,y=1.531215in,left,base]{\color{textcolor}\rmfamily\fontsize{10.000000}{12.000000}\selectfont \(\displaystyle 1.3\)}%
\end{pgfscope}%
\begin{pgfscope}%
\pgfsetbuttcap%
\pgfsetroundjoin%
\definecolor{currentfill}{rgb}{0.000000,0.000000,0.000000}%
\pgfsetfillcolor{currentfill}%
\pgfsetlinewidth{0.803000pt}%
\definecolor{currentstroke}{rgb}{0.000000,0.000000,0.000000}%
\pgfsetstrokecolor{currentstroke}%
\pgfsetdash{}{0pt}%
\pgfsys@defobject{currentmarker}{\pgfqpoint{-0.048611in}{0.000000in}}{\pgfqpoint{0.000000in}{0.000000in}}{%
\pgfpathmoveto{\pgfqpoint{0.000000in}{0.000000in}}%
\pgfpathlineto{\pgfqpoint{-0.048611in}{0.000000in}}%
\pgfusepath{stroke,fill}%
}%
\begin{pgfscope}%
\pgfsys@transformshift{0.800000in}{2.144207in}%
\pgfsys@useobject{currentmarker}{}%
\end{pgfscope}%
\end{pgfscope}%
\begin{pgfscope}%
\definecolor{textcolor}{rgb}{0.000000,0.000000,0.000000}%
\pgfsetstrokecolor{textcolor}%
\pgfsetfillcolor{textcolor}%
\pgftext[x=0.525308in,y=2.095982in,left,base]{\color{textcolor}\rmfamily\fontsize{10.000000}{12.000000}\selectfont \(\displaystyle 1.4\)}%
\end{pgfscope}%
\begin{pgfscope}%
\pgfsetbuttcap%
\pgfsetroundjoin%
\definecolor{currentfill}{rgb}{0.000000,0.000000,0.000000}%
\pgfsetfillcolor{currentfill}%
\pgfsetlinewidth{0.803000pt}%
\definecolor{currentstroke}{rgb}{0.000000,0.000000,0.000000}%
\pgfsetstrokecolor{currentstroke}%
\pgfsetdash{}{0pt}%
\pgfsys@defobject{currentmarker}{\pgfqpoint{-0.048611in}{0.000000in}}{\pgfqpoint{0.000000in}{0.000000in}}{%
\pgfpathmoveto{\pgfqpoint{0.000000in}{0.000000in}}%
\pgfpathlineto{\pgfqpoint{-0.048611in}{0.000000in}}%
\pgfusepath{stroke,fill}%
}%
\begin{pgfscope}%
\pgfsys@transformshift{0.800000in}{2.708974in}%
\pgfsys@useobject{currentmarker}{}%
\end{pgfscope}%
\end{pgfscope}%
\begin{pgfscope}%
\definecolor{textcolor}{rgb}{0.000000,0.000000,0.000000}%
\pgfsetstrokecolor{textcolor}%
\pgfsetfillcolor{textcolor}%
\pgftext[x=0.525308in,y=2.660749in,left,base]{\color{textcolor}\rmfamily\fontsize{10.000000}{12.000000}\selectfont \(\displaystyle 1.5\)}%
\end{pgfscope}%
\begin{pgfscope}%
\pgfsetbuttcap%
\pgfsetroundjoin%
\definecolor{currentfill}{rgb}{0.000000,0.000000,0.000000}%
\pgfsetfillcolor{currentfill}%
\pgfsetlinewidth{0.803000pt}%
\definecolor{currentstroke}{rgb}{0.000000,0.000000,0.000000}%
\pgfsetstrokecolor{currentstroke}%
\pgfsetdash{}{0pt}%
\pgfsys@defobject{currentmarker}{\pgfqpoint{-0.048611in}{0.000000in}}{\pgfqpoint{0.000000in}{0.000000in}}{%
\pgfpathmoveto{\pgfqpoint{0.000000in}{0.000000in}}%
\pgfpathlineto{\pgfqpoint{-0.048611in}{0.000000in}}%
\pgfusepath{stroke,fill}%
}%
\begin{pgfscope}%
\pgfsys@transformshift{0.800000in}{3.273741in}%
\pgfsys@useobject{currentmarker}{}%
\end{pgfscope}%
\end{pgfscope}%
\begin{pgfscope}%
\definecolor{textcolor}{rgb}{0.000000,0.000000,0.000000}%
\pgfsetstrokecolor{textcolor}%
\pgfsetfillcolor{textcolor}%
\pgftext[x=0.525308in,y=3.225516in,left,base]{\color{textcolor}\rmfamily\fontsize{10.000000}{12.000000}\selectfont \(\displaystyle 1.6\)}%
\end{pgfscope}%
\begin{pgfscope}%
\pgfsetbuttcap%
\pgfsetroundjoin%
\definecolor{currentfill}{rgb}{0.000000,0.000000,0.000000}%
\pgfsetfillcolor{currentfill}%
\pgfsetlinewidth{0.803000pt}%
\definecolor{currentstroke}{rgb}{0.000000,0.000000,0.000000}%
\pgfsetstrokecolor{currentstroke}%
\pgfsetdash{}{0pt}%
\pgfsys@defobject{currentmarker}{\pgfqpoint{-0.048611in}{0.000000in}}{\pgfqpoint{0.000000in}{0.000000in}}{%
\pgfpathmoveto{\pgfqpoint{0.000000in}{0.000000in}}%
\pgfpathlineto{\pgfqpoint{-0.048611in}{0.000000in}}%
\pgfusepath{stroke,fill}%
}%
\begin{pgfscope}%
\pgfsys@transformshift{0.800000in}{3.838508in}%
\pgfsys@useobject{currentmarker}{}%
\end{pgfscope}%
\end{pgfscope}%
\begin{pgfscope}%
\definecolor{textcolor}{rgb}{0.000000,0.000000,0.000000}%
\pgfsetstrokecolor{textcolor}%
\pgfsetfillcolor{textcolor}%
\pgftext[x=0.525308in,y=3.790283in,left,base]{\color{textcolor}\rmfamily\fontsize{10.000000}{12.000000}\selectfont \(\displaystyle 1.7\)}%
\end{pgfscope}%
\begin{pgfscope}%
\definecolor{textcolor}{rgb}{0.000000,0.000000,0.000000}%
\pgfsetstrokecolor{textcolor}%
\pgfsetfillcolor{textcolor}%
\pgftext[x=0.469752in,y=2.376000in,,bottom,rotate=90.000000]{\color{textcolor}\rmfamily\fontsize{10.000000}{12.000000}\selectfont Square output voltage \(\displaystyle E^2\)}%
\end{pgfscope}%
\begin{pgfscope}%
\pgfpathrectangle{\pgfqpoint{0.800000in}{0.528000in}}{\pgfqpoint{4.960000in}{3.696000in}}%
\pgfusepath{clip}%
\pgfsetbuttcap%
\pgfsetmiterjoin%
\definecolor{currentfill}{rgb}{0.121569,0.466667,0.705882}%
\pgfsetfillcolor{currentfill}%
\pgfsetlinewidth{1.003750pt}%
\definecolor{currentstroke}{rgb}{0.121569,0.466667,0.705882}%
\pgfsetstrokecolor{currentstroke}%
\pgfsetdash{}{0pt}%
\pgfsys@defobject{currentmarker}{\pgfqpoint{-0.041667in}{-0.041667in}}{\pgfqpoint{0.041667in}{0.041667in}}{%
\pgfpathmoveto{\pgfqpoint{-0.041667in}{-0.041667in}}%
\pgfpathlineto{\pgfqpoint{0.041667in}{-0.041667in}}%
\pgfpathlineto{\pgfqpoint{0.041667in}{0.041667in}}%
\pgfpathlineto{\pgfqpoint{-0.041667in}{0.041667in}}%
\pgfpathclose%
\pgfusepath{stroke,fill}%
}%
\begin{pgfscope}%
\pgfsys@transformshift{1.025455in}{0.706537in}%
\pgfsys@useobject{currentmarker}{}%
\end{pgfscope}%
\begin{pgfscope}%
\pgfsys@transformshift{1.116162in}{0.844227in}%
\pgfsys@useobject{currentmarker}{}%
\end{pgfscope}%
\begin{pgfscope}%
\pgfsys@transformshift{1.193742in}{0.956390in}%
\pgfsys@useobject{currentmarker}{}%
\end{pgfscope}%
\begin{pgfscope}%
\pgfsys@transformshift{1.333916in}{1.142876in}%
\pgfsys@useobject{currentmarker}{}%
\end{pgfscope}%
\begin{pgfscope}%
\pgfsys@transformshift{1.451433in}{1.286383in}%
\pgfsys@useobject{currentmarker}{}%
\end{pgfscope}%
\begin{pgfscope}%
\pgfsys@transformshift{1.568167in}{1.421362in}%
\pgfsys@useobject{currentmarker}{}%
\end{pgfscope}%
\begin{pgfscope}%
\pgfsys@transformshift{1.639439in}{1.500599in}%
\pgfsys@useobject{currentmarker}{}%
\end{pgfscope}%
\begin{pgfscope}%
\pgfsys@transformshift{1.742207in}{1.609260in}%
\pgfsys@useobject{currentmarker}{}%
\end{pgfscope}%
\begin{pgfscope}%
\pgfsys@transformshift{1.820259in}{1.688610in}%
\pgfsys@useobject{currentmarker}{}%
\end{pgfscope}%
\begin{pgfscope}%
\pgfsys@transformshift{1.892270in}{1.760505in}%
\pgfsys@useobject{currentmarker}{}%
\end{pgfscope}%
\begin{pgfscope}%
\pgfsys@transformshift{1.964766in}{1.828390in}%
\pgfsys@useobject{currentmarker}{}%
\end{pgfscope}%
\begin{pgfscope}%
\pgfsys@transformshift{2.096706in}{1.949419in}%
\pgfsys@useobject{currentmarker}{}%
\end{pgfscope}%
\begin{pgfscope}%
\pgfsys@transformshift{2.221932in}{2.058645in}%
\pgfsys@useobject{currentmarker}{}%
\end{pgfscope}%
\begin{pgfscope}%
\pgfsys@transformshift{2.330003in}{2.149799in}%
\pgfsys@useobject{currentmarker}{}%
\end{pgfscope}%
\begin{pgfscope}%
\pgfsys@transformshift{2.442864in}{2.239710in}%
\pgfsys@useobject{currentmarker}{}%
\end{pgfscope}%
\begin{pgfscope}%
\pgfsys@transformshift{2.542673in}{2.318495in}%
\pgfsys@useobject{currentmarker}{}%
\end{pgfscope}%
\begin{pgfscope}%
\pgfsys@transformshift{2.769696in}{2.487473in}%
\pgfsys@useobject{currentmarker}{}%
\end{pgfscope}%
\begin{pgfscope}%
\pgfsys@transformshift{2.984402in}{2.638717in}%
\pgfsys@useobject{currentmarker}{}%
\end{pgfscope}%
\begin{pgfscope}%
\pgfsys@transformshift{3.346955in}{2.877840in}%
\pgfsys@useobject{currentmarker}{}%
\end{pgfscope}%
\begin{pgfscope}%
\pgfsys@transformshift{3.677527in}{3.080083in}%
\pgfsys@useobject{currentmarker}{}%
\end{pgfscope}%
\begin{pgfscope}%
\pgfsys@transformshift{3.961470in}{3.244656in}%
\pgfsys@useobject{currentmarker}{}%
\end{pgfscope}%
\begin{pgfscope}%
\pgfsys@transformshift{4.232459in}{3.397595in}%
\pgfsys@useobject{currentmarker}{}%
\end{pgfscope}%
\begin{pgfscope}%
\pgfsys@transformshift{4.489913in}{3.534212in}%
\pgfsys@useobject{currentmarker}{}%
\end{pgfscope}%
\begin{pgfscope}%
\pgfsys@transformshift{4.712684in}{3.650497in}%
\pgfsys@useobject{currentmarker}{}%
\end{pgfscope}%
\begin{pgfscope}%
\pgfsys@transformshift{4.951839in}{3.770171in}%
\pgfsys@useobject{currentmarker}{}%
\end{pgfscope}%
\begin{pgfscope}%
\pgfsys@transformshift{5.141232in}{3.869006in}%
\pgfsys@useobject{currentmarker}{}%
\end{pgfscope}%
\begin{pgfscope}%
\pgfsys@transformshift{5.340917in}{3.965807in}%
\pgfsys@useobject{currentmarker}{}%
\end{pgfscope}%
\begin{pgfscope}%
\pgfsys@transformshift{5.534545in}{4.056000in}%
\pgfsys@useobject{currentmarker}{}%
\end{pgfscope}%
\end{pgfscope}%
\begin{pgfscope}%
\pgfpathrectangle{\pgfqpoint{0.800000in}{0.528000in}}{\pgfqpoint{4.960000in}{3.696000in}}%
\pgfusepath{clip}%
\pgfsetrectcap%
\pgfsetroundjoin%
\pgfsetlinewidth{1.505625pt}%
\definecolor{currentstroke}{rgb}{1.000000,0.498039,0.054902}%
\pgfsetstrokecolor{currentstroke}%
\pgfsetdash{}{0pt}%
\pgfpathmoveto{\pgfqpoint{1.025455in}{0.696000in}}%
\pgfpathlineto{\pgfqpoint{1.075104in}{0.778088in}}%
\pgfpathlineto{\pgfqpoint{1.124754in}{0.855906in}}%
\pgfpathlineto{\pgfqpoint{1.178917in}{0.936603in}}%
\pgfpathlineto{\pgfqpoint{1.237594in}{1.019746in}}%
\pgfpathlineto{\pgfqpoint{1.296271in}{1.099013in}}%
\pgfpathlineto{\pgfqpoint{1.359461in}{1.180584in}}%
\pgfpathlineto{\pgfqpoint{1.427165in}{1.264143in}}%
\pgfpathlineto{\pgfqpoint{1.494869in}{1.344194in}}%
\pgfpathlineto{\pgfqpoint{1.567087in}{1.426152in}}%
\pgfpathlineto{\pgfqpoint{1.643818in}{1.509779in}}%
\pgfpathlineto{\pgfqpoint{1.725063in}{1.594865in}}%
\pgfpathlineto{\pgfqpoint{1.806308in}{1.676764in}}%
\pgfpathlineto{\pgfqpoint{1.892067in}{1.760105in}}%
\pgfpathlineto{\pgfqpoint{1.982339in}{1.844724in}}%
\pgfpathlineto{\pgfqpoint{2.077124in}{1.930474in}}%
\pgfpathlineto{\pgfqpoint{2.176424in}{2.017226in}}%
\pgfpathlineto{\pgfqpoint{2.280237in}{2.104866in}}%
\pgfpathlineto{\pgfqpoint{2.388563in}{2.193294in}}%
\pgfpathlineto{\pgfqpoint{2.501403in}{2.282419in}}%
\pgfpathlineto{\pgfqpoint{2.618757in}{2.372162in}}%
\pgfpathlineto{\pgfqpoint{2.740624in}{2.462454in}}%
\pgfpathlineto{\pgfqpoint{2.867005in}{2.553232in}}%
\pgfpathlineto{\pgfqpoint{2.997900in}{2.644441in}}%
\pgfpathlineto{\pgfqpoint{3.133308in}{2.736031in}}%
\pgfpathlineto{\pgfqpoint{3.273230in}{2.827958in}}%
\pgfpathlineto{\pgfqpoint{3.417665in}{2.920183in}}%
\pgfpathlineto{\pgfqpoint{3.566614in}{3.012671in}}%
\pgfpathlineto{\pgfqpoint{3.720076in}{3.105390in}}%
\pgfpathlineto{\pgfqpoint{3.882566in}{3.200931in}}%
\pgfpathlineto{\pgfqpoint{4.049570in}{3.296514in}}%
\pgfpathlineto{\pgfqpoint{4.221087in}{3.392126in}}%
\pgfpathlineto{\pgfqpoint{4.401631in}{3.490174in}}%
\pgfpathlineto{\pgfqpoint{4.586689in}{3.588107in}}%
\pgfpathlineto{\pgfqpoint{4.776260in}{3.685924in}}%
\pgfpathlineto{\pgfqpoint{4.974858in}{3.785869in}}%
\pgfpathlineto{\pgfqpoint{5.177971in}{3.885589in}}%
\pgfpathlineto{\pgfqpoint{5.390110in}{3.987228in}}%
\pgfpathlineto{\pgfqpoint{5.534545in}{4.055048in}}%
\pgfpathlineto{\pgfqpoint{5.534545in}{4.055048in}}%
\pgfusepath{stroke}%
\end{pgfscope}%
\begin{pgfscope}%
\pgfsetrectcap%
\pgfsetmiterjoin%
\pgfsetlinewidth{0.803000pt}%
\definecolor{currentstroke}{rgb}{0.000000,0.000000,0.000000}%
\pgfsetstrokecolor{currentstroke}%
\pgfsetdash{}{0pt}%
\pgfpathmoveto{\pgfqpoint{0.800000in}{0.528000in}}%
\pgfpathlineto{\pgfqpoint{0.800000in}{4.224000in}}%
\pgfusepath{stroke}%
\end{pgfscope}%
\begin{pgfscope}%
\pgfsetrectcap%
\pgfsetmiterjoin%
\pgfsetlinewidth{0.803000pt}%
\definecolor{currentstroke}{rgb}{0.000000,0.000000,0.000000}%
\pgfsetstrokecolor{currentstroke}%
\pgfsetdash{}{0pt}%
\pgfpathmoveto{\pgfqpoint{5.760000in}{0.528000in}}%
\pgfpathlineto{\pgfqpoint{5.760000in}{4.224000in}}%
\pgfusepath{stroke}%
\end{pgfscope}%
\begin{pgfscope}%
\pgfsetrectcap%
\pgfsetmiterjoin%
\pgfsetlinewidth{0.803000pt}%
\definecolor{currentstroke}{rgb}{0.000000,0.000000,0.000000}%
\pgfsetstrokecolor{currentstroke}%
\pgfsetdash{}{0pt}%
\pgfpathmoveto{\pgfqpoint{0.800000in}{0.528000in}}%
\pgfpathlineto{\pgfqpoint{5.760000in}{0.528000in}}%
\pgfusepath{stroke}%
\end{pgfscope}%
\begin{pgfscope}%
\pgfsetrectcap%
\pgfsetmiterjoin%
\pgfsetlinewidth{0.803000pt}%
\definecolor{currentstroke}{rgb}{0.000000,0.000000,0.000000}%
\pgfsetstrokecolor{currentstroke}%
\pgfsetdash{}{0pt}%
\pgfpathmoveto{\pgfqpoint{0.800000in}{4.224000in}}%
\pgfpathlineto{\pgfqpoint{5.760000in}{4.224000in}}%
\pgfusepath{stroke}%
\end{pgfscope}%
\begin{pgfscope}%
\pgfsetbuttcap%
\pgfsetmiterjoin%
\definecolor{currentfill}{rgb}{1.000000,1.000000,1.000000}%
\pgfsetfillcolor{currentfill}%
\pgfsetfillopacity{0.800000}%
\pgfsetlinewidth{1.003750pt}%
\definecolor{currentstroke}{rgb}{0.800000,0.800000,0.800000}%
\pgfsetstrokecolor{currentstroke}%
\pgfsetstrokeopacity{0.800000}%
\pgfsetdash{}{0pt}%
\pgfpathmoveto{\pgfqpoint{0.897222in}{3.725543in}}%
\pgfpathlineto{\pgfqpoint{2.722071in}{3.725543in}}%
\pgfpathquadraticcurveto{\pgfqpoint{2.749849in}{3.725543in}}{\pgfqpoint{2.749849in}{3.753321in}}%
\pgfpathlineto{\pgfqpoint{2.749849in}{4.126778in}}%
\pgfpathquadraticcurveto{\pgfqpoint{2.749849in}{4.154556in}}{\pgfqpoint{2.722071in}{4.154556in}}%
\pgfpathlineto{\pgfqpoint{0.897222in}{4.154556in}}%
\pgfpathquadraticcurveto{\pgfqpoint{0.869444in}{4.154556in}}{\pgfqpoint{0.869444in}{4.126778in}}%
\pgfpathlineto{\pgfqpoint{0.869444in}{3.753321in}}%
\pgfpathquadraticcurveto{\pgfqpoint{0.869444in}{3.725543in}}{\pgfqpoint{0.897222in}{3.725543in}}%
\pgfpathclose%
\pgfusepath{stroke,fill}%
\end{pgfscope}%
\begin{pgfscope}%
\pgfsetbuttcap%
\pgfsetmiterjoin%
\definecolor{currentfill}{rgb}{0.121569,0.466667,0.705882}%
\pgfsetfillcolor{currentfill}%
\pgfsetlinewidth{1.003750pt}%
\definecolor{currentstroke}{rgb}{0.121569,0.466667,0.705882}%
\pgfsetstrokecolor{currentstroke}%
\pgfsetdash{}{0pt}%
\pgfsys@defobject{currentmarker}{\pgfqpoint{-0.041667in}{-0.041667in}}{\pgfqpoint{0.041667in}{0.041667in}}{%
\pgfpathmoveto{\pgfqpoint{-0.041667in}{-0.041667in}}%
\pgfpathlineto{\pgfqpoint{0.041667in}{-0.041667in}}%
\pgfpathlineto{\pgfqpoint{0.041667in}{0.041667in}}%
\pgfpathlineto{\pgfqpoint{-0.041667in}{0.041667in}}%
\pgfpathclose%
\pgfusepath{stroke,fill}%
}%
\begin{pgfscope}%
\pgfsys@transformshift{1.063889in}{4.050389in}%
\pgfsys@useobject{currentmarker}{}%
\end{pgfscope}%
\end{pgfscope}%
\begin{pgfscope}%
\definecolor{textcolor}{rgb}{0.000000,0.000000,0.000000}%
\pgfsetstrokecolor{textcolor}%
\pgfsetfillcolor{textcolor}%
\pgftext[x=1.313889in,y=4.001778in,left,base]{\color{textcolor}\rmfamily\fontsize{10.000000}{12.000000}\selectfont Provided data}%
\end{pgfscope}%
\begin{pgfscope}%
\pgfsetrectcap%
\pgfsetroundjoin%
\pgfsetlinewidth{1.505625pt}%
\definecolor{currentstroke}{rgb}{1.000000,0.498039,0.054902}%
\pgfsetstrokecolor{currentstroke}%
\pgfsetdash{}{0pt}%
\pgfpathmoveto{\pgfqpoint{0.925000in}{3.856716in}}%
\pgfpathlineto{\pgfqpoint{1.202778in}{3.856716in}}%
\pgfusepath{stroke}%
\end{pgfscope}%
\begin{pgfscope}%
\definecolor{textcolor}{rgb}{0.000000,0.000000,0.000000}%
\pgfsetstrokecolor{textcolor}%
\pgfsetfillcolor{textcolor}%
\pgftext[x=1.313889in,y=3.808105in,left,base]{\color{textcolor}\rmfamily\fontsize{10.000000}{12.000000}\selectfont King's Law calibration}%
\end{pgfscope}%
\end{pgfpicture}%
\makeatother%
\endgroup%

    \caption{Calibration of King's Law coefficients $A = \num{8.44e-1}$ \si{V^2}, $B = \num{2.23e-1}$ \si{uSI} and $n=0.458$ with a L2-normalized residual rms $e_{rms} = \num{4.6e-4}$}
    \label{fig:AnemometerFitting}
\end{figure}

%% Rappeler les valeurs dans un tableau synthétoqie et expliquer l'expression de l'erreur e_rms


\section{Results}

\begin{figure}
    \centering
        \resizebox{0.8\linewidth}{!}{%% Creator: Matplotlib, PGF backend
%%
%% To include the figure in your LaTeX document, write
%%   \input{<filename>.pgf}
%%
%% Make sure the required packages are loaded in your preamble
%%   \usepackage{pgf}
%%
%% Figures using additional raster images can only be included by \input if
%% they are in the same directory as the main LaTeX file. For loading figures
%% from other directories you can use the `import` package
%%   \usepackage{import}
%% and then include the figures with
%%   \import{<path to file>}{<filename>.pgf}
%%
%% Matplotlib used the following preamble
%%   \usepackage{fontspec}
%%   \setmainfont{DejaVuSerif.ttf}[Path=C:/Users/garaya/AppData/Local/Continuum/anaconda3/lib/site-packages/matplotlib/mpl-data/fonts/ttf/]
%%   \setsansfont{DejaVuSans.ttf}[Path=C:/Users/garaya/AppData/Local/Continuum/anaconda3/lib/site-packages/matplotlib/mpl-data/fonts/ttf/]
%%   \setmonofont{DejaVuSansMono.ttf}[Path=C:/Users/garaya/AppData/Local/Continuum/anaconda3/lib/site-packages/matplotlib/mpl-data/fonts/ttf/]
%%
\begingroup%
\makeatletter%
\begin{pgfpicture}%
\pgfpathrectangle{\pgfpointorigin}{\pgfqpoint{6.400000in}{4.780000in}}%
\pgfusepath{use as bounding box, clip}%
\begin{pgfscope}%
\pgfsetbuttcap%
\pgfsetmiterjoin%
\definecolor{currentfill}{rgb}{1.000000,1.000000,1.000000}%
\pgfsetfillcolor{currentfill}%
\pgfsetlinewidth{0.000000pt}%
\definecolor{currentstroke}{rgb}{1.000000,1.000000,1.000000}%
\pgfsetstrokecolor{currentstroke}%
\pgfsetdash{}{0pt}%
\pgfpathmoveto{\pgfqpoint{0.000000in}{0.000000in}}%
\pgfpathlineto{\pgfqpoint{6.400000in}{0.000000in}}%
\pgfpathlineto{\pgfqpoint{6.400000in}{4.780000in}}%
\pgfpathlineto{\pgfqpoint{0.000000in}{4.780000in}}%
\pgfpathclose%
\pgfusepath{fill}%
\end{pgfscope}%
\begin{pgfscope}%
\pgfsetbuttcap%
\pgfsetmiterjoin%
\definecolor{currentfill}{rgb}{1.000000,1.000000,1.000000}%
\pgfsetfillcolor{currentfill}%
\pgfsetlinewidth{0.000000pt}%
\definecolor{currentstroke}{rgb}{0.000000,0.000000,0.000000}%
\pgfsetstrokecolor{currentstroke}%
\pgfsetstrokeopacity{0.000000}%
\pgfsetdash{}{0pt}%
\pgfpathmoveto{\pgfqpoint{0.798865in}{0.725653in}}%
\pgfpathlineto{\pgfqpoint{6.190000in}{0.725653in}}%
\pgfpathlineto{\pgfqpoint{6.190000in}{4.570875in}}%
\pgfpathlineto{\pgfqpoint{0.798865in}{4.570875in}}%
\pgfpathclose%
\pgfusepath{fill}%
\end{pgfscope}%
\begin{pgfscope}%
\pgfsetbuttcap%
\pgfsetroundjoin%
\definecolor{currentfill}{rgb}{0.000000,0.000000,0.000000}%
\pgfsetfillcolor{currentfill}%
\pgfsetlinewidth{0.803000pt}%
\definecolor{currentstroke}{rgb}{0.000000,0.000000,0.000000}%
\pgfsetstrokecolor{currentstroke}%
\pgfsetdash{}{0pt}%
\pgfsys@defobject{currentmarker}{\pgfqpoint{0.000000in}{-0.048611in}}{\pgfqpoint{0.000000in}{0.000000in}}{%
\pgfpathmoveto{\pgfqpoint{0.000000in}{0.000000in}}%
\pgfpathlineto{\pgfqpoint{0.000000in}{-0.048611in}}%
\pgfusepath{stroke,fill}%
}%
\begin{pgfscope}%
\pgfsys@transformshift{1.962835in}{0.725653in}%
\pgfsys@useobject{currentmarker}{}%
\end{pgfscope}%
\end{pgfscope}%
\begin{pgfscope}%
\definecolor{textcolor}{rgb}{0.000000,0.000000,0.000000}%
\pgfsetstrokecolor{textcolor}%
\pgfsetfillcolor{textcolor}%
\pgftext[x=1.962835in,y=0.628430in,,top]{\color{textcolor}\rmfamily\fontsize{14.000000}{16.800000}\selectfont \(\displaystyle 10^{-2}\)}%
\end{pgfscope}%
\begin{pgfscope}%
\pgfsetbuttcap%
\pgfsetroundjoin%
\definecolor{currentfill}{rgb}{0.000000,0.000000,0.000000}%
\pgfsetfillcolor{currentfill}%
\pgfsetlinewidth{0.803000pt}%
\definecolor{currentstroke}{rgb}{0.000000,0.000000,0.000000}%
\pgfsetstrokecolor{currentstroke}%
\pgfsetdash{}{0pt}%
\pgfsys@defobject{currentmarker}{\pgfqpoint{0.000000in}{-0.048611in}}{\pgfqpoint{0.000000in}{0.000000in}}{%
\pgfpathmoveto{\pgfqpoint{0.000000in}{0.000000in}}%
\pgfpathlineto{\pgfqpoint{0.000000in}{-0.048611in}}%
\pgfusepath{stroke,fill}%
}%
\begin{pgfscope}%
\pgfsys@transformshift{3.953892in}{0.725653in}%
\pgfsys@useobject{currentmarker}{}%
\end{pgfscope}%
\end{pgfscope}%
\begin{pgfscope}%
\definecolor{textcolor}{rgb}{0.000000,0.000000,0.000000}%
\pgfsetstrokecolor{textcolor}%
\pgfsetfillcolor{textcolor}%
\pgftext[x=3.953892in,y=0.628430in,,top]{\color{textcolor}\rmfamily\fontsize{14.000000}{16.800000}\selectfont \(\displaystyle 10^{-1}\)}%
\end{pgfscope}%
\begin{pgfscope}%
\pgfsetbuttcap%
\pgfsetroundjoin%
\definecolor{currentfill}{rgb}{0.000000,0.000000,0.000000}%
\pgfsetfillcolor{currentfill}%
\pgfsetlinewidth{0.803000pt}%
\definecolor{currentstroke}{rgb}{0.000000,0.000000,0.000000}%
\pgfsetstrokecolor{currentstroke}%
\pgfsetdash{}{0pt}%
\pgfsys@defobject{currentmarker}{\pgfqpoint{0.000000in}{-0.048611in}}{\pgfqpoint{0.000000in}{0.000000in}}{%
\pgfpathmoveto{\pgfqpoint{0.000000in}{0.000000in}}%
\pgfpathlineto{\pgfqpoint{0.000000in}{-0.048611in}}%
\pgfusepath{stroke,fill}%
}%
\begin{pgfscope}%
\pgfsys@transformshift{5.944948in}{0.725653in}%
\pgfsys@useobject{currentmarker}{}%
\end{pgfscope}%
\end{pgfscope}%
\begin{pgfscope}%
\definecolor{textcolor}{rgb}{0.000000,0.000000,0.000000}%
\pgfsetstrokecolor{textcolor}%
\pgfsetfillcolor{textcolor}%
\pgftext[x=5.944948in,y=0.628430in,,top]{\color{textcolor}\rmfamily\fontsize{14.000000}{16.800000}\selectfont \(\displaystyle 10^{0}\)}%
\end{pgfscope}%
\begin{pgfscope}%
\pgfsetbuttcap%
\pgfsetroundjoin%
\definecolor{currentfill}{rgb}{0.000000,0.000000,0.000000}%
\pgfsetfillcolor{currentfill}%
\pgfsetlinewidth{0.602250pt}%
\definecolor{currentstroke}{rgb}{0.000000,0.000000,0.000000}%
\pgfsetstrokecolor{currentstroke}%
\pgfsetdash{}{0pt}%
\pgfsys@defobject{currentmarker}{\pgfqpoint{0.000000in}{-0.027778in}}{\pgfqpoint{0.000000in}{0.000000in}}{%
\pgfpathmoveto{\pgfqpoint{0.000000in}{0.000000in}}%
\pgfpathlineto{\pgfqpoint{0.000000in}{-0.027778in}}%
\pgfusepath{stroke,fill}%
}%
\begin{pgfscope}%
\pgfsys@transformshift{0.921754in}{0.725653in}%
\pgfsys@useobject{currentmarker}{}%
\end{pgfscope}%
\end{pgfscope}%
\begin{pgfscope}%
\pgfsetbuttcap%
\pgfsetroundjoin%
\definecolor{currentfill}{rgb}{0.000000,0.000000,0.000000}%
\pgfsetfillcolor{currentfill}%
\pgfsetlinewidth{0.602250pt}%
\definecolor{currentstroke}{rgb}{0.000000,0.000000,0.000000}%
\pgfsetstrokecolor{currentstroke}%
\pgfsetdash{}{0pt}%
\pgfsys@defobject{currentmarker}{\pgfqpoint{0.000000in}{-0.027778in}}{\pgfqpoint{0.000000in}{0.000000in}}{%
\pgfpathmoveto{\pgfqpoint{0.000000in}{0.000000in}}%
\pgfpathlineto{\pgfqpoint{0.000000in}{-0.027778in}}%
\pgfusepath{stroke,fill}%
}%
\begin{pgfscope}%
\pgfsys@transformshift{1.170514in}{0.725653in}%
\pgfsys@useobject{currentmarker}{}%
\end{pgfscope}%
\end{pgfscope}%
\begin{pgfscope}%
\pgfsetbuttcap%
\pgfsetroundjoin%
\definecolor{currentfill}{rgb}{0.000000,0.000000,0.000000}%
\pgfsetfillcolor{currentfill}%
\pgfsetlinewidth{0.602250pt}%
\definecolor{currentstroke}{rgb}{0.000000,0.000000,0.000000}%
\pgfsetstrokecolor{currentstroke}%
\pgfsetdash{}{0pt}%
\pgfsys@defobject{currentmarker}{\pgfqpoint{0.000000in}{-0.027778in}}{\pgfqpoint{0.000000in}{0.000000in}}{%
\pgfpathmoveto{\pgfqpoint{0.000000in}{0.000000in}}%
\pgfpathlineto{\pgfqpoint{0.000000in}{-0.027778in}}%
\pgfusepath{stroke,fill}%
}%
\begin{pgfscope}%
\pgfsys@transformshift{1.363467in}{0.725653in}%
\pgfsys@useobject{currentmarker}{}%
\end{pgfscope}%
\end{pgfscope}%
\begin{pgfscope}%
\pgfsetbuttcap%
\pgfsetroundjoin%
\definecolor{currentfill}{rgb}{0.000000,0.000000,0.000000}%
\pgfsetfillcolor{currentfill}%
\pgfsetlinewidth{0.602250pt}%
\definecolor{currentstroke}{rgb}{0.000000,0.000000,0.000000}%
\pgfsetstrokecolor{currentstroke}%
\pgfsetdash{}{0pt}%
\pgfsys@defobject{currentmarker}{\pgfqpoint{0.000000in}{-0.027778in}}{\pgfqpoint{0.000000in}{0.000000in}}{%
\pgfpathmoveto{\pgfqpoint{0.000000in}{0.000000in}}%
\pgfpathlineto{\pgfqpoint{0.000000in}{-0.027778in}}%
\pgfusepath{stroke,fill}%
}%
\begin{pgfscope}%
\pgfsys@transformshift{1.521122in}{0.725653in}%
\pgfsys@useobject{currentmarker}{}%
\end{pgfscope}%
\end{pgfscope}%
\begin{pgfscope}%
\pgfsetbuttcap%
\pgfsetroundjoin%
\definecolor{currentfill}{rgb}{0.000000,0.000000,0.000000}%
\pgfsetfillcolor{currentfill}%
\pgfsetlinewidth{0.602250pt}%
\definecolor{currentstroke}{rgb}{0.000000,0.000000,0.000000}%
\pgfsetstrokecolor{currentstroke}%
\pgfsetdash{}{0pt}%
\pgfsys@defobject{currentmarker}{\pgfqpoint{0.000000in}{-0.027778in}}{\pgfqpoint{0.000000in}{0.000000in}}{%
\pgfpathmoveto{\pgfqpoint{0.000000in}{0.000000in}}%
\pgfpathlineto{\pgfqpoint{0.000000in}{-0.027778in}}%
\pgfusepath{stroke,fill}%
}%
\begin{pgfscope}%
\pgfsys@transformshift{1.654417in}{0.725653in}%
\pgfsys@useobject{currentmarker}{}%
\end{pgfscope}%
\end{pgfscope}%
\begin{pgfscope}%
\pgfsetbuttcap%
\pgfsetroundjoin%
\definecolor{currentfill}{rgb}{0.000000,0.000000,0.000000}%
\pgfsetfillcolor{currentfill}%
\pgfsetlinewidth{0.602250pt}%
\definecolor{currentstroke}{rgb}{0.000000,0.000000,0.000000}%
\pgfsetstrokecolor{currentstroke}%
\pgfsetdash{}{0pt}%
\pgfsys@defobject{currentmarker}{\pgfqpoint{0.000000in}{-0.027778in}}{\pgfqpoint{0.000000in}{0.000000in}}{%
\pgfpathmoveto{\pgfqpoint{0.000000in}{0.000000in}}%
\pgfpathlineto{\pgfqpoint{0.000000in}{-0.027778in}}%
\pgfusepath{stroke,fill}%
}%
\begin{pgfscope}%
\pgfsys@transformshift{1.769882in}{0.725653in}%
\pgfsys@useobject{currentmarker}{}%
\end{pgfscope}%
\end{pgfscope}%
\begin{pgfscope}%
\pgfsetbuttcap%
\pgfsetroundjoin%
\definecolor{currentfill}{rgb}{0.000000,0.000000,0.000000}%
\pgfsetfillcolor{currentfill}%
\pgfsetlinewidth{0.602250pt}%
\definecolor{currentstroke}{rgb}{0.000000,0.000000,0.000000}%
\pgfsetstrokecolor{currentstroke}%
\pgfsetdash{}{0pt}%
\pgfsys@defobject{currentmarker}{\pgfqpoint{0.000000in}{-0.027778in}}{\pgfqpoint{0.000000in}{0.000000in}}{%
\pgfpathmoveto{\pgfqpoint{0.000000in}{0.000000in}}%
\pgfpathlineto{\pgfqpoint{0.000000in}{-0.027778in}}%
\pgfusepath{stroke,fill}%
}%
\begin{pgfscope}%
\pgfsys@transformshift{1.871729in}{0.725653in}%
\pgfsys@useobject{currentmarker}{}%
\end{pgfscope}%
\end{pgfscope}%
\begin{pgfscope}%
\pgfsetbuttcap%
\pgfsetroundjoin%
\definecolor{currentfill}{rgb}{0.000000,0.000000,0.000000}%
\pgfsetfillcolor{currentfill}%
\pgfsetlinewidth{0.602250pt}%
\definecolor{currentstroke}{rgb}{0.000000,0.000000,0.000000}%
\pgfsetstrokecolor{currentstroke}%
\pgfsetdash{}{0pt}%
\pgfsys@defobject{currentmarker}{\pgfqpoint{0.000000in}{-0.027778in}}{\pgfqpoint{0.000000in}{0.000000in}}{%
\pgfpathmoveto{\pgfqpoint{0.000000in}{0.000000in}}%
\pgfpathlineto{\pgfqpoint{0.000000in}{-0.027778in}}%
\pgfusepath{stroke,fill}%
}%
\begin{pgfscope}%
\pgfsys@transformshift{2.562203in}{0.725653in}%
\pgfsys@useobject{currentmarker}{}%
\end{pgfscope}%
\end{pgfscope}%
\begin{pgfscope}%
\pgfsetbuttcap%
\pgfsetroundjoin%
\definecolor{currentfill}{rgb}{0.000000,0.000000,0.000000}%
\pgfsetfillcolor{currentfill}%
\pgfsetlinewidth{0.602250pt}%
\definecolor{currentstroke}{rgb}{0.000000,0.000000,0.000000}%
\pgfsetstrokecolor{currentstroke}%
\pgfsetdash{}{0pt}%
\pgfsys@defobject{currentmarker}{\pgfqpoint{0.000000in}{-0.027778in}}{\pgfqpoint{0.000000in}{0.000000in}}{%
\pgfpathmoveto{\pgfqpoint{0.000000in}{0.000000in}}%
\pgfpathlineto{\pgfqpoint{0.000000in}{-0.027778in}}%
\pgfusepath{stroke,fill}%
}%
\begin{pgfscope}%
\pgfsys@transformshift{2.912811in}{0.725653in}%
\pgfsys@useobject{currentmarker}{}%
\end{pgfscope}%
\end{pgfscope}%
\begin{pgfscope}%
\pgfsetbuttcap%
\pgfsetroundjoin%
\definecolor{currentfill}{rgb}{0.000000,0.000000,0.000000}%
\pgfsetfillcolor{currentfill}%
\pgfsetlinewidth{0.602250pt}%
\definecolor{currentstroke}{rgb}{0.000000,0.000000,0.000000}%
\pgfsetstrokecolor{currentstroke}%
\pgfsetdash{}{0pt}%
\pgfsys@defobject{currentmarker}{\pgfqpoint{0.000000in}{-0.027778in}}{\pgfqpoint{0.000000in}{0.000000in}}{%
\pgfpathmoveto{\pgfqpoint{0.000000in}{0.000000in}}%
\pgfpathlineto{\pgfqpoint{0.000000in}{-0.027778in}}%
\pgfusepath{stroke,fill}%
}%
\begin{pgfscope}%
\pgfsys@transformshift{3.161571in}{0.725653in}%
\pgfsys@useobject{currentmarker}{}%
\end{pgfscope}%
\end{pgfscope}%
\begin{pgfscope}%
\pgfsetbuttcap%
\pgfsetroundjoin%
\definecolor{currentfill}{rgb}{0.000000,0.000000,0.000000}%
\pgfsetfillcolor{currentfill}%
\pgfsetlinewidth{0.602250pt}%
\definecolor{currentstroke}{rgb}{0.000000,0.000000,0.000000}%
\pgfsetstrokecolor{currentstroke}%
\pgfsetdash{}{0pt}%
\pgfsys@defobject{currentmarker}{\pgfqpoint{0.000000in}{-0.027778in}}{\pgfqpoint{0.000000in}{0.000000in}}{%
\pgfpathmoveto{\pgfqpoint{0.000000in}{0.000000in}}%
\pgfpathlineto{\pgfqpoint{0.000000in}{-0.027778in}}%
\pgfusepath{stroke,fill}%
}%
\begin{pgfscope}%
\pgfsys@transformshift{3.354524in}{0.725653in}%
\pgfsys@useobject{currentmarker}{}%
\end{pgfscope}%
\end{pgfscope}%
\begin{pgfscope}%
\pgfsetbuttcap%
\pgfsetroundjoin%
\definecolor{currentfill}{rgb}{0.000000,0.000000,0.000000}%
\pgfsetfillcolor{currentfill}%
\pgfsetlinewidth{0.602250pt}%
\definecolor{currentstroke}{rgb}{0.000000,0.000000,0.000000}%
\pgfsetstrokecolor{currentstroke}%
\pgfsetdash{}{0pt}%
\pgfsys@defobject{currentmarker}{\pgfqpoint{0.000000in}{-0.027778in}}{\pgfqpoint{0.000000in}{0.000000in}}{%
\pgfpathmoveto{\pgfqpoint{0.000000in}{0.000000in}}%
\pgfpathlineto{\pgfqpoint{0.000000in}{-0.027778in}}%
\pgfusepath{stroke,fill}%
}%
\begin{pgfscope}%
\pgfsys@transformshift{3.512178in}{0.725653in}%
\pgfsys@useobject{currentmarker}{}%
\end{pgfscope}%
\end{pgfscope}%
\begin{pgfscope}%
\pgfsetbuttcap%
\pgfsetroundjoin%
\definecolor{currentfill}{rgb}{0.000000,0.000000,0.000000}%
\pgfsetfillcolor{currentfill}%
\pgfsetlinewidth{0.602250pt}%
\definecolor{currentstroke}{rgb}{0.000000,0.000000,0.000000}%
\pgfsetstrokecolor{currentstroke}%
\pgfsetdash{}{0pt}%
\pgfsys@defobject{currentmarker}{\pgfqpoint{0.000000in}{-0.027778in}}{\pgfqpoint{0.000000in}{0.000000in}}{%
\pgfpathmoveto{\pgfqpoint{0.000000in}{0.000000in}}%
\pgfpathlineto{\pgfqpoint{0.000000in}{-0.027778in}}%
\pgfusepath{stroke,fill}%
}%
\begin{pgfscope}%
\pgfsys@transformshift{3.645473in}{0.725653in}%
\pgfsys@useobject{currentmarker}{}%
\end{pgfscope}%
\end{pgfscope}%
\begin{pgfscope}%
\pgfsetbuttcap%
\pgfsetroundjoin%
\definecolor{currentfill}{rgb}{0.000000,0.000000,0.000000}%
\pgfsetfillcolor{currentfill}%
\pgfsetlinewidth{0.602250pt}%
\definecolor{currentstroke}{rgb}{0.000000,0.000000,0.000000}%
\pgfsetstrokecolor{currentstroke}%
\pgfsetdash{}{0pt}%
\pgfsys@defobject{currentmarker}{\pgfqpoint{0.000000in}{-0.027778in}}{\pgfqpoint{0.000000in}{0.000000in}}{%
\pgfpathmoveto{\pgfqpoint{0.000000in}{0.000000in}}%
\pgfpathlineto{\pgfqpoint{0.000000in}{-0.027778in}}%
\pgfusepath{stroke,fill}%
}%
\begin{pgfscope}%
\pgfsys@transformshift{3.760938in}{0.725653in}%
\pgfsys@useobject{currentmarker}{}%
\end{pgfscope}%
\end{pgfscope}%
\begin{pgfscope}%
\pgfsetbuttcap%
\pgfsetroundjoin%
\definecolor{currentfill}{rgb}{0.000000,0.000000,0.000000}%
\pgfsetfillcolor{currentfill}%
\pgfsetlinewidth{0.602250pt}%
\definecolor{currentstroke}{rgb}{0.000000,0.000000,0.000000}%
\pgfsetstrokecolor{currentstroke}%
\pgfsetdash{}{0pt}%
\pgfsys@defobject{currentmarker}{\pgfqpoint{0.000000in}{-0.027778in}}{\pgfqpoint{0.000000in}{0.000000in}}{%
\pgfpathmoveto{\pgfqpoint{0.000000in}{0.000000in}}%
\pgfpathlineto{\pgfqpoint{0.000000in}{-0.027778in}}%
\pgfusepath{stroke,fill}%
}%
\begin{pgfscope}%
\pgfsys@transformshift{3.862786in}{0.725653in}%
\pgfsys@useobject{currentmarker}{}%
\end{pgfscope}%
\end{pgfscope}%
\begin{pgfscope}%
\pgfsetbuttcap%
\pgfsetroundjoin%
\definecolor{currentfill}{rgb}{0.000000,0.000000,0.000000}%
\pgfsetfillcolor{currentfill}%
\pgfsetlinewidth{0.602250pt}%
\definecolor{currentstroke}{rgb}{0.000000,0.000000,0.000000}%
\pgfsetstrokecolor{currentstroke}%
\pgfsetdash{}{0pt}%
\pgfsys@defobject{currentmarker}{\pgfqpoint{0.000000in}{-0.027778in}}{\pgfqpoint{0.000000in}{0.000000in}}{%
\pgfpathmoveto{\pgfqpoint{0.000000in}{0.000000in}}%
\pgfpathlineto{\pgfqpoint{0.000000in}{-0.027778in}}%
\pgfusepath{stroke,fill}%
}%
\begin{pgfscope}%
\pgfsys@transformshift{4.553260in}{0.725653in}%
\pgfsys@useobject{currentmarker}{}%
\end{pgfscope}%
\end{pgfscope}%
\begin{pgfscope}%
\pgfsetbuttcap%
\pgfsetroundjoin%
\definecolor{currentfill}{rgb}{0.000000,0.000000,0.000000}%
\pgfsetfillcolor{currentfill}%
\pgfsetlinewidth{0.602250pt}%
\definecolor{currentstroke}{rgb}{0.000000,0.000000,0.000000}%
\pgfsetstrokecolor{currentstroke}%
\pgfsetdash{}{0pt}%
\pgfsys@defobject{currentmarker}{\pgfqpoint{0.000000in}{-0.027778in}}{\pgfqpoint{0.000000in}{0.000000in}}{%
\pgfpathmoveto{\pgfqpoint{0.000000in}{0.000000in}}%
\pgfpathlineto{\pgfqpoint{0.000000in}{-0.027778in}}%
\pgfusepath{stroke,fill}%
}%
\begin{pgfscope}%
\pgfsys@transformshift{4.903867in}{0.725653in}%
\pgfsys@useobject{currentmarker}{}%
\end{pgfscope}%
\end{pgfscope}%
\begin{pgfscope}%
\pgfsetbuttcap%
\pgfsetroundjoin%
\definecolor{currentfill}{rgb}{0.000000,0.000000,0.000000}%
\pgfsetfillcolor{currentfill}%
\pgfsetlinewidth{0.602250pt}%
\definecolor{currentstroke}{rgb}{0.000000,0.000000,0.000000}%
\pgfsetstrokecolor{currentstroke}%
\pgfsetdash{}{0pt}%
\pgfsys@defobject{currentmarker}{\pgfqpoint{0.000000in}{-0.027778in}}{\pgfqpoint{0.000000in}{0.000000in}}{%
\pgfpathmoveto{\pgfqpoint{0.000000in}{0.000000in}}%
\pgfpathlineto{\pgfqpoint{0.000000in}{-0.027778in}}%
\pgfusepath{stroke,fill}%
}%
\begin{pgfscope}%
\pgfsys@transformshift{5.152627in}{0.725653in}%
\pgfsys@useobject{currentmarker}{}%
\end{pgfscope}%
\end{pgfscope}%
\begin{pgfscope}%
\pgfsetbuttcap%
\pgfsetroundjoin%
\definecolor{currentfill}{rgb}{0.000000,0.000000,0.000000}%
\pgfsetfillcolor{currentfill}%
\pgfsetlinewidth{0.602250pt}%
\definecolor{currentstroke}{rgb}{0.000000,0.000000,0.000000}%
\pgfsetstrokecolor{currentstroke}%
\pgfsetdash{}{0pt}%
\pgfsys@defobject{currentmarker}{\pgfqpoint{0.000000in}{-0.027778in}}{\pgfqpoint{0.000000in}{0.000000in}}{%
\pgfpathmoveto{\pgfqpoint{0.000000in}{0.000000in}}%
\pgfpathlineto{\pgfqpoint{0.000000in}{-0.027778in}}%
\pgfusepath{stroke,fill}%
}%
\begin{pgfscope}%
\pgfsys@transformshift{5.345581in}{0.725653in}%
\pgfsys@useobject{currentmarker}{}%
\end{pgfscope}%
\end{pgfscope}%
\begin{pgfscope}%
\pgfsetbuttcap%
\pgfsetroundjoin%
\definecolor{currentfill}{rgb}{0.000000,0.000000,0.000000}%
\pgfsetfillcolor{currentfill}%
\pgfsetlinewidth{0.602250pt}%
\definecolor{currentstroke}{rgb}{0.000000,0.000000,0.000000}%
\pgfsetstrokecolor{currentstroke}%
\pgfsetdash{}{0pt}%
\pgfsys@defobject{currentmarker}{\pgfqpoint{0.000000in}{-0.027778in}}{\pgfqpoint{0.000000in}{0.000000in}}{%
\pgfpathmoveto{\pgfqpoint{0.000000in}{0.000000in}}%
\pgfpathlineto{\pgfqpoint{0.000000in}{-0.027778in}}%
\pgfusepath{stroke,fill}%
}%
\begin{pgfscope}%
\pgfsys@transformshift{5.503235in}{0.725653in}%
\pgfsys@useobject{currentmarker}{}%
\end{pgfscope}%
\end{pgfscope}%
\begin{pgfscope}%
\pgfsetbuttcap%
\pgfsetroundjoin%
\definecolor{currentfill}{rgb}{0.000000,0.000000,0.000000}%
\pgfsetfillcolor{currentfill}%
\pgfsetlinewidth{0.602250pt}%
\definecolor{currentstroke}{rgb}{0.000000,0.000000,0.000000}%
\pgfsetstrokecolor{currentstroke}%
\pgfsetdash{}{0pt}%
\pgfsys@defobject{currentmarker}{\pgfqpoint{0.000000in}{-0.027778in}}{\pgfqpoint{0.000000in}{0.000000in}}{%
\pgfpathmoveto{\pgfqpoint{0.000000in}{0.000000in}}%
\pgfpathlineto{\pgfqpoint{0.000000in}{-0.027778in}}%
\pgfusepath{stroke,fill}%
}%
\begin{pgfscope}%
\pgfsys@transformshift{5.636530in}{0.725653in}%
\pgfsys@useobject{currentmarker}{}%
\end{pgfscope}%
\end{pgfscope}%
\begin{pgfscope}%
\pgfsetbuttcap%
\pgfsetroundjoin%
\definecolor{currentfill}{rgb}{0.000000,0.000000,0.000000}%
\pgfsetfillcolor{currentfill}%
\pgfsetlinewidth{0.602250pt}%
\definecolor{currentstroke}{rgb}{0.000000,0.000000,0.000000}%
\pgfsetstrokecolor{currentstroke}%
\pgfsetdash{}{0pt}%
\pgfsys@defobject{currentmarker}{\pgfqpoint{0.000000in}{-0.027778in}}{\pgfqpoint{0.000000in}{0.000000in}}{%
\pgfpathmoveto{\pgfqpoint{0.000000in}{0.000000in}}%
\pgfpathlineto{\pgfqpoint{0.000000in}{-0.027778in}}%
\pgfusepath{stroke,fill}%
}%
\begin{pgfscope}%
\pgfsys@transformshift{5.751995in}{0.725653in}%
\pgfsys@useobject{currentmarker}{}%
\end{pgfscope}%
\end{pgfscope}%
\begin{pgfscope}%
\pgfsetbuttcap%
\pgfsetroundjoin%
\definecolor{currentfill}{rgb}{0.000000,0.000000,0.000000}%
\pgfsetfillcolor{currentfill}%
\pgfsetlinewidth{0.602250pt}%
\definecolor{currentstroke}{rgb}{0.000000,0.000000,0.000000}%
\pgfsetstrokecolor{currentstroke}%
\pgfsetdash{}{0pt}%
\pgfsys@defobject{currentmarker}{\pgfqpoint{0.000000in}{-0.027778in}}{\pgfqpoint{0.000000in}{0.000000in}}{%
\pgfpathmoveto{\pgfqpoint{0.000000in}{0.000000in}}%
\pgfpathlineto{\pgfqpoint{0.000000in}{-0.027778in}}%
\pgfusepath{stroke,fill}%
}%
\begin{pgfscope}%
\pgfsys@transformshift{5.853843in}{0.725653in}%
\pgfsys@useobject{currentmarker}{}%
\end{pgfscope}%
\end{pgfscope}%
\begin{pgfscope}%
\definecolor{textcolor}{rgb}{0.000000,0.000000,0.000000}%
\pgfsetstrokecolor{textcolor}%
\pgfsetfillcolor{textcolor}%
\pgftext[x=3.494432in,y=0.384697in,,top]{\color{textcolor}\rmfamily\fontsize{14.000000}{16.800000}\selectfont Distance from wall \(\displaystyle \eta = y/h\) - log scale}%
\end{pgfscope}%
\begin{pgfscope}%
\pgfsetbuttcap%
\pgfsetroundjoin%
\definecolor{currentfill}{rgb}{0.000000,0.000000,0.000000}%
\pgfsetfillcolor{currentfill}%
\pgfsetlinewidth{0.803000pt}%
\definecolor{currentstroke}{rgb}{0.000000,0.000000,0.000000}%
\pgfsetstrokecolor{currentstroke}%
\pgfsetdash{}{0pt}%
\pgfsys@defobject{currentmarker}{\pgfqpoint{-0.048611in}{0.000000in}}{\pgfqpoint{0.000000in}{0.000000in}}{%
\pgfpathmoveto{\pgfqpoint{0.000000in}{0.000000in}}%
\pgfpathlineto{\pgfqpoint{-0.048611in}{0.000000in}}%
\pgfusepath{stroke,fill}%
}%
\begin{pgfscope}%
\pgfsys@transformshift{0.798865in}{0.815849in}%
\pgfsys@useobject{currentmarker}{}%
\end{pgfscope}%
\end{pgfscope}%
\begin{pgfscope}%
\definecolor{textcolor}{rgb}{0.000000,0.000000,0.000000}%
\pgfsetstrokecolor{textcolor}%
\pgfsetfillcolor{textcolor}%
\pgftext[x=0.451415in,y=0.741983in,left,base]{\color{textcolor}\rmfamily\fontsize{14.000000}{16.800000}\selectfont \(\displaystyle 2.0\)}%
\end{pgfscope}%
\begin{pgfscope}%
\pgfsetbuttcap%
\pgfsetroundjoin%
\definecolor{currentfill}{rgb}{0.000000,0.000000,0.000000}%
\pgfsetfillcolor{currentfill}%
\pgfsetlinewidth{0.803000pt}%
\definecolor{currentstroke}{rgb}{0.000000,0.000000,0.000000}%
\pgfsetstrokecolor{currentstroke}%
\pgfsetdash{}{0pt}%
\pgfsys@defobject{currentmarker}{\pgfqpoint{-0.048611in}{0.000000in}}{\pgfqpoint{0.000000in}{0.000000in}}{%
\pgfpathmoveto{\pgfqpoint{0.000000in}{0.000000in}}%
\pgfpathlineto{\pgfqpoint{-0.048611in}{0.000000in}}%
\pgfusepath{stroke,fill}%
}%
\begin{pgfscope}%
\pgfsys@transformshift{0.798865in}{1.369946in}%
\pgfsys@useobject{currentmarker}{}%
\end{pgfscope}%
\end{pgfscope}%
\begin{pgfscope}%
\definecolor{textcolor}{rgb}{0.000000,0.000000,0.000000}%
\pgfsetstrokecolor{textcolor}%
\pgfsetfillcolor{textcolor}%
\pgftext[x=0.451415in,y=1.296080in,left,base]{\color{textcolor}\rmfamily\fontsize{14.000000}{16.800000}\selectfont \(\displaystyle 2.5\)}%
\end{pgfscope}%
\begin{pgfscope}%
\pgfsetbuttcap%
\pgfsetroundjoin%
\definecolor{currentfill}{rgb}{0.000000,0.000000,0.000000}%
\pgfsetfillcolor{currentfill}%
\pgfsetlinewidth{0.803000pt}%
\definecolor{currentstroke}{rgb}{0.000000,0.000000,0.000000}%
\pgfsetstrokecolor{currentstroke}%
\pgfsetdash{}{0pt}%
\pgfsys@defobject{currentmarker}{\pgfqpoint{-0.048611in}{0.000000in}}{\pgfqpoint{0.000000in}{0.000000in}}{%
\pgfpathmoveto{\pgfqpoint{0.000000in}{0.000000in}}%
\pgfpathlineto{\pgfqpoint{-0.048611in}{0.000000in}}%
\pgfusepath{stroke,fill}%
}%
\begin{pgfscope}%
\pgfsys@transformshift{0.798865in}{1.924043in}%
\pgfsys@useobject{currentmarker}{}%
\end{pgfscope}%
\end{pgfscope}%
\begin{pgfscope}%
\definecolor{textcolor}{rgb}{0.000000,0.000000,0.000000}%
\pgfsetstrokecolor{textcolor}%
\pgfsetfillcolor{textcolor}%
\pgftext[x=0.451415in,y=1.850177in,left,base]{\color{textcolor}\rmfamily\fontsize{14.000000}{16.800000}\selectfont \(\displaystyle 3.0\)}%
\end{pgfscope}%
\begin{pgfscope}%
\pgfsetbuttcap%
\pgfsetroundjoin%
\definecolor{currentfill}{rgb}{0.000000,0.000000,0.000000}%
\pgfsetfillcolor{currentfill}%
\pgfsetlinewidth{0.803000pt}%
\definecolor{currentstroke}{rgb}{0.000000,0.000000,0.000000}%
\pgfsetstrokecolor{currentstroke}%
\pgfsetdash{}{0pt}%
\pgfsys@defobject{currentmarker}{\pgfqpoint{-0.048611in}{0.000000in}}{\pgfqpoint{0.000000in}{0.000000in}}{%
\pgfpathmoveto{\pgfqpoint{0.000000in}{0.000000in}}%
\pgfpathlineto{\pgfqpoint{-0.048611in}{0.000000in}}%
\pgfusepath{stroke,fill}%
}%
\begin{pgfscope}%
\pgfsys@transformshift{0.798865in}{2.478140in}%
\pgfsys@useobject{currentmarker}{}%
\end{pgfscope}%
\end{pgfscope}%
\begin{pgfscope}%
\definecolor{textcolor}{rgb}{0.000000,0.000000,0.000000}%
\pgfsetstrokecolor{textcolor}%
\pgfsetfillcolor{textcolor}%
\pgftext[x=0.451415in,y=2.404274in,left,base]{\color{textcolor}\rmfamily\fontsize{14.000000}{16.800000}\selectfont \(\displaystyle 3.5\)}%
\end{pgfscope}%
\begin{pgfscope}%
\pgfsetbuttcap%
\pgfsetroundjoin%
\definecolor{currentfill}{rgb}{0.000000,0.000000,0.000000}%
\pgfsetfillcolor{currentfill}%
\pgfsetlinewidth{0.803000pt}%
\definecolor{currentstroke}{rgb}{0.000000,0.000000,0.000000}%
\pgfsetstrokecolor{currentstroke}%
\pgfsetdash{}{0pt}%
\pgfsys@defobject{currentmarker}{\pgfqpoint{-0.048611in}{0.000000in}}{\pgfqpoint{0.000000in}{0.000000in}}{%
\pgfpathmoveto{\pgfqpoint{0.000000in}{0.000000in}}%
\pgfpathlineto{\pgfqpoint{-0.048611in}{0.000000in}}%
\pgfusepath{stroke,fill}%
}%
\begin{pgfscope}%
\pgfsys@transformshift{0.798865in}{3.032237in}%
\pgfsys@useobject{currentmarker}{}%
\end{pgfscope}%
\end{pgfscope}%
\begin{pgfscope}%
\definecolor{textcolor}{rgb}{0.000000,0.000000,0.000000}%
\pgfsetstrokecolor{textcolor}%
\pgfsetfillcolor{textcolor}%
\pgftext[x=0.451415in,y=2.958371in,left,base]{\color{textcolor}\rmfamily\fontsize{14.000000}{16.800000}\selectfont \(\displaystyle 4.0\)}%
\end{pgfscope}%
\begin{pgfscope}%
\pgfsetbuttcap%
\pgfsetroundjoin%
\definecolor{currentfill}{rgb}{0.000000,0.000000,0.000000}%
\pgfsetfillcolor{currentfill}%
\pgfsetlinewidth{0.803000pt}%
\definecolor{currentstroke}{rgb}{0.000000,0.000000,0.000000}%
\pgfsetstrokecolor{currentstroke}%
\pgfsetdash{}{0pt}%
\pgfsys@defobject{currentmarker}{\pgfqpoint{-0.048611in}{0.000000in}}{\pgfqpoint{0.000000in}{0.000000in}}{%
\pgfpathmoveto{\pgfqpoint{0.000000in}{0.000000in}}%
\pgfpathlineto{\pgfqpoint{-0.048611in}{0.000000in}}%
\pgfusepath{stroke,fill}%
}%
\begin{pgfscope}%
\pgfsys@transformshift{0.798865in}{3.586335in}%
\pgfsys@useobject{currentmarker}{}%
\end{pgfscope}%
\end{pgfscope}%
\begin{pgfscope}%
\definecolor{textcolor}{rgb}{0.000000,0.000000,0.000000}%
\pgfsetstrokecolor{textcolor}%
\pgfsetfillcolor{textcolor}%
\pgftext[x=0.451415in,y=3.512469in,left,base]{\color{textcolor}\rmfamily\fontsize{14.000000}{16.800000}\selectfont \(\displaystyle 4.5\)}%
\end{pgfscope}%
\begin{pgfscope}%
\pgfsetbuttcap%
\pgfsetroundjoin%
\definecolor{currentfill}{rgb}{0.000000,0.000000,0.000000}%
\pgfsetfillcolor{currentfill}%
\pgfsetlinewidth{0.803000pt}%
\definecolor{currentstroke}{rgb}{0.000000,0.000000,0.000000}%
\pgfsetstrokecolor{currentstroke}%
\pgfsetdash{}{0pt}%
\pgfsys@defobject{currentmarker}{\pgfqpoint{-0.048611in}{0.000000in}}{\pgfqpoint{0.000000in}{0.000000in}}{%
\pgfpathmoveto{\pgfqpoint{0.000000in}{0.000000in}}%
\pgfpathlineto{\pgfqpoint{-0.048611in}{0.000000in}}%
\pgfusepath{stroke,fill}%
}%
\begin{pgfscope}%
\pgfsys@transformshift{0.798865in}{4.140432in}%
\pgfsys@useobject{currentmarker}{}%
\end{pgfscope}%
\end{pgfscope}%
\begin{pgfscope}%
\definecolor{textcolor}{rgb}{0.000000,0.000000,0.000000}%
\pgfsetstrokecolor{textcolor}%
\pgfsetfillcolor{textcolor}%
\pgftext[x=0.451415in,y=4.066566in,left,base]{\color{textcolor}\rmfamily\fontsize{14.000000}{16.800000}\selectfont \(\displaystyle 5.0\)}%
\end{pgfscope}%
\begin{pgfscope}%
\definecolor{textcolor}{rgb}{0.000000,0.000000,0.000000}%
\pgfsetstrokecolor{textcolor}%
\pgfsetfillcolor{textcolor}%
\pgftext[x=0.395859in,y=2.648264in,,bottom,rotate=90.000000]{\color{textcolor}\rmfamily\fontsize{14.000000}{16.800000}\selectfont Mean velocity (m/s)}%
\end{pgfscope}%
\begin{pgfscope}%
\pgfpathrectangle{\pgfqpoint{0.798865in}{0.725653in}}{\pgfqpoint{5.391135in}{3.845222in}}%
\pgfusepath{clip}%
\pgfsetbuttcap%
\pgfsetroundjoin%
\pgfsetlinewidth{1.505625pt}%
\definecolor{currentstroke}{rgb}{0.000000,0.000000,0.000000}%
\pgfsetstrokecolor{currentstroke}%
\pgfsetdash{{5.550000pt}{2.400000pt}}{0.000000pt}%
\pgfpathmoveto{\pgfqpoint{1.043917in}{0.989614in}}%
\pgfpathlineto{\pgfqpoint{1.043917in}{4.396092in}}%
\pgfusepath{stroke}%
\end{pgfscope}%
\begin{pgfscope}%
\pgfpathrectangle{\pgfqpoint{0.798865in}{0.725653in}}{\pgfqpoint{5.391135in}{3.845222in}}%
\pgfusepath{clip}%
\pgfsetbuttcap%
\pgfsetroundjoin%
\pgfsetlinewidth{1.505625pt}%
\definecolor{currentstroke}{rgb}{0.000000,0.000000,0.000000}%
\pgfsetstrokecolor{currentstroke}%
\pgfsetdash{{5.550000pt}{2.400000pt}}{0.000000pt}%
\pgfpathmoveto{\pgfqpoint{4.903867in}{0.989614in}}%
\pgfpathlineto{\pgfqpoint{4.903867in}{4.396092in}}%
\pgfusepath{stroke}%
\end{pgfscope}%
\begin{pgfscope}%
\pgfpathrectangle{\pgfqpoint{0.798865in}{0.725653in}}{\pgfqpoint{5.391135in}{3.845222in}}%
\pgfusepath{clip}%
\pgfsetrectcap%
\pgfsetroundjoin%
\pgfsetlinewidth{1.505625pt}%
\definecolor{currentstroke}{rgb}{0.121569,0.466667,0.705882}%
\pgfsetstrokecolor{currentstroke}%
\pgfsetdash{}{0pt}%
\pgfpathmoveto{\pgfqpoint{1.239727in}{0.900436in}}%
\pgfpathlineto{\pgfqpoint{5.944948in}{4.288329in}}%
\pgfpathlineto{\pgfqpoint{5.944948in}{4.288329in}}%
\pgfusepath{stroke}%
\end{pgfscope}%
\begin{pgfscope}%
\pgfpathrectangle{\pgfqpoint{0.798865in}{0.725653in}}{\pgfqpoint{5.391135in}{3.845222in}}%
\pgfusepath{clip}%
\pgfsetbuttcap%
\pgfsetroundjoin%
\pgfsetlinewidth{1.505625pt}%
\definecolor{currentstroke}{rgb}{1.000000,0.498039,0.054902}%
\pgfsetstrokecolor{currentstroke}%
\pgfsetdash{{1.500000pt}{2.475000pt}}{0.000000pt}%
\pgfpathmoveto{\pgfqpoint{1.239727in}{0.989614in}}%
\pgfpathlineto{\pgfqpoint{1.363467in}{1.062201in}}%
\pgfpathlineto{\pgfqpoint{1.612228in}{1.165817in}}%
\pgfpathlineto{\pgfqpoint{1.805181in}{1.313096in}}%
\pgfpathlineto{\pgfqpoint{3.003916in}{2.162970in}}%
\pgfpathlineto{\pgfqpoint{3.603284in}{2.615224in}}%
\pgfpathlineto{\pgfqpoint{4.146845in}{3.004976in}}%
\pgfpathlineto{\pgfqpoint{4.478020in}{3.211654in}}%
\pgfpathlineto{\pgfqpoint{4.716898in}{3.371234in}}%
\pgfpathlineto{\pgfqpoint{4.903867in}{3.532476in}}%
\pgfpathlineto{\pgfqpoint{4.994973in}{3.644626in}}%
\pgfpathlineto{\pgfqpoint{5.345581in}{3.925220in}}%
\pgfpathlineto{\pgfqpoint{5.594341in}{4.157720in}}%
\pgfpathlineto{\pgfqpoint{5.787294in}{4.335141in}}%
\pgfpathlineto{\pgfqpoint{5.944948in}{4.396092in}}%
\pgfusepath{stroke}%
\end{pgfscope}%
\begin{pgfscope}%
\pgfpathrectangle{\pgfqpoint{0.798865in}{0.725653in}}{\pgfqpoint{5.391135in}{3.845222in}}%
\pgfusepath{clip}%
\pgfsetbuttcap%
\pgfsetroundjoin%
\definecolor{currentfill}{rgb}{1.000000,0.498039,0.054902}%
\pgfsetfillcolor{currentfill}%
\pgfsetlinewidth{1.003750pt}%
\definecolor{currentstroke}{rgb}{1.000000,0.498039,0.054902}%
\pgfsetstrokecolor{currentstroke}%
\pgfsetdash{}{0pt}%
\pgfsys@defobject{currentmarker}{\pgfqpoint{-0.034722in}{-0.034722in}}{\pgfqpoint{0.034722in}{0.034722in}}{%
\pgfpathmoveto{\pgfqpoint{0.000000in}{-0.034722in}}%
\pgfpathcurveto{\pgfqpoint{0.009208in}{-0.034722in}}{\pgfqpoint{0.018041in}{-0.031064in}}{\pgfqpoint{0.024552in}{-0.024552in}}%
\pgfpathcurveto{\pgfqpoint{0.031064in}{-0.018041in}}{\pgfqpoint{0.034722in}{-0.009208in}}{\pgfqpoint{0.034722in}{0.000000in}}%
\pgfpathcurveto{\pgfqpoint{0.034722in}{0.009208in}}{\pgfqpoint{0.031064in}{0.018041in}}{\pgfqpoint{0.024552in}{0.024552in}}%
\pgfpathcurveto{\pgfqpoint{0.018041in}{0.031064in}}{\pgfqpoint{0.009208in}{0.034722in}}{\pgfqpoint{0.000000in}{0.034722in}}%
\pgfpathcurveto{\pgfqpoint{-0.009208in}{0.034722in}}{\pgfqpoint{-0.018041in}{0.031064in}}{\pgfqpoint{-0.024552in}{0.024552in}}%
\pgfpathcurveto{\pgfqpoint{-0.031064in}{0.018041in}}{\pgfqpoint{-0.034722in}{0.009208in}}{\pgfqpoint{-0.034722in}{0.000000in}}%
\pgfpathcurveto{\pgfqpoint{-0.034722in}{-0.009208in}}{\pgfqpoint{-0.031064in}{-0.018041in}}{\pgfqpoint{-0.024552in}{-0.024552in}}%
\pgfpathcurveto{\pgfqpoint{-0.018041in}{-0.031064in}}{\pgfqpoint{-0.009208in}{-0.034722in}}{\pgfqpoint{0.000000in}{-0.034722in}}%
\pgfpathclose%
\pgfusepath{stroke,fill}%
}%
\begin{pgfscope}%
\pgfsys@transformshift{1.239727in}{0.989614in}%
\pgfsys@useobject{currentmarker}{}%
\end{pgfscope}%
\begin{pgfscope}%
\pgfsys@transformshift{1.363467in}{1.062201in}%
\pgfsys@useobject{currentmarker}{}%
\end{pgfscope}%
\begin{pgfscope}%
\pgfsys@transformshift{1.612228in}{1.165817in}%
\pgfsys@useobject{currentmarker}{}%
\end{pgfscope}%
\begin{pgfscope}%
\pgfsys@transformshift{1.805181in}{1.313096in}%
\pgfsys@useobject{currentmarker}{}%
\end{pgfscope}%
\begin{pgfscope}%
\pgfsys@transformshift{3.003916in}{2.162970in}%
\pgfsys@useobject{currentmarker}{}%
\end{pgfscope}%
\begin{pgfscope}%
\pgfsys@transformshift{3.603284in}{2.615224in}%
\pgfsys@useobject{currentmarker}{}%
\end{pgfscope}%
\begin{pgfscope}%
\pgfsys@transformshift{4.146845in}{3.004976in}%
\pgfsys@useobject{currentmarker}{}%
\end{pgfscope}%
\begin{pgfscope}%
\pgfsys@transformshift{4.478020in}{3.211654in}%
\pgfsys@useobject{currentmarker}{}%
\end{pgfscope}%
\begin{pgfscope}%
\pgfsys@transformshift{4.716898in}{3.371234in}%
\pgfsys@useobject{currentmarker}{}%
\end{pgfscope}%
\begin{pgfscope}%
\pgfsys@transformshift{4.903867in}{3.532476in}%
\pgfsys@useobject{currentmarker}{}%
\end{pgfscope}%
\begin{pgfscope}%
\pgfsys@transformshift{4.994973in}{3.644626in}%
\pgfsys@useobject{currentmarker}{}%
\end{pgfscope}%
\begin{pgfscope}%
\pgfsys@transformshift{5.345581in}{3.925220in}%
\pgfsys@useobject{currentmarker}{}%
\end{pgfscope}%
\begin{pgfscope}%
\pgfsys@transformshift{5.594341in}{4.157720in}%
\pgfsys@useobject{currentmarker}{}%
\end{pgfscope}%
\begin{pgfscope}%
\pgfsys@transformshift{5.787294in}{4.335141in}%
\pgfsys@useobject{currentmarker}{}%
\end{pgfscope}%
\begin{pgfscope}%
\pgfsys@transformshift{5.944948in}{4.396092in}%
\pgfsys@useobject{currentmarker}{}%
\end{pgfscope}%
\end{pgfscope}%
\begin{pgfscope}%
\pgfsetrectcap%
\pgfsetmiterjoin%
\pgfsetlinewidth{0.803000pt}%
\definecolor{currentstroke}{rgb}{0.000000,0.000000,0.000000}%
\pgfsetstrokecolor{currentstroke}%
\pgfsetdash{}{0pt}%
\pgfpathmoveto{\pgfqpoint{0.798865in}{0.725653in}}%
\pgfpathlineto{\pgfqpoint{0.798865in}{4.570875in}}%
\pgfusepath{stroke}%
\end{pgfscope}%
\begin{pgfscope}%
\pgfsetrectcap%
\pgfsetmiterjoin%
\pgfsetlinewidth{0.803000pt}%
\definecolor{currentstroke}{rgb}{0.000000,0.000000,0.000000}%
\pgfsetstrokecolor{currentstroke}%
\pgfsetdash{}{0pt}%
\pgfpathmoveto{\pgfqpoint{6.190000in}{0.725653in}}%
\pgfpathlineto{\pgfqpoint{6.190000in}{4.570875in}}%
\pgfusepath{stroke}%
\end{pgfscope}%
\begin{pgfscope}%
\pgfsetrectcap%
\pgfsetmiterjoin%
\pgfsetlinewidth{0.803000pt}%
\definecolor{currentstroke}{rgb}{0.000000,0.000000,0.000000}%
\pgfsetstrokecolor{currentstroke}%
\pgfsetdash{}{0pt}%
\pgfpathmoveto{\pgfqpoint{0.798865in}{0.725653in}}%
\pgfpathlineto{\pgfqpoint{6.190000in}{0.725653in}}%
\pgfusepath{stroke}%
\end{pgfscope}%
\begin{pgfscope}%
\pgfsetrectcap%
\pgfsetmiterjoin%
\pgfsetlinewidth{0.803000pt}%
\definecolor{currentstroke}{rgb}{0.000000,0.000000,0.000000}%
\pgfsetstrokecolor{currentstroke}%
\pgfsetdash{}{0pt}%
\pgfpathmoveto{\pgfqpoint{0.798865in}{4.570875in}}%
\pgfpathlineto{\pgfqpoint{6.190000in}{4.570875in}}%
\pgfusepath{stroke}%
\end{pgfscope}%
\begin{pgfscope}%
\definecolor{textcolor}{rgb}{0.000000,0.000000,0.000000}%
\pgfsetstrokecolor{textcolor}%
\pgfsetfillcolor{textcolor}%
\pgftext[x=1.201571in,y=3.032237in,left,base]{\color{textcolor}\rmfamily\fontsize{14.000000}{16.800000}\selectfont \(\displaystyle y^{+} = 30\)}%
\end{pgfscope}%
\begin{pgfscope}%
\definecolor{textcolor}{rgb}{0.000000,0.000000,0.000000}%
\pgfsetstrokecolor{textcolor}%
\pgfsetfillcolor{textcolor}%
\pgftext[x=5.012096in,y=2.478140in,left,base]{\color{textcolor}\rmfamily\fontsize{14.000000}{16.800000}\selectfont \(\displaystyle \eta = \frac{y}{h} = 0.3\)}%
\end{pgfscope}%
\end{pgfpicture}%
\makeatother%
\endgroup%
}
    \caption{Mean velocity profile displayed in a log-lin scale (dashed lines and orange points). Fitting with a log-law $\bar{U} = A\ln{y}+B$ with $A = \num{1.31}$ and $B=\num{3.25}$ with a normalized RMS error of \num{1.22e-2}}
    \label{fig:Mean_Velocity_Profile}
\end{figure}

\begin{figure}
    \centering
        \resizebox{0.8\linewidth}{!}{%% Creator: Matplotlib, PGF backend
%%
%% To include the figure in your LaTeX document, write
%%   \input{<filename>.pgf}
%%
%% Make sure the required packages are loaded in your preamble
%%   \usepackage{pgf}
%%
%% Figures using additional raster images can only be included by \input if
%% they are in the same directory as the main LaTeX file. For loading figures
%% from other directories you can use the `import` package
%%   \usepackage{import}
%% and then include the figures with
%%   \import{<path to file>}{<filename>.pgf}
%%
%% Matplotlib used the following preamble
%%   \usepackage{fontspec}
%%   \setmainfont{DejaVuSerif.ttf}[Path=C:/ProgramData/Anaconda3/lib/site-packages/matplotlib/mpl-data/fonts/ttf/]
%%   \setsansfont{DejaVuSans.ttf}[Path=C:/ProgramData/Anaconda3/lib/site-packages/matplotlib/mpl-data/fonts/ttf/]
%%   \setmonofont{DejaVuSansMono.ttf}[Path=C:/ProgramData/Anaconda3/lib/site-packages/matplotlib/mpl-data/fonts/ttf/]
%%
\begingroup%
\makeatletter%
\begin{pgfpicture}%
\pgfpathrectangle{\pgfpointorigin}{\pgfqpoint{6.400000in}{4.800000in}}%
\pgfusepath{use as bounding box, clip}%
\begin{pgfscope}%
\pgfsetbuttcap%
\pgfsetmiterjoin%
\definecolor{currentfill}{rgb}{1.000000,1.000000,1.000000}%
\pgfsetfillcolor{currentfill}%
\pgfsetlinewidth{0.000000pt}%
\definecolor{currentstroke}{rgb}{1.000000,1.000000,1.000000}%
\pgfsetstrokecolor{currentstroke}%
\pgfsetdash{}{0pt}%
\pgfpathmoveto{\pgfqpoint{0.000000in}{0.000000in}}%
\pgfpathlineto{\pgfqpoint{6.400000in}{0.000000in}}%
\pgfpathlineto{\pgfqpoint{6.400000in}{4.800000in}}%
\pgfpathlineto{\pgfqpoint{0.000000in}{4.800000in}}%
\pgfpathclose%
\pgfusepath{fill}%
\end{pgfscope}%
\begin{pgfscope}%
\pgfsetbuttcap%
\pgfsetmiterjoin%
\definecolor{currentfill}{rgb}{1.000000,1.000000,1.000000}%
\pgfsetfillcolor{currentfill}%
\pgfsetlinewidth{0.000000pt}%
\definecolor{currentstroke}{rgb}{0.000000,0.000000,0.000000}%
\pgfsetstrokecolor{currentstroke}%
\pgfsetstrokeopacity{0.000000}%
\pgfsetdash{}{0pt}%
\pgfpathmoveto{\pgfqpoint{0.662778in}{0.668778in}}%
\pgfpathlineto{\pgfqpoint{6.250000in}{0.668778in}}%
\pgfpathlineto{\pgfqpoint{6.250000in}{4.650000in}}%
\pgfpathlineto{\pgfqpoint{0.662778in}{4.650000in}}%
\pgfpathclose%
\pgfusepath{fill}%
\end{pgfscope}%
\begin{pgfscope}%
\pgfsetbuttcap%
\pgfsetroundjoin%
\definecolor{currentfill}{rgb}{0.000000,0.000000,0.000000}%
\pgfsetfillcolor{currentfill}%
\pgfsetlinewidth{0.803000pt}%
\definecolor{currentstroke}{rgb}{0.000000,0.000000,0.000000}%
\pgfsetstrokecolor{currentstroke}%
\pgfsetdash{}{0pt}%
\pgfsys@defobject{currentmarker}{\pgfqpoint{0.000000in}{-0.048611in}}{\pgfqpoint{0.000000in}{0.000000in}}{%
\pgfpathmoveto{\pgfqpoint{0.000000in}{0.000000in}}%
\pgfpathlineto{\pgfqpoint{0.000000in}{-0.048611in}}%
\pgfusepath{stroke,fill}%
}%
\begin{pgfscope}%
\pgfsys@transformshift{0.671838in}{0.668778in}%
\pgfsys@useobject{currentmarker}{}%
\end{pgfscope}%
\end{pgfscope}%
\begin{pgfscope}%
\definecolor{textcolor}{rgb}{0.000000,0.000000,0.000000}%
\pgfsetstrokecolor{textcolor}%
\pgfsetfillcolor{textcolor}%
\pgftext[x=0.671838in,y=0.571556in,,top]{\color{textcolor}\sffamily\fontsize{10.000000}{12.000000}\selectfont \(\displaystyle {10^{-1}}\)}%
\end{pgfscope}%
\begin{pgfscope}%
\pgfsetbuttcap%
\pgfsetroundjoin%
\definecolor{currentfill}{rgb}{0.000000,0.000000,0.000000}%
\pgfsetfillcolor{currentfill}%
\pgfsetlinewidth{0.803000pt}%
\definecolor{currentstroke}{rgb}{0.000000,0.000000,0.000000}%
\pgfsetstrokecolor{currentstroke}%
\pgfsetdash{}{0pt}%
\pgfsys@defobject{currentmarker}{\pgfqpoint{0.000000in}{-0.048611in}}{\pgfqpoint{0.000000in}{0.000000in}}{%
\pgfpathmoveto{\pgfqpoint{0.000000in}{0.000000in}}%
\pgfpathlineto{\pgfqpoint{0.000000in}{-0.048611in}}%
\pgfusepath{stroke,fill}%
}%
\begin{pgfscope}%
\pgfsys@transformshift{2.821187in}{0.668778in}%
\pgfsys@useobject{currentmarker}{}%
\end{pgfscope}%
\end{pgfscope}%
\begin{pgfscope}%
\definecolor{textcolor}{rgb}{0.000000,0.000000,0.000000}%
\pgfsetstrokecolor{textcolor}%
\pgfsetfillcolor{textcolor}%
\pgftext[x=2.821187in,y=0.571556in,,top]{\color{textcolor}\sffamily\fontsize{10.000000}{12.000000}\selectfont \(\displaystyle {10^{0}}\)}%
\end{pgfscope}%
\begin{pgfscope}%
\pgfsetbuttcap%
\pgfsetroundjoin%
\definecolor{currentfill}{rgb}{0.000000,0.000000,0.000000}%
\pgfsetfillcolor{currentfill}%
\pgfsetlinewidth{0.803000pt}%
\definecolor{currentstroke}{rgb}{0.000000,0.000000,0.000000}%
\pgfsetstrokecolor{currentstroke}%
\pgfsetdash{}{0pt}%
\pgfsys@defobject{currentmarker}{\pgfqpoint{0.000000in}{-0.048611in}}{\pgfqpoint{0.000000in}{0.000000in}}{%
\pgfpathmoveto{\pgfqpoint{0.000000in}{0.000000in}}%
\pgfpathlineto{\pgfqpoint{0.000000in}{-0.048611in}}%
\pgfusepath{stroke,fill}%
}%
\begin{pgfscope}%
\pgfsys@transformshift{4.970535in}{0.668778in}%
\pgfsys@useobject{currentmarker}{}%
\end{pgfscope}%
\end{pgfscope}%
\begin{pgfscope}%
\definecolor{textcolor}{rgb}{0.000000,0.000000,0.000000}%
\pgfsetstrokecolor{textcolor}%
\pgfsetfillcolor{textcolor}%
\pgftext[x=4.970535in,y=0.571556in,,top]{\color{textcolor}\sffamily\fontsize{10.000000}{12.000000}\selectfont \(\displaystyle {10^{1}}\)}%
\end{pgfscope}%
\begin{pgfscope}%
\pgfsetbuttcap%
\pgfsetroundjoin%
\definecolor{currentfill}{rgb}{0.000000,0.000000,0.000000}%
\pgfsetfillcolor{currentfill}%
\pgfsetlinewidth{0.602250pt}%
\definecolor{currentstroke}{rgb}{0.000000,0.000000,0.000000}%
\pgfsetstrokecolor{currentstroke}%
\pgfsetdash{}{0pt}%
\pgfsys@defobject{currentmarker}{\pgfqpoint{0.000000in}{-0.027778in}}{\pgfqpoint{0.000000in}{0.000000in}}{%
\pgfpathmoveto{\pgfqpoint{0.000000in}{0.000000in}}%
\pgfpathlineto{\pgfqpoint{0.000000in}{-0.027778in}}%
\pgfusepath{stroke,fill}%
}%
\begin{pgfscope}%
\pgfsys@transformshift{1.318857in}{0.668778in}%
\pgfsys@useobject{currentmarker}{}%
\end{pgfscope}%
\end{pgfscope}%
\begin{pgfscope}%
\pgfsetbuttcap%
\pgfsetroundjoin%
\definecolor{currentfill}{rgb}{0.000000,0.000000,0.000000}%
\pgfsetfillcolor{currentfill}%
\pgfsetlinewidth{0.602250pt}%
\definecolor{currentstroke}{rgb}{0.000000,0.000000,0.000000}%
\pgfsetstrokecolor{currentstroke}%
\pgfsetdash{}{0pt}%
\pgfsys@defobject{currentmarker}{\pgfqpoint{0.000000in}{-0.027778in}}{\pgfqpoint{0.000000in}{0.000000in}}{%
\pgfpathmoveto{\pgfqpoint{0.000000in}{0.000000in}}%
\pgfpathlineto{\pgfqpoint{0.000000in}{-0.027778in}}%
\pgfusepath{stroke,fill}%
}%
\begin{pgfscope}%
\pgfsys@transformshift{1.697338in}{0.668778in}%
\pgfsys@useobject{currentmarker}{}%
\end{pgfscope}%
\end{pgfscope}%
\begin{pgfscope}%
\pgfsetbuttcap%
\pgfsetroundjoin%
\definecolor{currentfill}{rgb}{0.000000,0.000000,0.000000}%
\pgfsetfillcolor{currentfill}%
\pgfsetlinewidth{0.602250pt}%
\definecolor{currentstroke}{rgb}{0.000000,0.000000,0.000000}%
\pgfsetstrokecolor{currentstroke}%
\pgfsetdash{}{0pt}%
\pgfsys@defobject{currentmarker}{\pgfqpoint{0.000000in}{-0.027778in}}{\pgfqpoint{0.000000in}{0.000000in}}{%
\pgfpathmoveto{\pgfqpoint{0.000000in}{0.000000in}}%
\pgfpathlineto{\pgfqpoint{0.000000in}{-0.027778in}}%
\pgfusepath{stroke,fill}%
}%
\begin{pgfscope}%
\pgfsys@transformshift{1.965875in}{0.668778in}%
\pgfsys@useobject{currentmarker}{}%
\end{pgfscope}%
\end{pgfscope}%
\begin{pgfscope}%
\pgfsetbuttcap%
\pgfsetroundjoin%
\definecolor{currentfill}{rgb}{0.000000,0.000000,0.000000}%
\pgfsetfillcolor{currentfill}%
\pgfsetlinewidth{0.602250pt}%
\definecolor{currentstroke}{rgb}{0.000000,0.000000,0.000000}%
\pgfsetstrokecolor{currentstroke}%
\pgfsetdash{}{0pt}%
\pgfsys@defobject{currentmarker}{\pgfqpoint{0.000000in}{-0.027778in}}{\pgfqpoint{0.000000in}{0.000000in}}{%
\pgfpathmoveto{\pgfqpoint{0.000000in}{0.000000in}}%
\pgfpathlineto{\pgfqpoint{0.000000in}{-0.027778in}}%
\pgfusepath{stroke,fill}%
}%
\begin{pgfscope}%
\pgfsys@transformshift{2.174169in}{0.668778in}%
\pgfsys@useobject{currentmarker}{}%
\end{pgfscope}%
\end{pgfscope}%
\begin{pgfscope}%
\pgfsetbuttcap%
\pgfsetroundjoin%
\definecolor{currentfill}{rgb}{0.000000,0.000000,0.000000}%
\pgfsetfillcolor{currentfill}%
\pgfsetlinewidth{0.602250pt}%
\definecolor{currentstroke}{rgb}{0.000000,0.000000,0.000000}%
\pgfsetstrokecolor{currentstroke}%
\pgfsetdash{}{0pt}%
\pgfsys@defobject{currentmarker}{\pgfqpoint{0.000000in}{-0.027778in}}{\pgfqpoint{0.000000in}{0.000000in}}{%
\pgfpathmoveto{\pgfqpoint{0.000000in}{0.000000in}}%
\pgfpathlineto{\pgfqpoint{0.000000in}{-0.027778in}}%
\pgfusepath{stroke,fill}%
}%
\begin{pgfscope}%
\pgfsys@transformshift{2.344357in}{0.668778in}%
\pgfsys@useobject{currentmarker}{}%
\end{pgfscope}%
\end{pgfscope}%
\begin{pgfscope}%
\pgfsetbuttcap%
\pgfsetroundjoin%
\definecolor{currentfill}{rgb}{0.000000,0.000000,0.000000}%
\pgfsetfillcolor{currentfill}%
\pgfsetlinewidth{0.602250pt}%
\definecolor{currentstroke}{rgb}{0.000000,0.000000,0.000000}%
\pgfsetstrokecolor{currentstroke}%
\pgfsetdash{}{0pt}%
\pgfsys@defobject{currentmarker}{\pgfqpoint{0.000000in}{-0.027778in}}{\pgfqpoint{0.000000in}{0.000000in}}{%
\pgfpathmoveto{\pgfqpoint{0.000000in}{0.000000in}}%
\pgfpathlineto{\pgfqpoint{0.000000in}{-0.027778in}}%
\pgfusepath{stroke,fill}%
}%
\begin{pgfscope}%
\pgfsys@transformshift{2.488249in}{0.668778in}%
\pgfsys@useobject{currentmarker}{}%
\end{pgfscope}%
\end{pgfscope}%
\begin{pgfscope}%
\pgfsetbuttcap%
\pgfsetroundjoin%
\definecolor{currentfill}{rgb}{0.000000,0.000000,0.000000}%
\pgfsetfillcolor{currentfill}%
\pgfsetlinewidth{0.602250pt}%
\definecolor{currentstroke}{rgb}{0.000000,0.000000,0.000000}%
\pgfsetstrokecolor{currentstroke}%
\pgfsetdash{}{0pt}%
\pgfsys@defobject{currentmarker}{\pgfqpoint{0.000000in}{-0.027778in}}{\pgfqpoint{0.000000in}{0.000000in}}{%
\pgfpathmoveto{\pgfqpoint{0.000000in}{0.000000in}}%
\pgfpathlineto{\pgfqpoint{0.000000in}{-0.027778in}}%
\pgfusepath{stroke,fill}%
}%
\begin{pgfscope}%
\pgfsys@transformshift{2.612894in}{0.668778in}%
\pgfsys@useobject{currentmarker}{}%
\end{pgfscope}%
\end{pgfscope}%
\begin{pgfscope}%
\pgfsetbuttcap%
\pgfsetroundjoin%
\definecolor{currentfill}{rgb}{0.000000,0.000000,0.000000}%
\pgfsetfillcolor{currentfill}%
\pgfsetlinewidth{0.602250pt}%
\definecolor{currentstroke}{rgb}{0.000000,0.000000,0.000000}%
\pgfsetstrokecolor{currentstroke}%
\pgfsetdash{}{0pt}%
\pgfsys@defobject{currentmarker}{\pgfqpoint{0.000000in}{-0.027778in}}{\pgfqpoint{0.000000in}{0.000000in}}{%
\pgfpathmoveto{\pgfqpoint{0.000000in}{0.000000in}}%
\pgfpathlineto{\pgfqpoint{0.000000in}{-0.027778in}}%
\pgfusepath{stroke,fill}%
}%
\begin{pgfscope}%
\pgfsys@transformshift{2.722838in}{0.668778in}%
\pgfsys@useobject{currentmarker}{}%
\end{pgfscope}%
\end{pgfscope}%
\begin{pgfscope}%
\pgfsetbuttcap%
\pgfsetroundjoin%
\definecolor{currentfill}{rgb}{0.000000,0.000000,0.000000}%
\pgfsetfillcolor{currentfill}%
\pgfsetlinewidth{0.602250pt}%
\definecolor{currentstroke}{rgb}{0.000000,0.000000,0.000000}%
\pgfsetstrokecolor{currentstroke}%
\pgfsetdash{}{0pt}%
\pgfsys@defobject{currentmarker}{\pgfqpoint{0.000000in}{-0.027778in}}{\pgfqpoint{0.000000in}{0.000000in}}{%
\pgfpathmoveto{\pgfqpoint{0.000000in}{0.000000in}}%
\pgfpathlineto{\pgfqpoint{0.000000in}{-0.027778in}}%
\pgfusepath{stroke,fill}%
}%
\begin{pgfscope}%
\pgfsys@transformshift{3.468205in}{0.668778in}%
\pgfsys@useobject{currentmarker}{}%
\end{pgfscope}%
\end{pgfscope}%
\begin{pgfscope}%
\pgfsetbuttcap%
\pgfsetroundjoin%
\definecolor{currentfill}{rgb}{0.000000,0.000000,0.000000}%
\pgfsetfillcolor{currentfill}%
\pgfsetlinewidth{0.602250pt}%
\definecolor{currentstroke}{rgb}{0.000000,0.000000,0.000000}%
\pgfsetstrokecolor{currentstroke}%
\pgfsetdash{}{0pt}%
\pgfsys@defobject{currentmarker}{\pgfqpoint{0.000000in}{-0.027778in}}{\pgfqpoint{0.000000in}{0.000000in}}{%
\pgfpathmoveto{\pgfqpoint{0.000000in}{0.000000in}}%
\pgfpathlineto{\pgfqpoint{0.000000in}{-0.027778in}}%
\pgfusepath{stroke,fill}%
}%
\begin{pgfscope}%
\pgfsys@transformshift{3.846687in}{0.668778in}%
\pgfsys@useobject{currentmarker}{}%
\end{pgfscope}%
\end{pgfscope}%
\begin{pgfscope}%
\pgfsetbuttcap%
\pgfsetroundjoin%
\definecolor{currentfill}{rgb}{0.000000,0.000000,0.000000}%
\pgfsetfillcolor{currentfill}%
\pgfsetlinewidth{0.602250pt}%
\definecolor{currentstroke}{rgb}{0.000000,0.000000,0.000000}%
\pgfsetstrokecolor{currentstroke}%
\pgfsetdash{}{0pt}%
\pgfsys@defobject{currentmarker}{\pgfqpoint{0.000000in}{-0.027778in}}{\pgfqpoint{0.000000in}{0.000000in}}{%
\pgfpathmoveto{\pgfqpoint{0.000000in}{0.000000in}}%
\pgfpathlineto{\pgfqpoint{0.000000in}{-0.027778in}}%
\pgfusepath{stroke,fill}%
}%
\begin{pgfscope}%
\pgfsys@transformshift{4.115224in}{0.668778in}%
\pgfsys@useobject{currentmarker}{}%
\end{pgfscope}%
\end{pgfscope}%
\begin{pgfscope}%
\pgfsetbuttcap%
\pgfsetroundjoin%
\definecolor{currentfill}{rgb}{0.000000,0.000000,0.000000}%
\pgfsetfillcolor{currentfill}%
\pgfsetlinewidth{0.602250pt}%
\definecolor{currentstroke}{rgb}{0.000000,0.000000,0.000000}%
\pgfsetstrokecolor{currentstroke}%
\pgfsetdash{}{0pt}%
\pgfsys@defobject{currentmarker}{\pgfqpoint{0.000000in}{-0.027778in}}{\pgfqpoint{0.000000in}{0.000000in}}{%
\pgfpathmoveto{\pgfqpoint{0.000000in}{0.000000in}}%
\pgfpathlineto{\pgfqpoint{0.000000in}{-0.027778in}}%
\pgfusepath{stroke,fill}%
}%
\begin{pgfscope}%
\pgfsys@transformshift{4.323517in}{0.668778in}%
\pgfsys@useobject{currentmarker}{}%
\end{pgfscope}%
\end{pgfscope}%
\begin{pgfscope}%
\pgfsetbuttcap%
\pgfsetroundjoin%
\definecolor{currentfill}{rgb}{0.000000,0.000000,0.000000}%
\pgfsetfillcolor{currentfill}%
\pgfsetlinewidth{0.602250pt}%
\definecolor{currentstroke}{rgb}{0.000000,0.000000,0.000000}%
\pgfsetstrokecolor{currentstroke}%
\pgfsetdash{}{0pt}%
\pgfsys@defobject{currentmarker}{\pgfqpoint{0.000000in}{-0.027778in}}{\pgfqpoint{0.000000in}{0.000000in}}{%
\pgfpathmoveto{\pgfqpoint{0.000000in}{0.000000in}}%
\pgfpathlineto{\pgfqpoint{0.000000in}{-0.027778in}}%
\pgfusepath{stroke,fill}%
}%
\begin{pgfscope}%
\pgfsys@transformshift{4.493705in}{0.668778in}%
\pgfsys@useobject{currentmarker}{}%
\end{pgfscope}%
\end{pgfscope}%
\begin{pgfscope}%
\pgfsetbuttcap%
\pgfsetroundjoin%
\definecolor{currentfill}{rgb}{0.000000,0.000000,0.000000}%
\pgfsetfillcolor{currentfill}%
\pgfsetlinewidth{0.602250pt}%
\definecolor{currentstroke}{rgb}{0.000000,0.000000,0.000000}%
\pgfsetstrokecolor{currentstroke}%
\pgfsetdash{}{0pt}%
\pgfsys@defobject{currentmarker}{\pgfqpoint{0.000000in}{-0.027778in}}{\pgfqpoint{0.000000in}{0.000000in}}{%
\pgfpathmoveto{\pgfqpoint{0.000000in}{0.000000in}}%
\pgfpathlineto{\pgfqpoint{0.000000in}{-0.027778in}}%
\pgfusepath{stroke,fill}%
}%
\begin{pgfscope}%
\pgfsys@transformshift{4.637597in}{0.668778in}%
\pgfsys@useobject{currentmarker}{}%
\end{pgfscope}%
\end{pgfscope}%
\begin{pgfscope}%
\pgfsetbuttcap%
\pgfsetroundjoin%
\definecolor{currentfill}{rgb}{0.000000,0.000000,0.000000}%
\pgfsetfillcolor{currentfill}%
\pgfsetlinewidth{0.602250pt}%
\definecolor{currentstroke}{rgb}{0.000000,0.000000,0.000000}%
\pgfsetstrokecolor{currentstroke}%
\pgfsetdash{}{0pt}%
\pgfsys@defobject{currentmarker}{\pgfqpoint{0.000000in}{-0.027778in}}{\pgfqpoint{0.000000in}{0.000000in}}{%
\pgfpathmoveto{\pgfqpoint{0.000000in}{0.000000in}}%
\pgfpathlineto{\pgfqpoint{0.000000in}{-0.027778in}}%
\pgfusepath{stroke,fill}%
}%
\begin{pgfscope}%
\pgfsys@transformshift{4.762242in}{0.668778in}%
\pgfsys@useobject{currentmarker}{}%
\end{pgfscope}%
\end{pgfscope}%
\begin{pgfscope}%
\pgfsetbuttcap%
\pgfsetroundjoin%
\definecolor{currentfill}{rgb}{0.000000,0.000000,0.000000}%
\pgfsetfillcolor{currentfill}%
\pgfsetlinewidth{0.602250pt}%
\definecolor{currentstroke}{rgb}{0.000000,0.000000,0.000000}%
\pgfsetstrokecolor{currentstroke}%
\pgfsetdash{}{0pt}%
\pgfsys@defobject{currentmarker}{\pgfqpoint{0.000000in}{-0.027778in}}{\pgfqpoint{0.000000in}{0.000000in}}{%
\pgfpathmoveto{\pgfqpoint{0.000000in}{0.000000in}}%
\pgfpathlineto{\pgfqpoint{0.000000in}{-0.027778in}}%
\pgfusepath{stroke,fill}%
}%
\begin{pgfscope}%
\pgfsys@transformshift{4.872187in}{0.668778in}%
\pgfsys@useobject{currentmarker}{}%
\end{pgfscope}%
\end{pgfscope}%
\begin{pgfscope}%
\pgfsetbuttcap%
\pgfsetroundjoin%
\definecolor{currentfill}{rgb}{0.000000,0.000000,0.000000}%
\pgfsetfillcolor{currentfill}%
\pgfsetlinewidth{0.602250pt}%
\definecolor{currentstroke}{rgb}{0.000000,0.000000,0.000000}%
\pgfsetstrokecolor{currentstroke}%
\pgfsetdash{}{0pt}%
\pgfsys@defobject{currentmarker}{\pgfqpoint{0.000000in}{-0.027778in}}{\pgfqpoint{0.000000in}{0.000000in}}{%
\pgfpathmoveto{\pgfqpoint{0.000000in}{0.000000in}}%
\pgfpathlineto{\pgfqpoint{0.000000in}{-0.027778in}}%
\pgfusepath{stroke,fill}%
}%
\begin{pgfscope}%
\pgfsys@transformshift{5.617554in}{0.668778in}%
\pgfsys@useobject{currentmarker}{}%
\end{pgfscope}%
\end{pgfscope}%
\begin{pgfscope}%
\pgfsetbuttcap%
\pgfsetroundjoin%
\definecolor{currentfill}{rgb}{0.000000,0.000000,0.000000}%
\pgfsetfillcolor{currentfill}%
\pgfsetlinewidth{0.602250pt}%
\definecolor{currentstroke}{rgb}{0.000000,0.000000,0.000000}%
\pgfsetstrokecolor{currentstroke}%
\pgfsetdash{}{0pt}%
\pgfsys@defobject{currentmarker}{\pgfqpoint{0.000000in}{-0.027778in}}{\pgfqpoint{0.000000in}{0.000000in}}{%
\pgfpathmoveto{\pgfqpoint{0.000000in}{0.000000in}}%
\pgfpathlineto{\pgfqpoint{0.000000in}{-0.027778in}}%
\pgfusepath{stroke,fill}%
}%
\begin{pgfscope}%
\pgfsys@transformshift{5.996035in}{0.668778in}%
\pgfsys@useobject{currentmarker}{}%
\end{pgfscope}%
\end{pgfscope}%
\begin{pgfscope}%
\definecolor{textcolor}{rgb}{0.000000,0.000000,0.000000}%
\pgfsetstrokecolor{textcolor}%
\pgfsetfillcolor{textcolor}%
\pgftext[x=3.456389in,y=0.381587in,,top]{\color{textcolor}\sffamily\fontsize{10.000000}{12.000000}\selectfont Distance from wall \(\displaystyle y\) (mm)}%
\end{pgfscope}%
\begin{pgfscope}%
\pgfsetbuttcap%
\pgfsetroundjoin%
\definecolor{currentfill}{rgb}{0.000000,0.000000,0.000000}%
\pgfsetfillcolor{currentfill}%
\pgfsetlinewidth{0.803000pt}%
\definecolor{currentstroke}{rgb}{0.000000,0.000000,0.000000}%
\pgfsetstrokecolor{currentstroke}%
\pgfsetdash{}{0pt}%
\pgfsys@defobject{currentmarker}{\pgfqpoint{-0.048611in}{0.000000in}}{\pgfqpoint{0.000000in}{0.000000in}}{%
\pgfpathmoveto{\pgfqpoint{0.000000in}{0.000000in}}%
\pgfpathlineto{\pgfqpoint{-0.048611in}{0.000000in}}%
\pgfusepath{stroke,fill}%
}%
\begin{pgfscope}%
\pgfsys@transformshift{0.662778in}{0.843608in}%
\pgfsys@useobject{currentmarker}{}%
\end{pgfscope}%
\end{pgfscope}%
\begin{pgfscope}%
\definecolor{textcolor}{rgb}{0.000000,0.000000,0.000000}%
\pgfsetstrokecolor{textcolor}%
\pgfsetfillcolor{textcolor}%
\pgftext[x=0.344676in,y=0.790846in,left,base]{\color{textcolor}\sffamily\fontsize{10.000000}{12.000000}\selectfont 0.2}%
\end{pgfscope}%
\begin{pgfscope}%
\pgfsetbuttcap%
\pgfsetroundjoin%
\definecolor{currentfill}{rgb}{0.000000,0.000000,0.000000}%
\pgfsetfillcolor{currentfill}%
\pgfsetlinewidth{0.803000pt}%
\definecolor{currentstroke}{rgb}{0.000000,0.000000,0.000000}%
\pgfsetstrokecolor{currentstroke}%
\pgfsetdash{}{0pt}%
\pgfsys@defobject{currentmarker}{\pgfqpoint{-0.048611in}{0.000000in}}{\pgfqpoint{0.000000in}{0.000000in}}{%
\pgfpathmoveto{\pgfqpoint{0.000000in}{0.000000in}}%
\pgfpathlineto{\pgfqpoint{-0.048611in}{0.000000in}}%
\pgfusepath{stroke,fill}%
}%
\begin{pgfscope}%
\pgfsys@transformshift{0.662778in}{1.672637in}%
\pgfsys@useobject{currentmarker}{}%
\end{pgfscope}%
\end{pgfscope}%
\begin{pgfscope}%
\definecolor{textcolor}{rgb}{0.000000,0.000000,0.000000}%
\pgfsetstrokecolor{textcolor}%
\pgfsetfillcolor{textcolor}%
\pgftext[x=0.344676in,y=1.619875in,left,base]{\color{textcolor}\sffamily\fontsize{10.000000}{12.000000}\selectfont 0.3}%
\end{pgfscope}%
\begin{pgfscope}%
\pgfsetbuttcap%
\pgfsetroundjoin%
\definecolor{currentfill}{rgb}{0.000000,0.000000,0.000000}%
\pgfsetfillcolor{currentfill}%
\pgfsetlinewidth{0.803000pt}%
\definecolor{currentstroke}{rgb}{0.000000,0.000000,0.000000}%
\pgfsetstrokecolor{currentstroke}%
\pgfsetdash{}{0pt}%
\pgfsys@defobject{currentmarker}{\pgfqpoint{-0.048611in}{0.000000in}}{\pgfqpoint{0.000000in}{0.000000in}}{%
\pgfpathmoveto{\pgfqpoint{0.000000in}{0.000000in}}%
\pgfpathlineto{\pgfqpoint{-0.048611in}{0.000000in}}%
\pgfusepath{stroke,fill}%
}%
\begin{pgfscope}%
\pgfsys@transformshift{0.662778in}{2.501666in}%
\pgfsys@useobject{currentmarker}{}%
\end{pgfscope}%
\end{pgfscope}%
\begin{pgfscope}%
\definecolor{textcolor}{rgb}{0.000000,0.000000,0.000000}%
\pgfsetstrokecolor{textcolor}%
\pgfsetfillcolor{textcolor}%
\pgftext[x=0.344676in,y=2.448905in,left,base]{\color{textcolor}\sffamily\fontsize{10.000000}{12.000000}\selectfont 0.4}%
\end{pgfscope}%
\begin{pgfscope}%
\pgfsetbuttcap%
\pgfsetroundjoin%
\definecolor{currentfill}{rgb}{0.000000,0.000000,0.000000}%
\pgfsetfillcolor{currentfill}%
\pgfsetlinewidth{0.803000pt}%
\definecolor{currentstroke}{rgb}{0.000000,0.000000,0.000000}%
\pgfsetstrokecolor{currentstroke}%
\pgfsetdash{}{0pt}%
\pgfsys@defobject{currentmarker}{\pgfqpoint{-0.048611in}{0.000000in}}{\pgfqpoint{0.000000in}{0.000000in}}{%
\pgfpathmoveto{\pgfqpoint{0.000000in}{0.000000in}}%
\pgfpathlineto{\pgfqpoint{-0.048611in}{0.000000in}}%
\pgfusepath{stroke,fill}%
}%
\begin{pgfscope}%
\pgfsys@transformshift{0.662778in}{3.330695in}%
\pgfsys@useobject{currentmarker}{}%
\end{pgfscope}%
\end{pgfscope}%
\begin{pgfscope}%
\definecolor{textcolor}{rgb}{0.000000,0.000000,0.000000}%
\pgfsetstrokecolor{textcolor}%
\pgfsetfillcolor{textcolor}%
\pgftext[x=0.344676in,y=3.277934in,left,base]{\color{textcolor}\sffamily\fontsize{10.000000}{12.000000}\selectfont 0.5}%
\end{pgfscope}%
\begin{pgfscope}%
\pgfsetbuttcap%
\pgfsetroundjoin%
\definecolor{currentfill}{rgb}{0.000000,0.000000,0.000000}%
\pgfsetfillcolor{currentfill}%
\pgfsetlinewidth{0.803000pt}%
\definecolor{currentstroke}{rgb}{0.000000,0.000000,0.000000}%
\pgfsetstrokecolor{currentstroke}%
\pgfsetdash{}{0pt}%
\pgfsys@defobject{currentmarker}{\pgfqpoint{-0.048611in}{0.000000in}}{\pgfqpoint{0.000000in}{0.000000in}}{%
\pgfpathmoveto{\pgfqpoint{0.000000in}{0.000000in}}%
\pgfpathlineto{\pgfqpoint{-0.048611in}{0.000000in}}%
\pgfusepath{stroke,fill}%
}%
\begin{pgfscope}%
\pgfsys@transformshift{0.662778in}{4.159725in}%
\pgfsys@useobject{currentmarker}{}%
\end{pgfscope}%
\end{pgfscope}%
\begin{pgfscope}%
\definecolor{textcolor}{rgb}{0.000000,0.000000,0.000000}%
\pgfsetstrokecolor{textcolor}%
\pgfsetfillcolor{textcolor}%
\pgftext[x=0.344676in,y=4.106963in,left,base]{\color{textcolor}\sffamily\fontsize{10.000000}{12.000000}\selectfont 0.6}%
\end{pgfscope}%
\begin{pgfscope}%
\definecolor{textcolor}{rgb}{0.000000,0.000000,0.000000}%
\pgfsetstrokecolor{textcolor}%
\pgfsetfillcolor{textcolor}%
\pgftext[x=0.289121in,y=2.659389in,,bottom,rotate=90.000000]{\color{textcolor}\sffamily\fontsize{10.000000}{12.000000}\selectfont u RMS (m/s)}%
\end{pgfscope}%
\begin{pgfscope}%
\pgfpathrectangle{\pgfqpoint{0.662778in}{0.668778in}}{\pgfqpoint{5.587222in}{3.981222in}}%
\pgfusepath{clip}%
\pgfsetbuttcap%
\pgfsetroundjoin%
\pgfsetlinewidth{1.505625pt}%
\definecolor{currentstroke}{rgb}{0.121569,0.466667,0.705882}%
\pgfsetstrokecolor{currentstroke}%
\pgfsetdash{{1.500000pt}{2.475000pt}}{0.000000pt}%
\pgfpathmoveto{\pgfqpoint{0.916742in}{4.310857in}}%
\pgfpathlineto{\pgfqpoint{1.050320in}{4.365158in}}%
\pgfpathlineto{\pgfqpoint{1.318857in}{4.455854in}}%
\pgfpathlineto{\pgfqpoint{1.527150in}{4.469035in}}%
\pgfpathlineto{\pgfqpoint{2.821187in}{4.075578in}}%
\pgfpathlineto{\pgfqpoint{3.468205in}{3.450324in}}%
\pgfpathlineto{\pgfqpoint{4.054980in}{2.960036in}}%
\pgfpathlineto{\pgfqpoint{4.412484in}{2.685296in}}%
\pgfpathlineto{\pgfqpoint{4.670353in}{2.540216in}}%
\pgfpathlineto{\pgfqpoint{4.872187in}{2.397623in}}%
\pgfpathlineto{\pgfqpoint{4.970535in}{2.316046in}}%
\pgfpathlineto{\pgfqpoint{5.349017in}{1.874754in}}%
\pgfpathlineto{\pgfqpoint{5.617554in}{1.450042in}}%
\pgfpathlineto{\pgfqpoint{5.825847in}{1.037932in}}%
\pgfpathlineto{\pgfqpoint{5.996035in}{0.849742in}}%
\pgfusepath{stroke}%
\end{pgfscope}%
\begin{pgfscope}%
\pgfpathrectangle{\pgfqpoint{0.662778in}{0.668778in}}{\pgfqpoint{5.587222in}{3.981222in}}%
\pgfusepath{clip}%
\pgfsetbuttcap%
\pgfsetroundjoin%
\definecolor{currentfill}{rgb}{0.121569,0.466667,0.705882}%
\pgfsetfillcolor{currentfill}%
\pgfsetlinewidth{1.003750pt}%
\definecolor{currentstroke}{rgb}{0.121569,0.466667,0.705882}%
\pgfsetstrokecolor{currentstroke}%
\pgfsetdash{}{0pt}%
\pgfsys@defobject{currentmarker}{\pgfqpoint{-0.034722in}{-0.034722in}}{\pgfqpoint{0.034722in}{0.034722in}}{%
\pgfpathmoveto{\pgfqpoint{0.000000in}{-0.034722in}}%
\pgfpathcurveto{\pgfqpoint{0.009208in}{-0.034722in}}{\pgfqpoint{0.018041in}{-0.031064in}}{\pgfqpoint{0.024552in}{-0.024552in}}%
\pgfpathcurveto{\pgfqpoint{0.031064in}{-0.018041in}}{\pgfqpoint{0.034722in}{-0.009208in}}{\pgfqpoint{0.034722in}{0.000000in}}%
\pgfpathcurveto{\pgfqpoint{0.034722in}{0.009208in}}{\pgfqpoint{0.031064in}{0.018041in}}{\pgfqpoint{0.024552in}{0.024552in}}%
\pgfpathcurveto{\pgfqpoint{0.018041in}{0.031064in}}{\pgfqpoint{0.009208in}{0.034722in}}{\pgfqpoint{0.000000in}{0.034722in}}%
\pgfpathcurveto{\pgfqpoint{-0.009208in}{0.034722in}}{\pgfqpoint{-0.018041in}{0.031064in}}{\pgfqpoint{-0.024552in}{0.024552in}}%
\pgfpathcurveto{\pgfqpoint{-0.031064in}{0.018041in}}{\pgfqpoint{-0.034722in}{0.009208in}}{\pgfqpoint{-0.034722in}{0.000000in}}%
\pgfpathcurveto{\pgfqpoint{-0.034722in}{-0.009208in}}{\pgfqpoint{-0.031064in}{-0.018041in}}{\pgfqpoint{-0.024552in}{-0.024552in}}%
\pgfpathcurveto{\pgfqpoint{-0.018041in}{-0.031064in}}{\pgfqpoint{-0.009208in}{-0.034722in}}{\pgfqpoint{0.000000in}{-0.034722in}}%
\pgfpathclose%
\pgfusepath{stroke,fill}%
}%
\begin{pgfscope}%
\pgfsys@transformshift{0.916742in}{4.310857in}%
\pgfsys@useobject{currentmarker}{}%
\end{pgfscope}%
\begin{pgfscope}%
\pgfsys@transformshift{1.050320in}{4.365158in}%
\pgfsys@useobject{currentmarker}{}%
\end{pgfscope}%
\begin{pgfscope}%
\pgfsys@transformshift{1.318857in}{4.455854in}%
\pgfsys@useobject{currentmarker}{}%
\end{pgfscope}%
\begin{pgfscope}%
\pgfsys@transformshift{1.527150in}{4.469035in}%
\pgfsys@useobject{currentmarker}{}%
\end{pgfscope}%
\begin{pgfscope}%
\pgfsys@transformshift{2.821187in}{4.075578in}%
\pgfsys@useobject{currentmarker}{}%
\end{pgfscope}%
\begin{pgfscope}%
\pgfsys@transformshift{3.468205in}{3.450324in}%
\pgfsys@useobject{currentmarker}{}%
\end{pgfscope}%
\begin{pgfscope}%
\pgfsys@transformshift{4.054980in}{2.960036in}%
\pgfsys@useobject{currentmarker}{}%
\end{pgfscope}%
\begin{pgfscope}%
\pgfsys@transformshift{4.412484in}{2.685296in}%
\pgfsys@useobject{currentmarker}{}%
\end{pgfscope}%
\begin{pgfscope}%
\pgfsys@transformshift{4.670353in}{2.540216in}%
\pgfsys@useobject{currentmarker}{}%
\end{pgfscope}%
\begin{pgfscope}%
\pgfsys@transformshift{4.872187in}{2.397623in}%
\pgfsys@useobject{currentmarker}{}%
\end{pgfscope}%
\begin{pgfscope}%
\pgfsys@transformshift{4.970535in}{2.316046in}%
\pgfsys@useobject{currentmarker}{}%
\end{pgfscope}%
\begin{pgfscope}%
\pgfsys@transformshift{5.349017in}{1.874754in}%
\pgfsys@useobject{currentmarker}{}%
\end{pgfscope}%
\begin{pgfscope}%
\pgfsys@transformshift{5.617554in}{1.450042in}%
\pgfsys@useobject{currentmarker}{}%
\end{pgfscope}%
\begin{pgfscope}%
\pgfsys@transformshift{5.825847in}{1.037932in}%
\pgfsys@useobject{currentmarker}{}%
\end{pgfscope}%
\begin{pgfscope}%
\pgfsys@transformshift{5.996035in}{0.849742in}%
\pgfsys@useobject{currentmarker}{}%
\end{pgfscope}%
\end{pgfscope}%
\begin{pgfscope}%
\pgfsetrectcap%
\pgfsetmiterjoin%
\pgfsetlinewidth{0.803000pt}%
\definecolor{currentstroke}{rgb}{0.000000,0.000000,0.000000}%
\pgfsetstrokecolor{currentstroke}%
\pgfsetdash{}{0pt}%
\pgfpathmoveto{\pgfqpoint{0.662778in}{0.668778in}}%
\pgfpathlineto{\pgfqpoint{0.662778in}{4.650000in}}%
\pgfusepath{stroke}%
\end{pgfscope}%
\begin{pgfscope}%
\pgfsetrectcap%
\pgfsetmiterjoin%
\pgfsetlinewidth{0.803000pt}%
\definecolor{currentstroke}{rgb}{0.000000,0.000000,0.000000}%
\pgfsetstrokecolor{currentstroke}%
\pgfsetdash{}{0pt}%
\pgfpathmoveto{\pgfqpoint{6.250000in}{0.668778in}}%
\pgfpathlineto{\pgfqpoint{6.250000in}{4.650000in}}%
\pgfusepath{stroke}%
\end{pgfscope}%
\begin{pgfscope}%
\pgfsetrectcap%
\pgfsetmiterjoin%
\pgfsetlinewidth{0.803000pt}%
\definecolor{currentstroke}{rgb}{0.000000,0.000000,0.000000}%
\pgfsetstrokecolor{currentstroke}%
\pgfsetdash{}{0pt}%
\pgfpathmoveto{\pgfqpoint{0.662778in}{0.668778in}}%
\pgfpathlineto{\pgfqpoint{6.250000in}{0.668778in}}%
\pgfusepath{stroke}%
\end{pgfscope}%
\begin{pgfscope}%
\pgfsetrectcap%
\pgfsetmiterjoin%
\pgfsetlinewidth{0.803000pt}%
\definecolor{currentstroke}{rgb}{0.000000,0.000000,0.000000}%
\pgfsetstrokecolor{currentstroke}%
\pgfsetdash{}{0pt}%
\pgfpathmoveto{\pgfqpoint{0.662778in}{4.650000in}}%
\pgfpathlineto{\pgfqpoint{6.250000in}{4.650000in}}%
\pgfusepath{stroke}%
\end{pgfscope}%
\end{pgfpicture}%
\makeatother%
\endgroup%
}
    \caption{RMS Profile plotted in log-lin scale}
    \label{fig:RMS_Profile}
\end{figure}

\begin{figure}
    \centering
        \resizebox{0.8\linewidth}{!}{%% Creator: Matplotlib, PGF backend
%%
%% To include the figure in your LaTeX document, write
%%   \input{<filename>.pgf}
%%
%% Make sure the required packages are loaded in your preamble
%%   \usepackage{pgf}
%%
%% Figures using additional raster images can only be included by \input if
%% they are in the same directory as the main LaTeX file. For loading figures
%% from other directories you can use the `import` package
%%   \usepackage{import}
%% and then include the figures with
%%   \import{<path to file>}{<filename>.pgf}
%%
%% Matplotlib used the following preamble
%%   \usepackage{fontspec}
%%   \setmainfont{DejaVuSerif.ttf}[Path=C:/ProgramData/Anaconda3/lib/site-packages/matplotlib/mpl-data/fonts/ttf/]
%%   \setsansfont{DejaVuSans.ttf}[Path=C:/ProgramData/Anaconda3/lib/site-packages/matplotlib/mpl-data/fonts/ttf/]
%%   \setmonofont{DejaVuSansMono.ttf}[Path=C:/ProgramData/Anaconda3/lib/site-packages/matplotlib/mpl-data/fonts/ttf/]
%%
\begingroup%
\makeatletter%
\begin{pgfpicture}%
\pgfpathrectangle{\pgfpointorigin}{\pgfqpoint{6.400000in}{4.800000in}}%
\pgfusepath{use as bounding box, clip}%
\begin{pgfscope}%
\pgfsetbuttcap%
\pgfsetmiterjoin%
\definecolor{currentfill}{rgb}{1.000000,1.000000,1.000000}%
\pgfsetfillcolor{currentfill}%
\pgfsetlinewidth{0.000000pt}%
\definecolor{currentstroke}{rgb}{1.000000,1.000000,1.000000}%
\pgfsetstrokecolor{currentstroke}%
\pgfsetdash{}{0pt}%
\pgfpathmoveto{\pgfqpoint{0.000000in}{0.000000in}}%
\pgfpathlineto{\pgfqpoint{6.400000in}{0.000000in}}%
\pgfpathlineto{\pgfqpoint{6.400000in}{4.800000in}}%
\pgfpathlineto{\pgfqpoint{0.000000in}{4.800000in}}%
\pgfpathclose%
\pgfusepath{fill}%
\end{pgfscope}%
\begin{pgfscope}%
\pgfsetbuttcap%
\pgfsetmiterjoin%
\definecolor{currentfill}{rgb}{1.000000,1.000000,1.000000}%
\pgfsetfillcolor{currentfill}%
\pgfsetlinewidth{0.000000pt}%
\definecolor{currentstroke}{rgb}{0.000000,0.000000,0.000000}%
\pgfsetstrokecolor{currentstroke}%
\pgfsetstrokeopacity{0.000000}%
\pgfsetdash{}{0pt}%
\pgfpathmoveto{\pgfqpoint{0.813028in}{0.668778in}}%
\pgfpathlineto{\pgfqpoint{6.250000in}{0.668778in}}%
\pgfpathlineto{\pgfqpoint{6.250000in}{4.650000in}}%
\pgfpathlineto{\pgfqpoint{0.813028in}{4.650000in}}%
\pgfpathclose%
\pgfusepath{fill}%
\end{pgfscope}%
\begin{pgfscope}%
\pgfsetbuttcap%
\pgfsetroundjoin%
\definecolor{currentfill}{rgb}{0.000000,0.000000,0.000000}%
\pgfsetfillcolor{currentfill}%
\pgfsetlinewidth{0.803000pt}%
\definecolor{currentstroke}{rgb}{0.000000,0.000000,0.000000}%
\pgfsetstrokecolor{currentstroke}%
\pgfsetdash{}{0pt}%
\pgfsys@defobject{currentmarker}{\pgfqpoint{0.000000in}{-0.048611in}}{\pgfqpoint{0.000000in}{0.000000in}}{%
\pgfpathmoveto{\pgfqpoint{0.000000in}{0.000000in}}%
\pgfpathlineto{\pgfqpoint{0.000000in}{-0.048611in}}%
\pgfusepath{stroke,fill}%
}%
\begin{pgfscope}%
\pgfsys@transformshift{0.821845in}{0.668778in}%
\pgfsys@useobject{currentmarker}{}%
\end{pgfscope}%
\end{pgfscope}%
\begin{pgfscope}%
\definecolor{textcolor}{rgb}{0.000000,0.000000,0.000000}%
\pgfsetstrokecolor{textcolor}%
\pgfsetfillcolor{textcolor}%
\pgftext[x=0.821845in,y=0.571556in,,top]{\color{textcolor}\sffamily\fontsize{10.000000}{12.000000}\selectfont \(\displaystyle {10^{-1}}\)}%
\end{pgfscope}%
\begin{pgfscope}%
\pgfsetbuttcap%
\pgfsetroundjoin%
\definecolor{currentfill}{rgb}{0.000000,0.000000,0.000000}%
\pgfsetfillcolor{currentfill}%
\pgfsetlinewidth{0.803000pt}%
\definecolor{currentstroke}{rgb}{0.000000,0.000000,0.000000}%
\pgfsetstrokecolor{currentstroke}%
\pgfsetdash{}{0pt}%
\pgfsys@defobject{currentmarker}{\pgfqpoint{0.000000in}{-0.048611in}}{\pgfqpoint{0.000000in}{0.000000in}}{%
\pgfpathmoveto{\pgfqpoint{0.000000in}{0.000000in}}%
\pgfpathlineto{\pgfqpoint{0.000000in}{-0.048611in}}%
\pgfusepath{stroke,fill}%
}%
\begin{pgfscope}%
\pgfsys@transformshift{2.913394in}{0.668778in}%
\pgfsys@useobject{currentmarker}{}%
\end{pgfscope}%
\end{pgfscope}%
\begin{pgfscope}%
\definecolor{textcolor}{rgb}{0.000000,0.000000,0.000000}%
\pgfsetstrokecolor{textcolor}%
\pgfsetfillcolor{textcolor}%
\pgftext[x=2.913394in,y=0.571556in,,top]{\color{textcolor}\sffamily\fontsize{10.000000}{12.000000}\selectfont \(\displaystyle {10^{0}}\)}%
\end{pgfscope}%
\begin{pgfscope}%
\pgfsetbuttcap%
\pgfsetroundjoin%
\definecolor{currentfill}{rgb}{0.000000,0.000000,0.000000}%
\pgfsetfillcolor{currentfill}%
\pgfsetlinewidth{0.803000pt}%
\definecolor{currentstroke}{rgb}{0.000000,0.000000,0.000000}%
\pgfsetstrokecolor{currentstroke}%
\pgfsetdash{}{0pt}%
\pgfsys@defobject{currentmarker}{\pgfqpoint{0.000000in}{-0.048611in}}{\pgfqpoint{0.000000in}{0.000000in}}{%
\pgfpathmoveto{\pgfqpoint{0.000000in}{0.000000in}}%
\pgfpathlineto{\pgfqpoint{0.000000in}{-0.048611in}}%
\pgfusepath{stroke,fill}%
}%
\begin{pgfscope}%
\pgfsys@transformshift{5.004942in}{0.668778in}%
\pgfsys@useobject{currentmarker}{}%
\end{pgfscope}%
\end{pgfscope}%
\begin{pgfscope}%
\definecolor{textcolor}{rgb}{0.000000,0.000000,0.000000}%
\pgfsetstrokecolor{textcolor}%
\pgfsetfillcolor{textcolor}%
\pgftext[x=5.004942in,y=0.571556in,,top]{\color{textcolor}\sffamily\fontsize{10.000000}{12.000000}\selectfont \(\displaystyle {10^{1}}\)}%
\end{pgfscope}%
\begin{pgfscope}%
\pgfsetbuttcap%
\pgfsetroundjoin%
\definecolor{currentfill}{rgb}{0.000000,0.000000,0.000000}%
\pgfsetfillcolor{currentfill}%
\pgfsetlinewidth{0.602250pt}%
\definecolor{currentstroke}{rgb}{0.000000,0.000000,0.000000}%
\pgfsetstrokecolor{currentstroke}%
\pgfsetdash{}{0pt}%
\pgfsys@defobject{currentmarker}{\pgfqpoint{0.000000in}{-0.027778in}}{\pgfqpoint{0.000000in}{0.000000in}}{%
\pgfpathmoveto{\pgfqpoint{0.000000in}{0.000000in}}%
\pgfpathlineto{\pgfqpoint{0.000000in}{-0.027778in}}%
\pgfusepath{stroke,fill}%
}%
\begin{pgfscope}%
\pgfsys@transformshift{1.451464in}{0.668778in}%
\pgfsys@useobject{currentmarker}{}%
\end{pgfscope}%
\end{pgfscope}%
\begin{pgfscope}%
\pgfsetbuttcap%
\pgfsetroundjoin%
\definecolor{currentfill}{rgb}{0.000000,0.000000,0.000000}%
\pgfsetfillcolor{currentfill}%
\pgfsetlinewidth{0.602250pt}%
\definecolor{currentstroke}{rgb}{0.000000,0.000000,0.000000}%
\pgfsetstrokecolor{currentstroke}%
\pgfsetdash{}{0pt}%
\pgfsys@defobject{currentmarker}{\pgfqpoint{0.000000in}{-0.027778in}}{\pgfqpoint{0.000000in}{0.000000in}}{%
\pgfpathmoveto{\pgfqpoint{0.000000in}{0.000000in}}%
\pgfpathlineto{\pgfqpoint{0.000000in}{-0.027778in}}%
\pgfusepath{stroke,fill}%
}%
\begin{pgfscope}%
\pgfsys@transformshift{1.819767in}{0.668778in}%
\pgfsys@useobject{currentmarker}{}%
\end{pgfscope}%
\end{pgfscope}%
\begin{pgfscope}%
\pgfsetbuttcap%
\pgfsetroundjoin%
\definecolor{currentfill}{rgb}{0.000000,0.000000,0.000000}%
\pgfsetfillcolor{currentfill}%
\pgfsetlinewidth{0.602250pt}%
\definecolor{currentstroke}{rgb}{0.000000,0.000000,0.000000}%
\pgfsetstrokecolor{currentstroke}%
\pgfsetdash{}{0pt}%
\pgfsys@defobject{currentmarker}{\pgfqpoint{0.000000in}{-0.027778in}}{\pgfqpoint{0.000000in}{0.000000in}}{%
\pgfpathmoveto{\pgfqpoint{0.000000in}{0.000000in}}%
\pgfpathlineto{\pgfqpoint{0.000000in}{-0.027778in}}%
\pgfusepath{stroke,fill}%
}%
\begin{pgfscope}%
\pgfsys@transformshift{2.081083in}{0.668778in}%
\pgfsys@useobject{currentmarker}{}%
\end{pgfscope}%
\end{pgfscope}%
\begin{pgfscope}%
\pgfsetbuttcap%
\pgfsetroundjoin%
\definecolor{currentfill}{rgb}{0.000000,0.000000,0.000000}%
\pgfsetfillcolor{currentfill}%
\pgfsetlinewidth{0.602250pt}%
\definecolor{currentstroke}{rgb}{0.000000,0.000000,0.000000}%
\pgfsetstrokecolor{currentstroke}%
\pgfsetdash{}{0pt}%
\pgfsys@defobject{currentmarker}{\pgfqpoint{0.000000in}{-0.027778in}}{\pgfqpoint{0.000000in}{0.000000in}}{%
\pgfpathmoveto{\pgfqpoint{0.000000in}{0.000000in}}%
\pgfpathlineto{\pgfqpoint{0.000000in}{-0.027778in}}%
\pgfusepath{stroke,fill}%
}%
\begin{pgfscope}%
\pgfsys@transformshift{2.283775in}{0.668778in}%
\pgfsys@useobject{currentmarker}{}%
\end{pgfscope}%
\end{pgfscope}%
\begin{pgfscope}%
\pgfsetbuttcap%
\pgfsetroundjoin%
\definecolor{currentfill}{rgb}{0.000000,0.000000,0.000000}%
\pgfsetfillcolor{currentfill}%
\pgfsetlinewidth{0.602250pt}%
\definecolor{currentstroke}{rgb}{0.000000,0.000000,0.000000}%
\pgfsetstrokecolor{currentstroke}%
\pgfsetdash{}{0pt}%
\pgfsys@defobject{currentmarker}{\pgfqpoint{0.000000in}{-0.027778in}}{\pgfqpoint{0.000000in}{0.000000in}}{%
\pgfpathmoveto{\pgfqpoint{0.000000in}{0.000000in}}%
\pgfpathlineto{\pgfqpoint{0.000000in}{-0.027778in}}%
\pgfusepath{stroke,fill}%
}%
\begin{pgfscope}%
\pgfsys@transformshift{2.449386in}{0.668778in}%
\pgfsys@useobject{currentmarker}{}%
\end{pgfscope}%
\end{pgfscope}%
\begin{pgfscope}%
\pgfsetbuttcap%
\pgfsetroundjoin%
\definecolor{currentfill}{rgb}{0.000000,0.000000,0.000000}%
\pgfsetfillcolor{currentfill}%
\pgfsetlinewidth{0.602250pt}%
\definecolor{currentstroke}{rgb}{0.000000,0.000000,0.000000}%
\pgfsetstrokecolor{currentstroke}%
\pgfsetdash{}{0pt}%
\pgfsys@defobject{currentmarker}{\pgfqpoint{0.000000in}{-0.027778in}}{\pgfqpoint{0.000000in}{0.000000in}}{%
\pgfpathmoveto{\pgfqpoint{0.000000in}{0.000000in}}%
\pgfpathlineto{\pgfqpoint{0.000000in}{-0.027778in}}%
\pgfusepath{stroke,fill}%
}%
\begin{pgfscope}%
\pgfsys@transformshift{2.589409in}{0.668778in}%
\pgfsys@useobject{currentmarker}{}%
\end{pgfscope}%
\end{pgfscope}%
\begin{pgfscope}%
\pgfsetbuttcap%
\pgfsetroundjoin%
\definecolor{currentfill}{rgb}{0.000000,0.000000,0.000000}%
\pgfsetfillcolor{currentfill}%
\pgfsetlinewidth{0.602250pt}%
\definecolor{currentstroke}{rgb}{0.000000,0.000000,0.000000}%
\pgfsetstrokecolor{currentstroke}%
\pgfsetdash{}{0pt}%
\pgfsys@defobject{currentmarker}{\pgfqpoint{0.000000in}{-0.027778in}}{\pgfqpoint{0.000000in}{0.000000in}}{%
\pgfpathmoveto{\pgfqpoint{0.000000in}{0.000000in}}%
\pgfpathlineto{\pgfqpoint{0.000000in}{-0.027778in}}%
\pgfusepath{stroke,fill}%
}%
\begin{pgfscope}%
\pgfsys@transformshift{2.710702in}{0.668778in}%
\pgfsys@useobject{currentmarker}{}%
\end{pgfscope}%
\end{pgfscope}%
\begin{pgfscope}%
\pgfsetbuttcap%
\pgfsetroundjoin%
\definecolor{currentfill}{rgb}{0.000000,0.000000,0.000000}%
\pgfsetfillcolor{currentfill}%
\pgfsetlinewidth{0.602250pt}%
\definecolor{currentstroke}{rgb}{0.000000,0.000000,0.000000}%
\pgfsetstrokecolor{currentstroke}%
\pgfsetdash{}{0pt}%
\pgfsys@defobject{currentmarker}{\pgfqpoint{0.000000in}{-0.027778in}}{\pgfqpoint{0.000000in}{0.000000in}}{%
\pgfpathmoveto{\pgfqpoint{0.000000in}{0.000000in}}%
\pgfpathlineto{\pgfqpoint{0.000000in}{-0.027778in}}%
\pgfusepath{stroke,fill}%
}%
\begin{pgfscope}%
\pgfsys@transformshift{2.817690in}{0.668778in}%
\pgfsys@useobject{currentmarker}{}%
\end{pgfscope}%
\end{pgfscope}%
\begin{pgfscope}%
\pgfsetbuttcap%
\pgfsetroundjoin%
\definecolor{currentfill}{rgb}{0.000000,0.000000,0.000000}%
\pgfsetfillcolor{currentfill}%
\pgfsetlinewidth{0.602250pt}%
\definecolor{currentstroke}{rgb}{0.000000,0.000000,0.000000}%
\pgfsetstrokecolor{currentstroke}%
\pgfsetdash{}{0pt}%
\pgfsys@defobject{currentmarker}{\pgfqpoint{0.000000in}{-0.027778in}}{\pgfqpoint{0.000000in}{0.000000in}}{%
\pgfpathmoveto{\pgfqpoint{0.000000in}{0.000000in}}%
\pgfpathlineto{\pgfqpoint{0.000000in}{-0.027778in}}%
\pgfusepath{stroke,fill}%
}%
\begin{pgfscope}%
\pgfsys@transformshift{3.543013in}{0.668778in}%
\pgfsys@useobject{currentmarker}{}%
\end{pgfscope}%
\end{pgfscope}%
\begin{pgfscope}%
\pgfsetbuttcap%
\pgfsetroundjoin%
\definecolor{currentfill}{rgb}{0.000000,0.000000,0.000000}%
\pgfsetfillcolor{currentfill}%
\pgfsetlinewidth{0.602250pt}%
\definecolor{currentstroke}{rgb}{0.000000,0.000000,0.000000}%
\pgfsetstrokecolor{currentstroke}%
\pgfsetdash{}{0pt}%
\pgfsys@defobject{currentmarker}{\pgfqpoint{0.000000in}{-0.027778in}}{\pgfqpoint{0.000000in}{0.000000in}}{%
\pgfpathmoveto{\pgfqpoint{0.000000in}{0.000000in}}%
\pgfpathlineto{\pgfqpoint{0.000000in}{-0.027778in}}%
\pgfusepath{stroke,fill}%
}%
\begin{pgfscope}%
\pgfsys@transformshift{3.911316in}{0.668778in}%
\pgfsys@useobject{currentmarker}{}%
\end{pgfscope}%
\end{pgfscope}%
\begin{pgfscope}%
\pgfsetbuttcap%
\pgfsetroundjoin%
\definecolor{currentfill}{rgb}{0.000000,0.000000,0.000000}%
\pgfsetfillcolor{currentfill}%
\pgfsetlinewidth{0.602250pt}%
\definecolor{currentstroke}{rgb}{0.000000,0.000000,0.000000}%
\pgfsetstrokecolor{currentstroke}%
\pgfsetdash{}{0pt}%
\pgfsys@defobject{currentmarker}{\pgfqpoint{0.000000in}{-0.027778in}}{\pgfqpoint{0.000000in}{0.000000in}}{%
\pgfpathmoveto{\pgfqpoint{0.000000in}{0.000000in}}%
\pgfpathlineto{\pgfqpoint{0.000000in}{-0.027778in}}%
\pgfusepath{stroke,fill}%
}%
\begin{pgfscope}%
\pgfsys@transformshift{4.172632in}{0.668778in}%
\pgfsys@useobject{currentmarker}{}%
\end{pgfscope}%
\end{pgfscope}%
\begin{pgfscope}%
\pgfsetbuttcap%
\pgfsetroundjoin%
\definecolor{currentfill}{rgb}{0.000000,0.000000,0.000000}%
\pgfsetfillcolor{currentfill}%
\pgfsetlinewidth{0.602250pt}%
\definecolor{currentstroke}{rgb}{0.000000,0.000000,0.000000}%
\pgfsetstrokecolor{currentstroke}%
\pgfsetdash{}{0pt}%
\pgfsys@defobject{currentmarker}{\pgfqpoint{0.000000in}{-0.027778in}}{\pgfqpoint{0.000000in}{0.000000in}}{%
\pgfpathmoveto{\pgfqpoint{0.000000in}{0.000000in}}%
\pgfpathlineto{\pgfqpoint{0.000000in}{-0.027778in}}%
\pgfusepath{stroke,fill}%
}%
\begin{pgfscope}%
\pgfsys@transformshift{4.375324in}{0.668778in}%
\pgfsys@useobject{currentmarker}{}%
\end{pgfscope}%
\end{pgfscope}%
\begin{pgfscope}%
\pgfsetbuttcap%
\pgfsetroundjoin%
\definecolor{currentfill}{rgb}{0.000000,0.000000,0.000000}%
\pgfsetfillcolor{currentfill}%
\pgfsetlinewidth{0.602250pt}%
\definecolor{currentstroke}{rgb}{0.000000,0.000000,0.000000}%
\pgfsetstrokecolor{currentstroke}%
\pgfsetdash{}{0pt}%
\pgfsys@defobject{currentmarker}{\pgfqpoint{0.000000in}{-0.027778in}}{\pgfqpoint{0.000000in}{0.000000in}}{%
\pgfpathmoveto{\pgfqpoint{0.000000in}{0.000000in}}%
\pgfpathlineto{\pgfqpoint{0.000000in}{-0.027778in}}%
\pgfusepath{stroke,fill}%
}%
\begin{pgfscope}%
\pgfsys@transformshift{4.540935in}{0.668778in}%
\pgfsys@useobject{currentmarker}{}%
\end{pgfscope}%
\end{pgfscope}%
\begin{pgfscope}%
\pgfsetbuttcap%
\pgfsetroundjoin%
\definecolor{currentfill}{rgb}{0.000000,0.000000,0.000000}%
\pgfsetfillcolor{currentfill}%
\pgfsetlinewidth{0.602250pt}%
\definecolor{currentstroke}{rgb}{0.000000,0.000000,0.000000}%
\pgfsetstrokecolor{currentstroke}%
\pgfsetdash{}{0pt}%
\pgfsys@defobject{currentmarker}{\pgfqpoint{0.000000in}{-0.027778in}}{\pgfqpoint{0.000000in}{0.000000in}}{%
\pgfpathmoveto{\pgfqpoint{0.000000in}{0.000000in}}%
\pgfpathlineto{\pgfqpoint{0.000000in}{-0.027778in}}%
\pgfusepath{stroke,fill}%
}%
\begin{pgfscope}%
\pgfsys@transformshift{4.680957in}{0.668778in}%
\pgfsys@useobject{currentmarker}{}%
\end{pgfscope}%
\end{pgfscope}%
\begin{pgfscope}%
\pgfsetbuttcap%
\pgfsetroundjoin%
\definecolor{currentfill}{rgb}{0.000000,0.000000,0.000000}%
\pgfsetfillcolor{currentfill}%
\pgfsetlinewidth{0.602250pt}%
\definecolor{currentstroke}{rgb}{0.000000,0.000000,0.000000}%
\pgfsetstrokecolor{currentstroke}%
\pgfsetdash{}{0pt}%
\pgfsys@defobject{currentmarker}{\pgfqpoint{0.000000in}{-0.027778in}}{\pgfqpoint{0.000000in}{0.000000in}}{%
\pgfpathmoveto{\pgfqpoint{0.000000in}{0.000000in}}%
\pgfpathlineto{\pgfqpoint{0.000000in}{-0.027778in}}%
\pgfusepath{stroke,fill}%
}%
\begin{pgfscope}%
\pgfsys@transformshift{4.802250in}{0.668778in}%
\pgfsys@useobject{currentmarker}{}%
\end{pgfscope}%
\end{pgfscope}%
\begin{pgfscope}%
\pgfsetbuttcap%
\pgfsetroundjoin%
\definecolor{currentfill}{rgb}{0.000000,0.000000,0.000000}%
\pgfsetfillcolor{currentfill}%
\pgfsetlinewidth{0.602250pt}%
\definecolor{currentstroke}{rgb}{0.000000,0.000000,0.000000}%
\pgfsetstrokecolor{currentstroke}%
\pgfsetdash{}{0pt}%
\pgfsys@defobject{currentmarker}{\pgfqpoint{0.000000in}{-0.027778in}}{\pgfqpoint{0.000000in}{0.000000in}}{%
\pgfpathmoveto{\pgfqpoint{0.000000in}{0.000000in}}%
\pgfpathlineto{\pgfqpoint{0.000000in}{-0.027778in}}%
\pgfusepath{stroke,fill}%
}%
\begin{pgfscope}%
\pgfsys@transformshift{4.909238in}{0.668778in}%
\pgfsys@useobject{currentmarker}{}%
\end{pgfscope}%
\end{pgfscope}%
\begin{pgfscope}%
\pgfsetbuttcap%
\pgfsetroundjoin%
\definecolor{currentfill}{rgb}{0.000000,0.000000,0.000000}%
\pgfsetfillcolor{currentfill}%
\pgfsetlinewidth{0.602250pt}%
\definecolor{currentstroke}{rgb}{0.000000,0.000000,0.000000}%
\pgfsetstrokecolor{currentstroke}%
\pgfsetdash{}{0pt}%
\pgfsys@defobject{currentmarker}{\pgfqpoint{0.000000in}{-0.027778in}}{\pgfqpoint{0.000000in}{0.000000in}}{%
\pgfpathmoveto{\pgfqpoint{0.000000in}{0.000000in}}%
\pgfpathlineto{\pgfqpoint{0.000000in}{-0.027778in}}%
\pgfusepath{stroke,fill}%
}%
\begin{pgfscope}%
\pgfsys@transformshift{5.634561in}{0.668778in}%
\pgfsys@useobject{currentmarker}{}%
\end{pgfscope}%
\end{pgfscope}%
\begin{pgfscope}%
\pgfsetbuttcap%
\pgfsetroundjoin%
\definecolor{currentfill}{rgb}{0.000000,0.000000,0.000000}%
\pgfsetfillcolor{currentfill}%
\pgfsetlinewidth{0.602250pt}%
\definecolor{currentstroke}{rgb}{0.000000,0.000000,0.000000}%
\pgfsetstrokecolor{currentstroke}%
\pgfsetdash{}{0pt}%
\pgfsys@defobject{currentmarker}{\pgfqpoint{0.000000in}{-0.027778in}}{\pgfqpoint{0.000000in}{0.000000in}}{%
\pgfpathmoveto{\pgfqpoint{0.000000in}{0.000000in}}%
\pgfpathlineto{\pgfqpoint{0.000000in}{-0.027778in}}%
\pgfusepath{stroke,fill}%
}%
\begin{pgfscope}%
\pgfsys@transformshift{6.002865in}{0.668778in}%
\pgfsys@useobject{currentmarker}{}%
\end{pgfscope}%
\end{pgfscope}%
\begin{pgfscope}%
\definecolor{textcolor}{rgb}{0.000000,0.000000,0.000000}%
\pgfsetstrokecolor{textcolor}%
\pgfsetfillcolor{textcolor}%
\pgftext[x=3.531514in,y=0.381587in,,top]{\color{textcolor}\sffamily\fontsize{10.000000}{12.000000}\selectfont Distance from wall \(\displaystyle y\) (mm)}%
\end{pgfscope}%
\begin{pgfscope}%
\pgfsetbuttcap%
\pgfsetroundjoin%
\definecolor{currentfill}{rgb}{0.000000,0.000000,0.000000}%
\pgfsetfillcolor{currentfill}%
\pgfsetlinewidth{0.803000pt}%
\definecolor{currentstroke}{rgb}{0.000000,0.000000,0.000000}%
\pgfsetstrokecolor{currentstroke}%
\pgfsetdash{}{0pt}%
\pgfsys@defobject{currentmarker}{\pgfqpoint{-0.048611in}{0.000000in}}{\pgfqpoint{0.000000in}{0.000000in}}{%
\pgfpathmoveto{\pgfqpoint{0.000000in}{0.000000in}}%
\pgfpathlineto{\pgfqpoint{-0.048611in}{0.000000in}}%
\pgfusepath{stroke,fill}%
}%
\begin{pgfscope}%
\pgfsys@transformshift{0.813028in}{0.947224in}%
\pgfsys@useobject{currentmarker}{}%
\end{pgfscope}%
\end{pgfscope}%
\begin{pgfscope}%
\definecolor{textcolor}{rgb}{0.000000,0.000000,0.000000}%
\pgfsetstrokecolor{textcolor}%
\pgfsetfillcolor{textcolor}%
\pgftext[x=0.378552in,y=0.894463in,left,base]{\color{textcolor}\sffamily\fontsize{10.000000}{12.000000}\selectfont −0.6}%
\end{pgfscope}%
\begin{pgfscope}%
\pgfsetbuttcap%
\pgfsetroundjoin%
\definecolor{currentfill}{rgb}{0.000000,0.000000,0.000000}%
\pgfsetfillcolor{currentfill}%
\pgfsetlinewidth{0.803000pt}%
\definecolor{currentstroke}{rgb}{0.000000,0.000000,0.000000}%
\pgfsetstrokecolor{currentstroke}%
\pgfsetdash{}{0pt}%
\pgfsys@defobject{currentmarker}{\pgfqpoint{-0.048611in}{0.000000in}}{\pgfqpoint{0.000000in}{0.000000in}}{%
\pgfpathmoveto{\pgfqpoint{0.000000in}{0.000000in}}%
\pgfpathlineto{\pgfqpoint{-0.048611in}{0.000000in}}%
\pgfusepath{stroke,fill}%
}%
\begin{pgfscope}%
\pgfsys@transformshift{0.813028in}{1.422978in}%
\pgfsys@useobject{currentmarker}{}%
\end{pgfscope}%
\end{pgfscope}%
\begin{pgfscope}%
\definecolor{textcolor}{rgb}{0.000000,0.000000,0.000000}%
\pgfsetstrokecolor{textcolor}%
\pgfsetfillcolor{textcolor}%
\pgftext[x=0.378552in,y=1.370216in,left,base]{\color{textcolor}\sffamily\fontsize{10.000000}{12.000000}\selectfont −0.5}%
\end{pgfscope}%
\begin{pgfscope}%
\pgfsetbuttcap%
\pgfsetroundjoin%
\definecolor{currentfill}{rgb}{0.000000,0.000000,0.000000}%
\pgfsetfillcolor{currentfill}%
\pgfsetlinewidth{0.803000pt}%
\definecolor{currentstroke}{rgb}{0.000000,0.000000,0.000000}%
\pgfsetstrokecolor{currentstroke}%
\pgfsetdash{}{0pt}%
\pgfsys@defobject{currentmarker}{\pgfqpoint{-0.048611in}{0.000000in}}{\pgfqpoint{0.000000in}{0.000000in}}{%
\pgfpathmoveto{\pgfqpoint{0.000000in}{0.000000in}}%
\pgfpathlineto{\pgfqpoint{-0.048611in}{0.000000in}}%
\pgfusepath{stroke,fill}%
}%
\begin{pgfscope}%
\pgfsys@transformshift{0.813028in}{1.898731in}%
\pgfsys@useobject{currentmarker}{}%
\end{pgfscope}%
\end{pgfscope}%
\begin{pgfscope}%
\definecolor{textcolor}{rgb}{0.000000,0.000000,0.000000}%
\pgfsetstrokecolor{textcolor}%
\pgfsetfillcolor{textcolor}%
\pgftext[x=0.378552in,y=1.845969in,left,base]{\color{textcolor}\sffamily\fontsize{10.000000}{12.000000}\selectfont −0.4}%
\end{pgfscope}%
\begin{pgfscope}%
\pgfsetbuttcap%
\pgfsetroundjoin%
\definecolor{currentfill}{rgb}{0.000000,0.000000,0.000000}%
\pgfsetfillcolor{currentfill}%
\pgfsetlinewidth{0.803000pt}%
\definecolor{currentstroke}{rgb}{0.000000,0.000000,0.000000}%
\pgfsetstrokecolor{currentstroke}%
\pgfsetdash{}{0pt}%
\pgfsys@defobject{currentmarker}{\pgfqpoint{-0.048611in}{0.000000in}}{\pgfqpoint{0.000000in}{0.000000in}}{%
\pgfpathmoveto{\pgfqpoint{0.000000in}{0.000000in}}%
\pgfpathlineto{\pgfqpoint{-0.048611in}{0.000000in}}%
\pgfusepath{stroke,fill}%
}%
\begin{pgfscope}%
\pgfsys@transformshift{0.813028in}{2.374484in}%
\pgfsys@useobject{currentmarker}{}%
\end{pgfscope}%
\end{pgfscope}%
\begin{pgfscope}%
\definecolor{textcolor}{rgb}{0.000000,0.000000,0.000000}%
\pgfsetstrokecolor{textcolor}%
\pgfsetfillcolor{textcolor}%
\pgftext[x=0.378552in,y=2.321723in,left,base]{\color{textcolor}\sffamily\fontsize{10.000000}{12.000000}\selectfont −0.3}%
\end{pgfscope}%
\begin{pgfscope}%
\pgfsetbuttcap%
\pgfsetroundjoin%
\definecolor{currentfill}{rgb}{0.000000,0.000000,0.000000}%
\pgfsetfillcolor{currentfill}%
\pgfsetlinewidth{0.803000pt}%
\definecolor{currentstroke}{rgb}{0.000000,0.000000,0.000000}%
\pgfsetstrokecolor{currentstroke}%
\pgfsetdash{}{0pt}%
\pgfsys@defobject{currentmarker}{\pgfqpoint{-0.048611in}{0.000000in}}{\pgfqpoint{0.000000in}{0.000000in}}{%
\pgfpathmoveto{\pgfqpoint{0.000000in}{0.000000in}}%
\pgfpathlineto{\pgfqpoint{-0.048611in}{0.000000in}}%
\pgfusepath{stroke,fill}%
}%
\begin{pgfscope}%
\pgfsys@transformshift{0.813028in}{2.850237in}%
\pgfsys@useobject{currentmarker}{}%
\end{pgfscope}%
\end{pgfscope}%
\begin{pgfscope}%
\definecolor{textcolor}{rgb}{0.000000,0.000000,0.000000}%
\pgfsetstrokecolor{textcolor}%
\pgfsetfillcolor{textcolor}%
\pgftext[x=0.378552in,y=2.797476in,left,base]{\color{textcolor}\sffamily\fontsize{10.000000}{12.000000}\selectfont −0.2}%
\end{pgfscope}%
\begin{pgfscope}%
\pgfsetbuttcap%
\pgfsetroundjoin%
\definecolor{currentfill}{rgb}{0.000000,0.000000,0.000000}%
\pgfsetfillcolor{currentfill}%
\pgfsetlinewidth{0.803000pt}%
\definecolor{currentstroke}{rgb}{0.000000,0.000000,0.000000}%
\pgfsetstrokecolor{currentstroke}%
\pgfsetdash{}{0pt}%
\pgfsys@defobject{currentmarker}{\pgfqpoint{-0.048611in}{0.000000in}}{\pgfqpoint{0.000000in}{0.000000in}}{%
\pgfpathmoveto{\pgfqpoint{0.000000in}{0.000000in}}%
\pgfpathlineto{\pgfqpoint{-0.048611in}{0.000000in}}%
\pgfusepath{stroke,fill}%
}%
\begin{pgfscope}%
\pgfsys@transformshift{0.813028in}{3.325991in}%
\pgfsys@useobject{currentmarker}{}%
\end{pgfscope}%
\end{pgfscope}%
\begin{pgfscope}%
\definecolor{textcolor}{rgb}{0.000000,0.000000,0.000000}%
\pgfsetstrokecolor{textcolor}%
\pgfsetfillcolor{textcolor}%
\pgftext[x=0.378552in,y=3.273229in,left,base]{\color{textcolor}\sffamily\fontsize{10.000000}{12.000000}\selectfont −0.1}%
\end{pgfscope}%
\begin{pgfscope}%
\pgfsetbuttcap%
\pgfsetroundjoin%
\definecolor{currentfill}{rgb}{0.000000,0.000000,0.000000}%
\pgfsetfillcolor{currentfill}%
\pgfsetlinewidth{0.803000pt}%
\definecolor{currentstroke}{rgb}{0.000000,0.000000,0.000000}%
\pgfsetstrokecolor{currentstroke}%
\pgfsetdash{}{0pt}%
\pgfsys@defobject{currentmarker}{\pgfqpoint{-0.048611in}{0.000000in}}{\pgfqpoint{0.000000in}{0.000000in}}{%
\pgfpathmoveto{\pgfqpoint{0.000000in}{0.000000in}}%
\pgfpathlineto{\pgfqpoint{-0.048611in}{0.000000in}}%
\pgfusepath{stroke,fill}%
}%
\begin{pgfscope}%
\pgfsys@transformshift{0.813028in}{3.801744in}%
\pgfsys@useobject{currentmarker}{}%
\end{pgfscope}%
\end{pgfscope}%
\begin{pgfscope}%
\definecolor{textcolor}{rgb}{0.000000,0.000000,0.000000}%
\pgfsetstrokecolor{textcolor}%
\pgfsetfillcolor{textcolor}%
\pgftext[x=0.494926in,y=3.748982in,left,base]{\color{textcolor}\sffamily\fontsize{10.000000}{12.000000}\selectfont 0.0}%
\end{pgfscope}%
\begin{pgfscope}%
\pgfsetbuttcap%
\pgfsetroundjoin%
\definecolor{currentfill}{rgb}{0.000000,0.000000,0.000000}%
\pgfsetfillcolor{currentfill}%
\pgfsetlinewidth{0.803000pt}%
\definecolor{currentstroke}{rgb}{0.000000,0.000000,0.000000}%
\pgfsetstrokecolor{currentstroke}%
\pgfsetdash{}{0pt}%
\pgfsys@defobject{currentmarker}{\pgfqpoint{-0.048611in}{0.000000in}}{\pgfqpoint{0.000000in}{0.000000in}}{%
\pgfpathmoveto{\pgfqpoint{0.000000in}{0.000000in}}%
\pgfpathlineto{\pgfqpoint{-0.048611in}{0.000000in}}%
\pgfusepath{stroke,fill}%
}%
\begin{pgfscope}%
\pgfsys@transformshift{0.813028in}{4.277497in}%
\pgfsys@useobject{currentmarker}{}%
\end{pgfscope}%
\end{pgfscope}%
\begin{pgfscope}%
\definecolor{textcolor}{rgb}{0.000000,0.000000,0.000000}%
\pgfsetstrokecolor{textcolor}%
\pgfsetfillcolor{textcolor}%
\pgftext[x=0.494926in,y=4.224736in,left,base]{\color{textcolor}\sffamily\fontsize{10.000000}{12.000000}\selectfont 0.1}%
\end{pgfscope}%
\begin{pgfscope}%
\definecolor{textcolor}{rgb}{0.000000,0.000000,0.000000}%
\pgfsetstrokecolor{textcolor}%
\pgfsetfillcolor{textcolor}%
\pgftext[x=0.322997in,y=2.659389in,,bottom,rotate=90.000000]{\color{textcolor}\sffamily\fontsize{10.000000}{12.000000}\selectfont Skewness \(\displaystyle S_u\)}%
\end{pgfscope}%
\begin{pgfscope}%
\pgfpathrectangle{\pgfqpoint{0.813028in}{0.668778in}}{\pgfqpoint{5.436972in}{3.981222in}}%
\pgfusepath{clip}%
\pgfsetbuttcap%
\pgfsetroundjoin%
\pgfsetlinewidth{1.505625pt}%
\definecolor{currentstroke}{rgb}{0.121569,0.466667,0.705882}%
\pgfsetstrokecolor{currentstroke}%
\pgfsetdash{{1.500000pt}{2.475000pt}}{0.000000pt}%
\pgfpathmoveto{\pgfqpoint{1.060163in}{4.469035in}}%
\pgfpathlineto{\pgfqpoint{1.190148in}{4.359850in}}%
\pgfpathlineto{\pgfqpoint{1.451464in}{4.182832in}}%
\pgfpathlineto{\pgfqpoint{1.654156in}{3.864415in}}%
\pgfpathlineto{\pgfqpoint{2.913394in}{2.591903in}}%
\pgfpathlineto{\pgfqpoint{3.543013in}{2.625444in}}%
\pgfpathlineto{\pgfqpoint{4.114008in}{3.047199in}}%
\pgfpathlineto{\pgfqpoint{4.461898in}{3.036114in}}%
\pgfpathlineto{\pgfqpoint{4.712833in}{2.921553in}}%
\pgfpathlineto{\pgfqpoint{4.909238in}{2.775401in}}%
\pgfpathlineto{\pgfqpoint{5.004942in}{2.619497in}}%
\pgfpathlineto{\pgfqpoint{5.373246in}{1.926610in}}%
\pgfpathlineto{\pgfqpoint{5.634561in}{1.152036in}}%
\pgfpathlineto{\pgfqpoint{5.837253in}{0.849742in}}%
\pgfpathlineto{\pgfqpoint{6.002865in}{1.542820in}}%
\pgfusepath{stroke}%
\end{pgfscope}%
\begin{pgfscope}%
\pgfpathrectangle{\pgfqpoint{0.813028in}{0.668778in}}{\pgfqpoint{5.436972in}{3.981222in}}%
\pgfusepath{clip}%
\pgfsetbuttcap%
\pgfsetroundjoin%
\definecolor{currentfill}{rgb}{0.121569,0.466667,0.705882}%
\pgfsetfillcolor{currentfill}%
\pgfsetlinewidth{1.003750pt}%
\definecolor{currentstroke}{rgb}{0.121569,0.466667,0.705882}%
\pgfsetstrokecolor{currentstroke}%
\pgfsetdash{}{0pt}%
\pgfsys@defobject{currentmarker}{\pgfqpoint{-0.034722in}{-0.034722in}}{\pgfqpoint{0.034722in}{0.034722in}}{%
\pgfpathmoveto{\pgfqpoint{0.000000in}{-0.034722in}}%
\pgfpathcurveto{\pgfqpoint{0.009208in}{-0.034722in}}{\pgfqpoint{0.018041in}{-0.031064in}}{\pgfqpoint{0.024552in}{-0.024552in}}%
\pgfpathcurveto{\pgfqpoint{0.031064in}{-0.018041in}}{\pgfqpoint{0.034722in}{-0.009208in}}{\pgfqpoint{0.034722in}{0.000000in}}%
\pgfpathcurveto{\pgfqpoint{0.034722in}{0.009208in}}{\pgfqpoint{0.031064in}{0.018041in}}{\pgfqpoint{0.024552in}{0.024552in}}%
\pgfpathcurveto{\pgfqpoint{0.018041in}{0.031064in}}{\pgfqpoint{0.009208in}{0.034722in}}{\pgfqpoint{0.000000in}{0.034722in}}%
\pgfpathcurveto{\pgfqpoint{-0.009208in}{0.034722in}}{\pgfqpoint{-0.018041in}{0.031064in}}{\pgfqpoint{-0.024552in}{0.024552in}}%
\pgfpathcurveto{\pgfqpoint{-0.031064in}{0.018041in}}{\pgfqpoint{-0.034722in}{0.009208in}}{\pgfqpoint{-0.034722in}{0.000000in}}%
\pgfpathcurveto{\pgfqpoint{-0.034722in}{-0.009208in}}{\pgfqpoint{-0.031064in}{-0.018041in}}{\pgfqpoint{-0.024552in}{-0.024552in}}%
\pgfpathcurveto{\pgfqpoint{-0.018041in}{-0.031064in}}{\pgfqpoint{-0.009208in}{-0.034722in}}{\pgfqpoint{0.000000in}{-0.034722in}}%
\pgfpathclose%
\pgfusepath{stroke,fill}%
}%
\begin{pgfscope}%
\pgfsys@transformshift{1.060163in}{4.469035in}%
\pgfsys@useobject{currentmarker}{}%
\end{pgfscope}%
\begin{pgfscope}%
\pgfsys@transformshift{1.190148in}{4.359850in}%
\pgfsys@useobject{currentmarker}{}%
\end{pgfscope}%
\begin{pgfscope}%
\pgfsys@transformshift{1.451464in}{4.182832in}%
\pgfsys@useobject{currentmarker}{}%
\end{pgfscope}%
\begin{pgfscope}%
\pgfsys@transformshift{1.654156in}{3.864415in}%
\pgfsys@useobject{currentmarker}{}%
\end{pgfscope}%
\begin{pgfscope}%
\pgfsys@transformshift{2.913394in}{2.591903in}%
\pgfsys@useobject{currentmarker}{}%
\end{pgfscope}%
\begin{pgfscope}%
\pgfsys@transformshift{3.543013in}{2.625444in}%
\pgfsys@useobject{currentmarker}{}%
\end{pgfscope}%
\begin{pgfscope}%
\pgfsys@transformshift{4.114008in}{3.047199in}%
\pgfsys@useobject{currentmarker}{}%
\end{pgfscope}%
\begin{pgfscope}%
\pgfsys@transformshift{4.461898in}{3.036114in}%
\pgfsys@useobject{currentmarker}{}%
\end{pgfscope}%
\begin{pgfscope}%
\pgfsys@transformshift{4.712833in}{2.921553in}%
\pgfsys@useobject{currentmarker}{}%
\end{pgfscope}%
\begin{pgfscope}%
\pgfsys@transformshift{4.909238in}{2.775401in}%
\pgfsys@useobject{currentmarker}{}%
\end{pgfscope}%
\begin{pgfscope}%
\pgfsys@transformshift{5.004942in}{2.619497in}%
\pgfsys@useobject{currentmarker}{}%
\end{pgfscope}%
\begin{pgfscope}%
\pgfsys@transformshift{5.373246in}{1.926610in}%
\pgfsys@useobject{currentmarker}{}%
\end{pgfscope}%
\begin{pgfscope}%
\pgfsys@transformshift{5.634561in}{1.152036in}%
\pgfsys@useobject{currentmarker}{}%
\end{pgfscope}%
\begin{pgfscope}%
\pgfsys@transformshift{5.837253in}{0.849742in}%
\pgfsys@useobject{currentmarker}{}%
\end{pgfscope}%
\begin{pgfscope}%
\pgfsys@transformshift{6.002865in}{1.542820in}%
\pgfsys@useobject{currentmarker}{}%
\end{pgfscope}%
\end{pgfscope}%
\begin{pgfscope}%
\pgfsetrectcap%
\pgfsetmiterjoin%
\pgfsetlinewidth{0.803000pt}%
\definecolor{currentstroke}{rgb}{0.000000,0.000000,0.000000}%
\pgfsetstrokecolor{currentstroke}%
\pgfsetdash{}{0pt}%
\pgfpathmoveto{\pgfqpoint{0.813028in}{0.668778in}}%
\pgfpathlineto{\pgfqpoint{0.813028in}{4.650000in}}%
\pgfusepath{stroke}%
\end{pgfscope}%
\begin{pgfscope}%
\pgfsetrectcap%
\pgfsetmiterjoin%
\pgfsetlinewidth{0.803000pt}%
\definecolor{currentstroke}{rgb}{0.000000,0.000000,0.000000}%
\pgfsetstrokecolor{currentstroke}%
\pgfsetdash{}{0pt}%
\pgfpathmoveto{\pgfqpoint{6.250000in}{0.668778in}}%
\pgfpathlineto{\pgfqpoint{6.250000in}{4.650000in}}%
\pgfusepath{stroke}%
\end{pgfscope}%
\begin{pgfscope}%
\pgfsetrectcap%
\pgfsetmiterjoin%
\pgfsetlinewidth{0.803000pt}%
\definecolor{currentstroke}{rgb}{0.000000,0.000000,0.000000}%
\pgfsetstrokecolor{currentstroke}%
\pgfsetdash{}{0pt}%
\pgfpathmoveto{\pgfqpoint{0.813028in}{0.668778in}}%
\pgfpathlineto{\pgfqpoint{6.250000in}{0.668778in}}%
\pgfusepath{stroke}%
\end{pgfscope}%
\begin{pgfscope}%
\pgfsetrectcap%
\pgfsetmiterjoin%
\pgfsetlinewidth{0.803000pt}%
\definecolor{currentstroke}{rgb}{0.000000,0.000000,0.000000}%
\pgfsetstrokecolor{currentstroke}%
\pgfsetdash{}{0pt}%
\pgfpathmoveto{\pgfqpoint{0.813028in}{4.650000in}}%
\pgfpathlineto{\pgfqpoint{6.250000in}{4.650000in}}%
\pgfusepath{stroke}%
\end{pgfscope}%
\end{pgfpicture}%
\makeatother%
\endgroup%
}
    \caption{Skewness profile $S_u$ displayed in log-lin scale}
    \label{fig:Skewness_Profile}
\end{figure}

\begin{figure}
    \centering
    \begin{tabular}{ll}
        a. & b. \\
         \resizebox{0.47\linewidth}{!}{%% Creator: Matplotlib, PGF backend
%%
%% To include the figure in your LaTeX document, write
%%   \input{<filename>.pgf}
%%
%% Make sure the required packages are loaded in your preamble
%%   \usepackage{pgf}
%%
%% Figures using additional raster images can only be included by \input if
%% they are in the same directory as the main LaTeX file. For loading figures
%% from other directories you can use the `import` package
%%   \usepackage{import}
%% and then include the figures with
%%   \import{<path to file>}{<filename>.pgf}
%%
%% Matplotlib used the following preamble
%%   \usepackage{fontspec}
%%   \setmainfont{DejaVuSerif.ttf}[Path=C:/Users/garaya/AppData/Local/Continuum/anaconda3/lib/site-packages/matplotlib/mpl-data/fonts/ttf/]
%%   \setsansfont{DejaVuSans.ttf}[Path=C:/Users/garaya/AppData/Local/Continuum/anaconda3/lib/site-packages/matplotlib/mpl-data/fonts/ttf/]
%%   \setmonofont{DejaVuSansMono.ttf}[Path=C:/Users/garaya/AppData/Local/Continuum/anaconda3/lib/site-packages/matplotlib/mpl-data/fonts/ttf/]
%%
\begingroup%
\makeatletter%
\begin{pgfpicture}%
\pgfpathrectangle{\pgfpointorigin}{\pgfqpoint{6.400000in}{4.800000in}}%
\pgfusepath{use as bounding box, clip}%
\begin{pgfscope}%
\pgfsetbuttcap%
\pgfsetmiterjoin%
\definecolor{currentfill}{rgb}{1.000000,1.000000,1.000000}%
\pgfsetfillcolor{currentfill}%
\pgfsetlinewidth{0.000000pt}%
\definecolor{currentstroke}{rgb}{1.000000,1.000000,1.000000}%
\pgfsetstrokecolor{currentstroke}%
\pgfsetdash{}{0pt}%
\pgfpathmoveto{\pgfqpoint{0.000000in}{0.000000in}}%
\pgfpathlineto{\pgfqpoint{6.400000in}{0.000000in}}%
\pgfpathlineto{\pgfqpoint{6.400000in}{4.800000in}}%
\pgfpathlineto{\pgfqpoint{0.000000in}{4.800000in}}%
\pgfpathclose%
\pgfusepath{fill}%
\end{pgfscope}%
\begin{pgfscope}%
\pgfsetbuttcap%
\pgfsetmiterjoin%
\definecolor{currentfill}{rgb}{1.000000,1.000000,1.000000}%
\pgfsetfillcolor{currentfill}%
\pgfsetlinewidth{0.000000pt}%
\definecolor{currentstroke}{rgb}{0.000000,0.000000,0.000000}%
\pgfsetstrokecolor{currentstroke}%
\pgfsetstrokeopacity{0.000000}%
\pgfsetdash{}{0pt}%
\pgfpathmoveto{\pgfqpoint{0.800000in}{0.528000in}}%
\pgfpathlineto{\pgfqpoint{5.760000in}{0.528000in}}%
\pgfpathlineto{\pgfqpoint{5.760000in}{4.224000in}}%
\pgfpathlineto{\pgfqpoint{0.800000in}{4.224000in}}%
\pgfpathclose%
\pgfusepath{fill}%
\end{pgfscope}%
\begin{pgfscope}%
\pgfsetbuttcap%
\pgfsetroundjoin%
\definecolor{currentfill}{rgb}{0.000000,0.000000,0.000000}%
\pgfsetfillcolor{currentfill}%
\pgfsetlinewidth{0.803000pt}%
\definecolor{currentstroke}{rgb}{0.000000,0.000000,0.000000}%
\pgfsetstrokecolor{currentstroke}%
\pgfsetdash{}{0pt}%
\pgfsys@defobject{currentmarker}{\pgfqpoint{0.000000in}{-0.048611in}}{\pgfqpoint{0.000000in}{0.000000in}}{%
\pgfpathmoveto{\pgfqpoint{0.000000in}{0.000000in}}%
\pgfpathlineto{\pgfqpoint{0.000000in}{-0.048611in}}%
\pgfusepath{stroke,fill}%
}%
\begin{pgfscope}%
\pgfsys@transformshift{1.365491in}{0.528000in}%
\pgfsys@useobject{currentmarker}{}%
\end{pgfscope}%
\end{pgfscope}%
\begin{pgfscope}%
\definecolor{textcolor}{rgb}{0.000000,0.000000,0.000000}%
\pgfsetstrokecolor{textcolor}%
\pgfsetfillcolor{textcolor}%
\pgftext[x=1.365491in,y=0.430778in,,top]{\color{textcolor}\rmfamily\fontsize{14.000000}{16.800000}\selectfont \(\displaystyle 1\)}%
\end{pgfscope}%
\begin{pgfscope}%
\pgfsetbuttcap%
\pgfsetroundjoin%
\definecolor{currentfill}{rgb}{0.000000,0.000000,0.000000}%
\pgfsetfillcolor{currentfill}%
\pgfsetlinewidth{0.803000pt}%
\definecolor{currentstroke}{rgb}{0.000000,0.000000,0.000000}%
\pgfsetstrokecolor{currentstroke}%
\pgfsetdash{}{0pt}%
\pgfsys@defobject{currentmarker}{\pgfqpoint{0.000000in}{-0.048611in}}{\pgfqpoint{0.000000in}{0.000000in}}{%
\pgfpathmoveto{\pgfqpoint{0.000000in}{0.000000in}}%
\pgfpathlineto{\pgfqpoint{0.000000in}{-0.048611in}}%
\pgfusepath{stroke,fill}%
}%
\begin{pgfscope}%
\pgfsys@transformshift{2.485266in}{0.528000in}%
\pgfsys@useobject{currentmarker}{}%
\end{pgfscope}%
\end{pgfscope}%
\begin{pgfscope}%
\definecolor{textcolor}{rgb}{0.000000,0.000000,0.000000}%
\pgfsetstrokecolor{textcolor}%
\pgfsetfillcolor{textcolor}%
\pgftext[x=2.485266in,y=0.430778in,,top]{\color{textcolor}\rmfamily\fontsize{14.000000}{16.800000}\selectfont \(\displaystyle 2\)}%
\end{pgfscope}%
\begin{pgfscope}%
\pgfsetbuttcap%
\pgfsetroundjoin%
\definecolor{currentfill}{rgb}{0.000000,0.000000,0.000000}%
\pgfsetfillcolor{currentfill}%
\pgfsetlinewidth{0.803000pt}%
\definecolor{currentstroke}{rgb}{0.000000,0.000000,0.000000}%
\pgfsetstrokecolor{currentstroke}%
\pgfsetdash{}{0pt}%
\pgfsys@defobject{currentmarker}{\pgfqpoint{0.000000in}{-0.048611in}}{\pgfqpoint{0.000000in}{0.000000in}}{%
\pgfpathmoveto{\pgfqpoint{0.000000in}{0.000000in}}%
\pgfpathlineto{\pgfqpoint{0.000000in}{-0.048611in}}%
\pgfusepath{stroke,fill}%
}%
\begin{pgfscope}%
\pgfsys@transformshift{3.605042in}{0.528000in}%
\pgfsys@useobject{currentmarker}{}%
\end{pgfscope}%
\end{pgfscope}%
\begin{pgfscope}%
\definecolor{textcolor}{rgb}{0.000000,0.000000,0.000000}%
\pgfsetstrokecolor{textcolor}%
\pgfsetfillcolor{textcolor}%
\pgftext[x=3.605042in,y=0.430778in,,top]{\color{textcolor}\rmfamily\fontsize{14.000000}{16.800000}\selectfont \(\displaystyle 3\)}%
\end{pgfscope}%
\begin{pgfscope}%
\pgfsetbuttcap%
\pgfsetroundjoin%
\definecolor{currentfill}{rgb}{0.000000,0.000000,0.000000}%
\pgfsetfillcolor{currentfill}%
\pgfsetlinewidth{0.803000pt}%
\definecolor{currentstroke}{rgb}{0.000000,0.000000,0.000000}%
\pgfsetstrokecolor{currentstroke}%
\pgfsetdash{}{0pt}%
\pgfsys@defobject{currentmarker}{\pgfqpoint{0.000000in}{-0.048611in}}{\pgfqpoint{0.000000in}{0.000000in}}{%
\pgfpathmoveto{\pgfqpoint{0.000000in}{0.000000in}}%
\pgfpathlineto{\pgfqpoint{0.000000in}{-0.048611in}}%
\pgfusepath{stroke,fill}%
}%
\begin{pgfscope}%
\pgfsys@transformshift{4.724817in}{0.528000in}%
\pgfsys@useobject{currentmarker}{}%
\end{pgfscope}%
\end{pgfscope}%
\begin{pgfscope}%
\definecolor{textcolor}{rgb}{0.000000,0.000000,0.000000}%
\pgfsetstrokecolor{textcolor}%
\pgfsetfillcolor{textcolor}%
\pgftext[x=4.724817in,y=0.430778in,,top]{\color{textcolor}\rmfamily\fontsize{14.000000}{16.800000}\selectfont \(\displaystyle 4\)}%
\end{pgfscope}%
\begin{pgfscope}%
\definecolor{textcolor}{rgb}{0.000000,0.000000,0.000000}%
\pgfsetstrokecolor{textcolor}%
\pgfsetfillcolor{textcolor}%
\pgftext[x=3.280000in,y=0.187044in,,top]{\color{textcolor}\rmfamily\fontsize{14.000000}{16.800000}\selectfont Speed range \(\displaystyle u\) (m/s)}%
\end{pgfscope}%
\begin{pgfscope}%
\pgfsetbuttcap%
\pgfsetroundjoin%
\definecolor{currentfill}{rgb}{0.000000,0.000000,0.000000}%
\pgfsetfillcolor{currentfill}%
\pgfsetlinewidth{0.803000pt}%
\definecolor{currentstroke}{rgb}{0.000000,0.000000,0.000000}%
\pgfsetstrokecolor{currentstroke}%
\pgfsetdash{}{0pt}%
\pgfsys@defobject{currentmarker}{\pgfqpoint{-0.048611in}{0.000000in}}{\pgfqpoint{0.000000in}{0.000000in}}{%
\pgfpathmoveto{\pgfqpoint{0.000000in}{0.000000in}}%
\pgfpathlineto{\pgfqpoint{-0.048611in}{0.000000in}}%
\pgfusepath{stroke,fill}%
}%
\begin{pgfscope}%
\pgfsys@transformshift{0.800000in}{0.696000in}%
\pgfsys@useobject{currentmarker}{}%
\end{pgfscope}%
\end{pgfscope}%
\begin{pgfscope}%
\definecolor{textcolor}{rgb}{0.000000,0.000000,0.000000}%
\pgfsetstrokecolor{textcolor}%
\pgfsetfillcolor{textcolor}%
\pgftext[x=0.452550in,y=0.622134in,left,base]{\color{textcolor}\rmfamily\fontsize{14.000000}{16.800000}\selectfont \(\displaystyle 0.0\)}%
\end{pgfscope}%
\begin{pgfscope}%
\pgfsetbuttcap%
\pgfsetroundjoin%
\definecolor{currentfill}{rgb}{0.000000,0.000000,0.000000}%
\pgfsetfillcolor{currentfill}%
\pgfsetlinewidth{0.803000pt}%
\definecolor{currentstroke}{rgb}{0.000000,0.000000,0.000000}%
\pgfsetstrokecolor{currentstroke}%
\pgfsetdash{}{0pt}%
\pgfsys@defobject{currentmarker}{\pgfqpoint{-0.048611in}{0.000000in}}{\pgfqpoint{0.000000in}{0.000000in}}{%
\pgfpathmoveto{\pgfqpoint{0.000000in}{0.000000in}}%
\pgfpathlineto{\pgfqpoint{-0.048611in}{0.000000in}}%
\pgfusepath{stroke,fill}%
}%
\begin{pgfscope}%
\pgfsys@transformshift{0.800000in}{1.302741in}%
\pgfsys@useobject{currentmarker}{}%
\end{pgfscope}%
\end{pgfscope}%
\begin{pgfscope}%
\definecolor{textcolor}{rgb}{0.000000,0.000000,0.000000}%
\pgfsetstrokecolor{textcolor}%
\pgfsetfillcolor{textcolor}%
\pgftext[x=0.452550in,y=1.228875in,left,base]{\color{textcolor}\rmfamily\fontsize{14.000000}{16.800000}\selectfont \(\displaystyle 0.1\)}%
\end{pgfscope}%
\begin{pgfscope}%
\pgfsetbuttcap%
\pgfsetroundjoin%
\definecolor{currentfill}{rgb}{0.000000,0.000000,0.000000}%
\pgfsetfillcolor{currentfill}%
\pgfsetlinewidth{0.803000pt}%
\definecolor{currentstroke}{rgb}{0.000000,0.000000,0.000000}%
\pgfsetstrokecolor{currentstroke}%
\pgfsetdash{}{0pt}%
\pgfsys@defobject{currentmarker}{\pgfqpoint{-0.048611in}{0.000000in}}{\pgfqpoint{0.000000in}{0.000000in}}{%
\pgfpathmoveto{\pgfqpoint{0.000000in}{0.000000in}}%
\pgfpathlineto{\pgfqpoint{-0.048611in}{0.000000in}}%
\pgfusepath{stroke,fill}%
}%
\begin{pgfscope}%
\pgfsys@transformshift{0.800000in}{1.909483in}%
\pgfsys@useobject{currentmarker}{}%
\end{pgfscope}%
\end{pgfscope}%
\begin{pgfscope}%
\definecolor{textcolor}{rgb}{0.000000,0.000000,0.000000}%
\pgfsetstrokecolor{textcolor}%
\pgfsetfillcolor{textcolor}%
\pgftext[x=0.452550in,y=1.835616in,left,base]{\color{textcolor}\rmfamily\fontsize{14.000000}{16.800000}\selectfont \(\displaystyle 0.2\)}%
\end{pgfscope}%
\begin{pgfscope}%
\pgfsetbuttcap%
\pgfsetroundjoin%
\definecolor{currentfill}{rgb}{0.000000,0.000000,0.000000}%
\pgfsetfillcolor{currentfill}%
\pgfsetlinewidth{0.803000pt}%
\definecolor{currentstroke}{rgb}{0.000000,0.000000,0.000000}%
\pgfsetstrokecolor{currentstroke}%
\pgfsetdash{}{0pt}%
\pgfsys@defobject{currentmarker}{\pgfqpoint{-0.048611in}{0.000000in}}{\pgfqpoint{0.000000in}{0.000000in}}{%
\pgfpathmoveto{\pgfqpoint{0.000000in}{0.000000in}}%
\pgfpathlineto{\pgfqpoint{-0.048611in}{0.000000in}}%
\pgfusepath{stroke,fill}%
}%
\begin{pgfscope}%
\pgfsys@transformshift{0.800000in}{2.516224in}%
\pgfsys@useobject{currentmarker}{}%
\end{pgfscope}%
\end{pgfscope}%
\begin{pgfscope}%
\definecolor{textcolor}{rgb}{0.000000,0.000000,0.000000}%
\pgfsetstrokecolor{textcolor}%
\pgfsetfillcolor{textcolor}%
\pgftext[x=0.452550in,y=2.442358in,left,base]{\color{textcolor}\rmfamily\fontsize{14.000000}{16.800000}\selectfont \(\displaystyle 0.3\)}%
\end{pgfscope}%
\begin{pgfscope}%
\pgfsetbuttcap%
\pgfsetroundjoin%
\definecolor{currentfill}{rgb}{0.000000,0.000000,0.000000}%
\pgfsetfillcolor{currentfill}%
\pgfsetlinewidth{0.803000pt}%
\definecolor{currentstroke}{rgb}{0.000000,0.000000,0.000000}%
\pgfsetstrokecolor{currentstroke}%
\pgfsetdash{}{0pt}%
\pgfsys@defobject{currentmarker}{\pgfqpoint{-0.048611in}{0.000000in}}{\pgfqpoint{0.000000in}{0.000000in}}{%
\pgfpathmoveto{\pgfqpoint{0.000000in}{0.000000in}}%
\pgfpathlineto{\pgfqpoint{-0.048611in}{0.000000in}}%
\pgfusepath{stroke,fill}%
}%
\begin{pgfscope}%
\pgfsys@transformshift{0.800000in}{3.122965in}%
\pgfsys@useobject{currentmarker}{}%
\end{pgfscope}%
\end{pgfscope}%
\begin{pgfscope}%
\definecolor{textcolor}{rgb}{0.000000,0.000000,0.000000}%
\pgfsetstrokecolor{textcolor}%
\pgfsetfillcolor{textcolor}%
\pgftext[x=0.452550in,y=3.049099in,left,base]{\color{textcolor}\rmfamily\fontsize{14.000000}{16.800000}\selectfont \(\displaystyle 0.4\)}%
\end{pgfscope}%
\begin{pgfscope}%
\pgfsetbuttcap%
\pgfsetroundjoin%
\definecolor{currentfill}{rgb}{0.000000,0.000000,0.000000}%
\pgfsetfillcolor{currentfill}%
\pgfsetlinewidth{0.803000pt}%
\definecolor{currentstroke}{rgb}{0.000000,0.000000,0.000000}%
\pgfsetstrokecolor{currentstroke}%
\pgfsetdash{}{0pt}%
\pgfsys@defobject{currentmarker}{\pgfqpoint{-0.048611in}{0.000000in}}{\pgfqpoint{0.000000in}{0.000000in}}{%
\pgfpathmoveto{\pgfqpoint{0.000000in}{0.000000in}}%
\pgfpathlineto{\pgfqpoint{-0.048611in}{0.000000in}}%
\pgfusepath{stroke,fill}%
}%
\begin{pgfscope}%
\pgfsys@transformshift{0.800000in}{3.729706in}%
\pgfsys@useobject{currentmarker}{}%
\end{pgfscope}%
\end{pgfscope}%
\begin{pgfscope}%
\definecolor{textcolor}{rgb}{0.000000,0.000000,0.000000}%
\pgfsetstrokecolor{textcolor}%
\pgfsetfillcolor{textcolor}%
\pgftext[x=0.452550in,y=3.655840in,left,base]{\color{textcolor}\rmfamily\fontsize{14.000000}{16.800000}\selectfont \(\displaystyle 0.5\)}%
\end{pgfscope}%
\begin{pgfscope}%
\definecolor{textcolor}{rgb}{0.000000,0.000000,0.000000}%
\pgfsetstrokecolor{textcolor}%
\pgfsetfillcolor{textcolor}%
\pgftext[x=0.396994in,y=2.376000in,,bottom,rotate=90.000000]{\color{textcolor}\rmfamily\fontsize{14.000000}{16.800000}\selectfont PDF (s/m)}%
\end{pgfscope}%
\begin{pgfscope}%
\pgfpathrectangle{\pgfqpoint{0.800000in}{0.528000in}}{\pgfqpoint{4.960000in}{3.696000in}}%
\pgfusepath{clip}%
\pgfsetrectcap%
\pgfsetroundjoin%
\pgfsetlinewidth{1.505625pt}%
\definecolor{currentstroke}{rgb}{0.121569,0.466667,0.705882}%
\pgfsetstrokecolor{currentstroke}%
\pgfsetdash{}{0pt}%
\pgfpathmoveto{\pgfqpoint{1.025455in}{0.697470in}}%
\pgfpathlineto{\pgfqpoint{1.070591in}{0.705187in}}%
\pgfpathlineto{\pgfqpoint{1.115727in}{0.745612in}}%
\pgfpathlineto{\pgfqpoint{1.160863in}{0.758842in}}%
\pgfpathlineto{\pgfqpoint{1.205999in}{0.792283in}}%
\pgfpathlineto{\pgfqpoint{1.251135in}{0.855125in}}%
\pgfpathlineto{\pgfqpoint{1.296271in}{0.981911in}}%
\pgfpathlineto{\pgfqpoint{1.341407in}{1.062759in}}%
\pgfpathlineto{\pgfqpoint{1.386543in}{1.215637in}}%
\pgfpathlineto{\pgfqpoint{1.431679in}{1.374394in}}%
\pgfpathlineto{\pgfqpoint{1.476815in}{1.471413in}}%
\pgfpathlineto{\pgfqpoint{1.521952in}{1.638255in}}%
\pgfpathlineto{\pgfqpoint{1.567088in}{1.868307in}}%
\pgfpathlineto{\pgfqpoint{1.612224in}{2.061609in}}%
\pgfpathlineto{\pgfqpoint{1.657360in}{2.222204in}}%
\pgfpathlineto{\pgfqpoint{1.702496in}{2.323632in}}%
\pgfpathlineto{\pgfqpoint{1.747632in}{2.383166in}}%
\pgfpathlineto{\pgfqpoint{1.792768in}{2.565811in}}%
\pgfpathlineto{\pgfqpoint{1.837904in}{2.688186in}}%
\pgfpathlineto{\pgfqpoint{1.883040in}{2.722363in}}%
\pgfpathlineto{\pgfqpoint{1.928176in}{2.893248in}}%
\pgfpathlineto{\pgfqpoint{1.973312in}{3.031059in}}%
\pgfpathlineto{\pgfqpoint{2.018449in}{3.063766in}}%
\pgfpathlineto{\pgfqpoint{2.063585in}{3.259640in}}%
\pgfpathlineto{\pgfqpoint{2.108721in}{3.375768in}}%
\pgfpathlineto{\pgfqpoint{2.153857in}{3.392305in}}%
\pgfpathlineto{\pgfqpoint{2.198993in}{3.394510in}}%
\pgfpathlineto{\pgfqpoint{2.244129in}{3.400390in}}%
\pgfpathlineto{\pgfqpoint{2.289265in}{3.463232in}}%
\pgfpathlineto{\pgfqpoint{2.334401in}{3.569805in}}%
\pgfpathlineto{\pgfqpoint{2.379537in}{3.586710in}}%
\pgfpathlineto{\pgfqpoint{2.424673in}{3.682258in}}%
\pgfpathlineto{\pgfqpoint{2.469809in}{3.652859in}}%
\pgfpathlineto{\pgfqpoint{2.514946in}{3.747672in}}%
\pgfpathlineto{\pgfqpoint{2.560082in}{3.801326in}}%
\pgfpathlineto{\pgfqpoint{2.605218in}{3.813821in}}%
\pgfpathlineto{\pgfqpoint{2.650354in}{3.805369in}}%
\pgfpathlineto{\pgfqpoint{2.695490in}{3.848733in}}%
\pgfpathlineto{\pgfqpoint{2.740626in}{3.950529in}}%
\pgfpathlineto{\pgfqpoint{2.785762in}{3.932522in}}%
\pgfpathlineto{\pgfqpoint{2.830898in}{4.028438in}}%
\pgfpathlineto{\pgfqpoint{2.876034in}{4.001611in}}%
\pgfpathlineto{\pgfqpoint{2.921170in}{4.012636in}}%
\pgfpathlineto{\pgfqpoint{2.966306in}{4.005286in}}%
\pgfpathlineto{\pgfqpoint{3.011443in}{4.056000in}}%
\pgfpathlineto{\pgfqpoint{3.056579in}{4.014106in}}%
\pgfpathlineto{\pgfqpoint{3.101715in}{4.017046in}}%
\pgfpathlineto{\pgfqpoint{3.146851in}{3.955674in}}%
\pgfpathlineto{\pgfqpoint{3.191987in}{3.888055in}}%
\pgfpathlineto{\pgfqpoint{3.237123in}{3.846528in}}%
\pgfpathlineto{\pgfqpoint{3.282259in}{3.718640in}}%
\pgfpathlineto{\pgfqpoint{3.327395in}{3.650286in}}%
\pgfpathlineto{\pgfqpoint{3.372531in}{3.599940in}}%
\pgfpathlineto{\pgfqpoint{3.417667in}{3.535996in}}%
\pgfpathlineto{\pgfqpoint{3.462803in}{3.432362in}}%
\pgfpathlineto{\pgfqpoint{3.507940in}{3.389365in}}%
\pgfpathlineto{\pgfqpoint{3.553076in}{3.200841in}}%
\pgfpathlineto{\pgfqpoint{3.598212in}{3.069645in}}%
\pgfpathlineto{\pgfqpoint{3.643348in}{2.999454in}}%
\pgfpathlineto{\pgfqpoint{3.688484in}{2.835184in}}%
\pgfpathlineto{\pgfqpoint{3.733620in}{2.738901in}}%
\pgfpathlineto{\pgfqpoint{3.778756in}{2.562503in}}%
\pgfpathlineto{\pgfqpoint{3.823892in}{2.360749in}}%
\pgfpathlineto{\pgfqpoint{3.869028in}{2.192804in}}%
\pgfpathlineto{\pgfqpoint{3.914164in}{2.034414in}}%
\pgfpathlineto{\pgfqpoint{3.959300in}{1.863897in}}%
\pgfpathlineto{\pgfqpoint{4.004437in}{1.760263in}}%
\pgfpathlineto{\pgfqpoint{4.049573in}{1.718737in}}%
\pgfpathlineto{\pgfqpoint{4.094709in}{1.586806in}}%
\pgfpathlineto{\pgfqpoint{4.139845in}{1.462226in}}%
\pgfpathlineto{\pgfqpoint{4.184981in}{1.365942in}}%
\pgfpathlineto{\pgfqpoint{4.230117in}{1.264881in}}%
\pgfpathlineto{\pgfqpoint{4.275253in}{1.148385in}}%
\pgfpathlineto{\pgfqpoint{4.320389in}{1.100979in}}%
\pgfpathlineto{\pgfqpoint{4.365525in}{1.048794in}}%
\pgfpathlineto{\pgfqpoint{4.410661in}{0.984850in}}%
\pgfpathlineto{\pgfqpoint{4.455797in}{0.944794in}}%
\pgfpathlineto{\pgfqpoint{4.500934in}{0.884525in}}%
\pgfpathlineto{\pgfqpoint{4.546070in}{0.867252in}}%
\pgfpathlineto{\pgfqpoint{4.591206in}{0.810291in}}%
\pgfpathlineto{\pgfqpoint{4.636342in}{0.789344in}}%
\pgfpathlineto{\pgfqpoint{4.681478in}{0.774644in}}%
\pgfpathlineto{\pgfqpoint{4.726614in}{0.760311in}}%
\pgfpathlineto{\pgfqpoint{4.771750in}{0.741937in}}%
\pgfpathlineto{\pgfqpoint{4.816886in}{0.729809in}}%
\pgfpathlineto{\pgfqpoint{4.862022in}{0.726870in}}%
\pgfpathlineto{\pgfqpoint{4.907158in}{0.722460in}}%
\pgfpathlineto{\pgfqpoint{4.952294in}{0.714007in}}%
\pgfpathlineto{\pgfqpoint{4.997430in}{0.714007in}}%
\pgfpathlineto{\pgfqpoint{5.042567in}{0.708862in}}%
\pgfpathlineto{\pgfqpoint{5.087703in}{0.705555in}}%
\pgfpathlineto{\pgfqpoint{5.132839in}{0.699675in}}%
\pgfpathlineto{\pgfqpoint{5.177975in}{0.700042in}}%
\pgfpathlineto{\pgfqpoint{5.223111in}{0.700042in}}%
\pgfpathlineto{\pgfqpoint{5.268247in}{0.696735in}}%
\pgfpathlineto{\pgfqpoint{5.313383in}{0.697102in}}%
\pgfpathlineto{\pgfqpoint{5.358519in}{0.696000in}}%
\pgfpathlineto{\pgfqpoint{5.403655in}{0.696367in}}%
\pgfpathlineto{\pgfqpoint{5.448791in}{0.696735in}}%
\pgfpathlineto{\pgfqpoint{5.493927in}{0.696367in}}%
\pgfusepath{stroke}%
\end{pgfscope}%
\begin{pgfscope}%
\pgfpathrectangle{\pgfqpoint{0.800000in}{0.528000in}}{\pgfqpoint{4.960000in}{3.696000in}}%
\pgfusepath{clip}%
\pgfsetrectcap%
\pgfsetroundjoin%
\pgfsetlinewidth{1.505625pt}%
\definecolor{currentstroke}{rgb}{1.000000,0.498039,0.054902}%
\pgfsetstrokecolor{currentstroke}%
\pgfsetdash{}{0pt}%
\pgfpathmoveto{\pgfqpoint{1.025455in}{0.700146in}}%
\pgfpathlineto{\pgfqpoint{1.052563in}{0.697080in}}%
\pgfpathlineto{\pgfqpoint{1.066118in}{0.697293in}}%
\pgfpathlineto{\pgfqpoint{1.079672in}{0.699709in}}%
\pgfpathlineto{\pgfqpoint{1.088708in}{0.703074in}}%
\pgfpathlineto{\pgfqpoint{1.097745in}{0.708241in}}%
\pgfpathlineto{\pgfqpoint{1.106781in}{0.715391in}}%
\pgfpathlineto{\pgfqpoint{1.124853in}{0.733866in}}%
\pgfpathlineto{\pgfqpoint{1.138408in}{0.746303in}}%
\pgfpathlineto{\pgfqpoint{1.147444in}{0.751436in}}%
\pgfpathlineto{\pgfqpoint{1.156480in}{0.753992in}}%
\pgfpathlineto{\pgfqpoint{1.179071in}{0.758058in}}%
\pgfpathlineto{\pgfqpoint{1.188107in}{0.762205in}}%
\pgfpathlineto{\pgfqpoint{1.201662in}{0.771529in}}%
\pgfpathlineto{\pgfqpoint{1.237807in}{0.800760in}}%
\pgfpathlineto{\pgfqpoint{1.246843in}{0.810411in}}%
\pgfpathlineto{\pgfqpoint{1.255879in}{0.823655in}}%
\pgfpathlineto{\pgfqpoint{1.264915in}{0.841890in}}%
\pgfpathlineto{\pgfqpoint{1.273952in}{0.865295in}}%
\pgfpathlineto{\pgfqpoint{1.305578in}{0.956250in}}%
\pgfpathlineto{\pgfqpoint{1.314615in}{0.975402in}}%
\pgfpathlineto{\pgfqpoint{1.328169in}{0.998723in}}%
\pgfpathlineto{\pgfqpoint{1.350760in}{1.035825in}}%
\pgfpathlineto{\pgfqpoint{1.364314in}{1.063661in}}%
\pgfpathlineto{\pgfqpoint{1.373350in}{1.086703in}}%
\pgfpathlineto{\pgfqpoint{1.382387in}{1.114601in}}%
\pgfpathlineto{\pgfqpoint{1.395941in}{1.166695in}}%
\pgfpathlineto{\pgfqpoint{1.432086in}{1.323394in}}%
\pgfpathlineto{\pgfqpoint{1.441122in}{1.350332in}}%
\pgfpathlineto{\pgfqpoint{1.450158in}{1.370129in}}%
\pgfpathlineto{\pgfqpoint{1.463713in}{1.391505in}}%
\pgfpathlineto{\pgfqpoint{1.472749in}{1.405710in}}%
\pgfpathlineto{\pgfqpoint{1.481785in}{1.424380in}}%
\pgfpathlineto{\pgfqpoint{1.490822in}{1.449414in}}%
\pgfpathlineto{\pgfqpoint{1.504376in}{1.497355in}}%
\pgfpathlineto{\pgfqpoint{1.545039in}{1.660212in}}%
\pgfpathlineto{\pgfqpoint{1.563112in}{1.743501in}}%
\pgfpathlineto{\pgfqpoint{1.603775in}{1.936163in}}%
\pgfpathlineto{\pgfqpoint{1.621847in}{2.009368in}}%
\pgfpathlineto{\pgfqpoint{1.653474in}{2.131702in}}%
\pgfpathlineto{\pgfqpoint{1.689619in}{2.286016in}}%
\pgfpathlineto{\pgfqpoint{1.698655in}{2.309861in}}%
\pgfpathlineto{\pgfqpoint{1.707692in}{2.323967in}}%
\pgfpathlineto{\pgfqpoint{1.716728in}{2.330691in}}%
\pgfpathlineto{\pgfqpoint{1.739319in}{2.340115in}}%
\pgfpathlineto{\pgfqpoint{1.748355in}{2.347963in}}%
\pgfpathlineto{\pgfqpoint{1.757391in}{2.360086in}}%
\pgfpathlineto{\pgfqpoint{1.766427in}{2.377800in}}%
\pgfpathlineto{\pgfqpoint{1.775464in}{2.402347in}}%
\pgfpathlineto{\pgfqpoint{1.784500in}{2.434112in}}%
\pgfpathlineto{\pgfqpoint{1.798054in}{2.492811in}}%
\pgfpathlineto{\pgfqpoint{1.816127in}{2.574165in}}%
\pgfpathlineto{\pgfqpoint{1.825163in}{2.606368in}}%
\pgfpathlineto{\pgfqpoint{1.834199in}{2.629867in}}%
\pgfpathlineto{\pgfqpoint{1.843236in}{2.646425in}}%
\pgfpathlineto{\pgfqpoint{1.861308in}{2.671466in}}%
\pgfpathlineto{\pgfqpoint{1.879381in}{2.696789in}}%
\pgfpathlineto{\pgfqpoint{1.888417in}{2.712516in}}%
\pgfpathlineto{\pgfqpoint{1.897453in}{2.732995in}}%
\pgfpathlineto{\pgfqpoint{1.911007in}{2.772821in}}%
\pgfpathlineto{\pgfqpoint{1.983298in}{3.004902in}}%
\pgfpathlineto{\pgfqpoint{1.987816in}{3.013585in}}%
\pgfpathlineto{\pgfqpoint{1.992334in}{3.019279in}}%
\pgfpathlineto{\pgfqpoint{1.996852in}{3.021986in}}%
\pgfpathlineto{\pgfqpoint{2.001370in}{3.022127in}}%
\pgfpathlineto{\pgfqpoint{2.014924in}{3.017076in}}%
\pgfpathlineto{\pgfqpoint{2.019443in}{3.017867in}}%
\pgfpathlineto{\pgfqpoint{2.023961in}{3.021885in}}%
\pgfpathlineto{\pgfqpoint{2.028479in}{3.029878in}}%
\pgfpathlineto{\pgfqpoint{2.032997in}{3.042149in}}%
\pgfpathlineto{\pgfqpoint{2.042033in}{3.078462in}}%
\pgfpathlineto{\pgfqpoint{2.055588in}{3.150633in}}%
\pgfpathlineto{\pgfqpoint{2.078178in}{3.272195in}}%
\pgfpathlineto{\pgfqpoint{2.087214in}{3.311972in}}%
\pgfpathlineto{\pgfqpoint{2.096251in}{3.342969in}}%
\pgfpathlineto{\pgfqpoint{2.105287in}{3.364147in}}%
\pgfpathlineto{\pgfqpoint{2.114323in}{3.376201in}}%
\pgfpathlineto{\pgfqpoint{2.123359in}{3.381641in}}%
\pgfpathlineto{\pgfqpoint{2.145950in}{3.387589in}}%
\pgfpathlineto{\pgfqpoint{2.154986in}{3.388912in}}%
\pgfpathlineto{\pgfqpoint{2.159504in}{3.388364in}}%
\pgfpathlineto{\pgfqpoint{2.168541in}{3.383870in}}%
\pgfpathlineto{\pgfqpoint{2.191131in}{3.364289in}}%
\pgfpathlineto{\pgfqpoint{2.195649in}{3.362628in}}%
\pgfpathlineto{\pgfqpoint{2.200168in}{3.362492in}}%
\pgfpathlineto{\pgfqpoint{2.204686in}{3.363913in}}%
\pgfpathlineto{\pgfqpoint{2.213722in}{3.370771in}}%
\pgfpathlineto{\pgfqpoint{2.236313in}{3.393785in}}%
\pgfpathlineto{\pgfqpoint{2.249867in}{3.402068in}}%
\pgfpathlineto{\pgfqpoint{2.263421in}{3.410484in}}%
\pgfpathlineto{\pgfqpoint{2.272458in}{3.418940in}}%
\pgfpathlineto{\pgfqpoint{2.290530in}{3.440334in}}%
\pgfpathlineto{\pgfqpoint{2.299566in}{3.453241in}}%
\pgfpathlineto{\pgfqpoint{2.308603in}{3.469631in}}%
\pgfpathlineto{\pgfqpoint{2.322157in}{3.501231in}}%
\pgfpathlineto{\pgfqpoint{2.344748in}{3.557090in}}%
\pgfpathlineto{\pgfqpoint{2.353784in}{3.575271in}}%
\pgfpathlineto{\pgfqpoint{2.362820in}{3.589506in}}%
\pgfpathlineto{\pgfqpoint{2.371856in}{3.598734in}}%
\pgfpathlineto{\pgfqpoint{2.376375in}{3.601000in}}%
\pgfpathlineto{\pgfqpoint{2.380893in}{3.601574in}}%
\pgfpathlineto{\pgfqpoint{2.385411in}{3.600591in}}%
\pgfpathlineto{\pgfqpoint{2.394447in}{3.595446in}}%
\pgfpathlineto{\pgfqpoint{2.403483in}{3.589861in}}%
\pgfpathlineto{\pgfqpoint{2.408001in}{3.588579in}}%
\pgfpathlineto{\pgfqpoint{2.412520in}{3.589190in}}%
\pgfpathlineto{\pgfqpoint{2.417038in}{3.592215in}}%
\pgfpathlineto{\pgfqpoint{2.421556in}{3.597941in}}%
\pgfpathlineto{\pgfqpoint{2.430592in}{3.617012in}}%
\pgfpathlineto{\pgfqpoint{2.453183in}{3.677500in}}%
\pgfpathlineto{\pgfqpoint{2.457701in}{3.685230in}}%
\pgfpathlineto{\pgfqpoint{2.462219in}{3.690274in}}%
\pgfpathlineto{\pgfqpoint{2.466737in}{3.692605in}}%
\pgfpathlineto{\pgfqpoint{2.471255in}{3.692449in}}%
\pgfpathlineto{\pgfqpoint{2.475773in}{3.690237in}}%
\pgfpathlineto{\pgfqpoint{2.484810in}{3.682045in}}%
\pgfpathlineto{\pgfqpoint{2.493846in}{3.673563in}}%
\pgfpathlineto{\pgfqpoint{2.498364in}{3.671051in}}%
\pgfpathlineto{\pgfqpoint{2.502882in}{3.670565in}}%
\pgfpathlineto{\pgfqpoint{2.507400in}{3.672613in}}%
\pgfpathlineto{\pgfqpoint{2.511918in}{3.677543in}}%
\pgfpathlineto{\pgfqpoint{2.516437in}{3.685497in}}%
\pgfpathlineto{\pgfqpoint{2.525473in}{3.709691in}}%
\pgfpathlineto{\pgfqpoint{2.548063in}{3.783737in}}%
\pgfpathlineto{\pgfqpoint{2.557100in}{3.801585in}}%
\pgfpathlineto{\pgfqpoint{2.561618in}{3.806665in}}%
\pgfpathlineto{\pgfqpoint{2.566136in}{3.809657in}}%
\pgfpathlineto{\pgfqpoint{2.575172in}{3.811975in}}%
\pgfpathlineto{\pgfqpoint{2.593245in}{3.815504in}}%
\pgfpathlineto{\pgfqpoint{2.602281in}{3.816734in}}%
\pgfpathlineto{\pgfqpoint{2.611317in}{3.815405in}}%
\pgfpathlineto{\pgfqpoint{2.620353in}{3.811121in}}%
\pgfpathlineto{\pgfqpoint{2.642944in}{3.796336in}}%
\pgfpathlineto{\pgfqpoint{2.647462in}{3.795713in}}%
\pgfpathlineto{\pgfqpoint{2.651980in}{3.796796in}}%
\pgfpathlineto{\pgfqpoint{2.656498in}{3.799730in}}%
\pgfpathlineto{\pgfqpoint{2.665535in}{3.810796in}}%
\pgfpathlineto{\pgfqpoint{2.688125in}{3.848267in}}%
\pgfpathlineto{\pgfqpoint{2.692643in}{3.852669in}}%
\pgfpathlineto{\pgfqpoint{2.697162in}{3.855155in}}%
\pgfpathlineto{\pgfqpoint{2.701680in}{3.855853in}}%
\pgfpathlineto{\pgfqpoint{2.719752in}{3.853260in}}%
\pgfpathlineto{\pgfqpoint{2.724270in}{3.855770in}}%
\pgfpathlineto{\pgfqpoint{2.728788in}{3.861074in}}%
\pgfpathlineto{\pgfqpoint{2.733307in}{3.869327in}}%
\pgfpathlineto{\pgfqpoint{2.742343in}{3.892964in}}%
\pgfpathlineto{\pgfqpoint{2.755897in}{3.932307in}}%
\pgfpathlineto{\pgfqpoint{2.764934in}{3.950575in}}%
\pgfpathlineto{\pgfqpoint{2.769452in}{3.956060in}}%
\pgfpathlineto{\pgfqpoint{2.773970in}{3.959225in}}%
\pgfpathlineto{\pgfqpoint{2.778488in}{3.960381in}}%
\pgfpathlineto{\pgfqpoint{2.783006in}{3.959932in}}%
\pgfpathlineto{\pgfqpoint{2.792042in}{3.956012in}}%
\pgfpathlineto{\pgfqpoint{2.805597in}{3.949359in}}%
\pgfpathlineto{\pgfqpoint{2.810115in}{3.948620in}}%
\pgfpathlineto{\pgfqpoint{2.814633in}{3.949217in}}%
\pgfpathlineto{\pgfqpoint{2.819151in}{3.951409in}}%
\pgfpathlineto{\pgfqpoint{2.823669in}{3.955319in}}%
\pgfpathlineto{\pgfqpoint{2.832705in}{3.967689in}}%
\pgfpathlineto{\pgfqpoint{2.850778in}{3.996954in}}%
\pgfpathlineto{\pgfqpoint{2.859814in}{4.007329in}}%
\pgfpathlineto{\pgfqpoint{2.873369in}{4.017631in}}%
\pgfpathlineto{\pgfqpoint{2.904995in}{4.038002in}}%
\pgfpathlineto{\pgfqpoint{2.909514in}{4.038053in}}%
\pgfpathlineto{\pgfqpoint{2.914032in}{4.036779in}}%
\pgfpathlineto{\pgfqpoint{2.923068in}{4.030970in}}%
\pgfpathlineto{\pgfqpoint{2.941140in}{4.016838in}}%
\pgfpathlineto{\pgfqpoint{2.950177in}{4.012821in}}%
\pgfpathlineto{\pgfqpoint{2.968249in}{4.008235in}}%
\pgfpathlineto{\pgfqpoint{2.977285in}{4.008749in}}%
\pgfpathlineto{\pgfqpoint{2.981804in}{4.010768in}}%
\pgfpathlineto{\pgfqpoint{2.990840in}{4.018918in}}%
\pgfpathlineto{\pgfqpoint{3.008912in}{4.041118in}}%
\pgfpathlineto{\pgfqpoint{3.017949in}{4.047042in}}%
\pgfpathlineto{\pgfqpoint{3.026985in}{4.047842in}}%
\pgfpathlineto{\pgfqpoint{3.036021in}{4.045120in}}%
\pgfpathlineto{\pgfqpoint{3.045057in}{4.040170in}}%
\pgfpathlineto{\pgfqpoint{3.058612in}{4.029817in}}%
\pgfpathlineto{\pgfqpoint{3.076684in}{4.015453in}}%
\pgfpathlineto{\pgfqpoint{3.085721in}{4.010512in}}%
\pgfpathlineto{\pgfqpoint{3.094757in}{4.008665in}}%
\pgfpathlineto{\pgfqpoint{3.103793in}{4.010933in}}%
\pgfpathlineto{\pgfqpoint{3.117347in}{4.019505in}}%
\pgfpathlineto{\pgfqpoint{3.126384in}{4.024191in}}%
\pgfpathlineto{\pgfqpoint{3.130902in}{4.025013in}}%
\pgfpathlineto{\pgfqpoint{3.135420in}{4.024379in}}%
\pgfpathlineto{\pgfqpoint{3.139938in}{4.022032in}}%
\pgfpathlineto{\pgfqpoint{3.144456in}{4.017802in}}%
\pgfpathlineto{\pgfqpoint{3.153492in}{4.003624in}}%
\pgfpathlineto{\pgfqpoint{3.162529in}{3.982932in}}%
\pgfpathlineto{\pgfqpoint{3.194156in}{3.900442in}}%
\pgfpathlineto{\pgfqpoint{3.203192in}{3.883512in}}%
\pgfpathlineto{\pgfqpoint{3.212228in}{3.871562in}}%
\pgfpathlineto{\pgfqpoint{3.221264in}{3.864819in}}%
\pgfpathlineto{\pgfqpoint{3.230301in}{3.862751in}}%
\pgfpathlineto{\pgfqpoint{3.252891in}{3.864320in}}%
\pgfpathlineto{\pgfqpoint{3.257409in}{3.861749in}}%
\pgfpathlineto{\pgfqpoint{3.261927in}{3.857007in}}%
\pgfpathlineto{\pgfqpoint{3.266446in}{3.849848in}}%
\pgfpathlineto{\pgfqpoint{3.275482in}{3.828299in}}%
\pgfpathlineto{\pgfqpoint{3.284518in}{3.798574in}}%
\pgfpathlineto{\pgfqpoint{3.311627in}{3.697995in}}%
\pgfpathlineto{\pgfqpoint{3.320663in}{3.675536in}}%
\pgfpathlineto{\pgfqpoint{3.329699in}{3.661667in}}%
\pgfpathlineto{\pgfqpoint{3.338736in}{3.654428in}}%
\pgfpathlineto{\pgfqpoint{3.356808in}{3.645113in}}%
\pgfpathlineto{\pgfqpoint{3.365844in}{3.637529in}}%
\pgfpathlineto{\pgfqpoint{3.374881in}{3.627044in}}%
\pgfpathlineto{\pgfqpoint{3.383917in}{3.613616in}}%
\pgfpathlineto{\pgfqpoint{3.397471in}{3.588582in}}%
\pgfpathlineto{\pgfqpoint{3.415544in}{3.553744in}}%
\pgfpathlineto{\pgfqpoint{3.424580in}{3.540735in}}%
\pgfpathlineto{\pgfqpoint{3.433616in}{3.531658in}}%
\pgfpathlineto{\pgfqpoint{3.447171in}{3.520047in}}%
\pgfpathlineto{\pgfqpoint{3.456207in}{3.507674in}}%
\pgfpathlineto{\pgfqpoint{3.465243in}{3.488418in}}%
\pgfpathlineto{\pgfqpoint{3.478798in}{3.449882in}}%
\pgfpathlineto{\pgfqpoint{3.487834in}{3.425009in}}%
\pgfpathlineto{\pgfqpoint{3.496870in}{3.406824in}}%
\pgfpathlineto{\pgfqpoint{3.501388in}{3.401080in}}%
\pgfpathlineto{\pgfqpoint{3.505906in}{3.397612in}}%
\pgfpathlineto{\pgfqpoint{3.510424in}{3.396158in}}%
\pgfpathlineto{\pgfqpoint{3.528497in}{3.395202in}}%
\pgfpathlineto{\pgfqpoint{3.533015in}{3.390737in}}%
\pgfpathlineto{\pgfqpoint{3.537533in}{3.382584in}}%
\pgfpathlineto{\pgfqpoint{3.542051in}{3.370280in}}%
\pgfpathlineto{\pgfqpoint{3.551088in}{3.333769in}}%
\pgfpathlineto{\pgfqpoint{3.569160in}{3.236923in}}%
\pgfpathlineto{\pgfqpoint{3.582715in}{3.171833in}}%
\pgfpathlineto{\pgfqpoint{3.591751in}{3.137551in}}%
\pgfpathlineto{\pgfqpoint{3.600787in}{3.109734in}}%
\pgfpathlineto{\pgfqpoint{3.609823in}{3.087301in}}%
\pgfpathlineto{\pgfqpoint{3.623378in}{3.061370in}}%
\pgfpathlineto{\pgfqpoint{3.655005in}{3.007690in}}%
\pgfpathlineto{\pgfqpoint{3.664041in}{2.986585in}}%
\pgfpathlineto{\pgfqpoint{3.677595in}{2.946435in}}%
\pgfpathlineto{\pgfqpoint{3.700186in}{2.874487in}}%
\pgfpathlineto{\pgfqpoint{3.713740in}{2.841031in}}%
\pgfpathlineto{\pgfqpoint{3.731813in}{2.800553in}}%
\pgfpathlineto{\pgfqpoint{3.745367in}{2.763151in}}%
\pgfpathlineto{\pgfqpoint{3.758921in}{2.717910in}}%
\pgfpathlineto{\pgfqpoint{3.781512in}{2.631222in}}%
\pgfpathlineto{\pgfqpoint{3.826693in}{2.455846in}}%
\pgfpathlineto{\pgfqpoint{3.849284in}{2.349130in}}%
\pgfpathlineto{\pgfqpoint{3.867357in}{2.268963in}}%
\pgfpathlineto{\pgfqpoint{3.880911in}{2.218811in}}%
\pgfpathlineto{\pgfqpoint{3.921574in}{2.083874in}}%
\pgfpathlineto{\pgfqpoint{3.944165in}{2.008810in}}%
\pgfpathlineto{\pgfqpoint{3.962237in}{1.938199in}}%
\pgfpathlineto{\pgfqpoint{3.989346in}{1.828249in}}%
\pgfpathlineto{\pgfqpoint{3.998382in}{1.801285in}}%
\pgfpathlineto{\pgfqpoint{4.007418in}{1.782544in}}%
\pgfpathlineto{\pgfqpoint{4.016455in}{1.771012in}}%
\pgfpathlineto{\pgfqpoint{4.025491in}{1.764480in}}%
\pgfpathlineto{\pgfqpoint{4.043563in}{1.754793in}}%
\pgfpathlineto{\pgfqpoint{4.052600in}{1.745895in}}%
\pgfpathlineto{\pgfqpoint{4.061636in}{1.731183in}}%
\pgfpathlineto{\pgfqpoint{4.070672in}{1.709981in}}%
\pgfpathlineto{\pgfqpoint{4.084227in}{1.669527in}}%
\pgfpathlineto{\pgfqpoint{4.102299in}{1.614837in}}%
\pgfpathlineto{\pgfqpoint{4.115854in}{1.581524in}}%
\pgfpathlineto{\pgfqpoint{4.138444in}{1.530737in}}%
\pgfpathlineto{\pgfqpoint{4.156517in}{1.481676in}}%
\pgfpathlineto{\pgfqpoint{4.179107in}{1.421123in}}%
\pgfpathlineto{\pgfqpoint{4.197180in}{1.380733in}}%
\pgfpathlineto{\pgfqpoint{4.237843in}{1.297418in}}%
\pgfpathlineto{\pgfqpoint{4.251397in}{1.261075in}}%
\pgfpathlineto{\pgfqpoint{4.283024in}{1.171021in}}%
\pgfpathlineto{\pgfqpoint{4.292060in}{1.151270in}}%
\pgfpathlineto{\pgfqpoint{4.301097in}{1.136562in}}%
\pgfpathlineto{\pgfqpoint{4.310133in}{1.126206in}}%
\pgfpathlineto{\pgfqpoint{4.323687in}{1.115175in}}%
\pgfpathlineto{\pgfqpoint{4.355314in}{1.092847in}}%
\pgfpathlineto{\pgfqpoint{4.364351in}{1.083436in}}%
\pgfpathlineto{\pgfqpoint{4.373387in}{1.071372in}}%
\pgfpathlineto{\pgfqpoint{4.386941in}{1.048891in}}%
\pgfpathlineto{\pgfqpoint{4.409532in}{1.008438in}}%
\pgfpathlineto{\pgfqpoint{4.418568in}{0.995442in}}%
\pgfpathlineto{\pgfqpoint{4.427604in}{0.985209in}}%
\pgfpathlineto{\pgfqpoint{4.441159in}{0.973755in}}%
\pgfpathlineto{\pgfqpoint{4.463749in}{0.956152in}}%
\pgfpathlineto{\pgfqpoint{4.477304in}{0.942022in}}%
\pgfpathlineto{\pgfqpoint{4.495376in}{0.919042in}}%
\pgfpathlineto{\pgfqpoint{4.517967in}{0.890236in}}%
\pgfpathlineto{\pgfqpoint{4.527003in}{0.881722in}}%
\pgfpathlineto{\pgfqpoint{4.536039in}{0.875898in}}%
\pgfpathlineto{\pgfqpoint{4.545076in}{0.872583in}}%
\pgfpathlineto{\pgfqpoint{4.563148in}{0.867765in}}%
\pgfpathlineto{\pgfqpoint{4.572184in}{0.862154in}}%
\pgfpathlineto{\pgfqpoint{4.581221in}{0.852860in}}%
\pgfpathlineto{\pgfqpoint{4.594775in}{0.834066in}}%
\pgfpathlineto{\pgfqpoint{4.608329in}{0.815667in}}%
\pgfpathlineto{\pgfqpoint{4.617366in}{0.806518in}}%
\pgfpathlineto{\pgfqpoint{4.626402in}{0.800465in}}%
\pgfpathlineto{\pgfqpoint{4.639956in}{0.795557in}}%
\pgfpathlineto{\pgfqpoint{4.658029in}{0.790200in}}%
\pgfpathlineto{\pgfqpoint{4.680619in}{0.780019in}}%
\pgfpathlineto{\pgfqpoint{4.698692in}{0.772646in}}%
\pgfpathlineto{\pgfqpoint{4.716764in}{0.768501in}}%
\pgfpathlineto{\pgfqpoint{4.743873in}{0.763057in}}%
\pgfpathlineto{\pgfqpoint{4.761946in}{0.757278in}}%
\pgfpathlineto{\pgfqpoint{4.775500in}{0.750848in}}%
\pgfpathlineto{\pgfqpoint{4.807127in}{0.733925in}}%
\pgfpathlineto{\pgfqpoint{4.820681in}{0.730252in}}%
\pgfpathlineto{\pgfqpoint{4.838754in}{0.728515in}}%
\pgfpathlineto{\pgfqpoint{4.892971in}{0.726203in}}%
\pgfpathlineto{\pgfqpoint{4.965261in}{0.716436in}}%
\pgfpathlineto{\pgfqpoint{4.987852in}{0.711945in}}%
\pgfpathlineto{\pgfqpoint{5.001406in}{0.711238in}}%
\pgfpathlineto{\pgfqpoint{5.046588in}{0.712342in}}%
\pgfpathlineto{\pgfqpoint{5.100805in}{0.705947in}}%
\pgfpathlineto{\pgfqpoint{5.123396in}{0.704195in}}%
\pgfpathlineto{\pgfqpoint{5.168577in}{0.698533in}}%
\pgfpathlineto{\pgfqpoint{5.191168in}{0.699420in}}%
\pgfpathlineto{\pgfqpoint{5.222795in}{0.701163in}}%
\pgfpathlineto{\pgfqpoint{5.249903in}{0.699731in}}%
\pgfpathlineto{\pgfqpoint{5.295085in}{0.696704in}}%
\pgfpathlineto{\pgfqpoint{5.457737in}{0.696306in}}%
\pgfpathlineto{\pgfqpoint{5.534545in}{0.696363in}}%
\pgfpathlineto{\pgfqpoint{5.534545in}{0.696363in}}%
\pgfusepath{stroke}%
\end{pgfscope}%
\begin{pgfscope}%
\pgfsetrectcap%
\pgfsetmiterjoin%
\pgfsetlinewidth{0.803000pt}%
\definecolor{currentstroke}{rgb}{0.000000,0.000000,0.000000}%
\pgfsetstrokecolor{currentstroke}%
\pgfsetdash{}{0pt}%
\pgfpathmoveto{\pgfqpoint{0.800000in}{0.528000in}}%
\pgfpathlineto{\pgfqpoint{0.800000in}{4.224000in}}%
\pgfusepath{stroke}%
\end{pgfscope}%
\begin{pgfscope}%
\pgfsetrectcap%
\pgfsetmiterjoin%
\pgfsetlinewidth{0.803000pt}%
\definecolor{currentstroke}{rgb}{0.000000,0.000000,0.000000}%
\pgfsetstrokecolor{currentstroke}%
\pgfsetdash{}{0pt}%
\pgfpathmoveto{\pgfqpoint{5.760000in}{0.528000in}}%
\pgfpathlineto{\pgfqpoint{5.760000in}{4.224000in}}%
\pgfusepath{stroke}%
\end{pgfscope}%
\begin{pgfscope}%
\pgfsetrectcap%
\pgfsetmiterjoin%
\pgfsetlinewidth{0.803000pt}%
\definecolor{currentstroke}{rgb}{0.000000,0.000000,0.000000}%
\pgfsetstrokecolor{currentstroke}%
\pgfsetdash{}{0pt}%
\pgfpathmoveto{\pgfqpoint{0.800000in}{0.528000in}}%
\pgfpathlineto{\pgfqpoint{5.760000in}{0.528000in}}%
\pgfusepath{stroke}%
\end{pgfscope}%
\begin{pgfscope}%
\pgfsetrectcap%
\pgfsetmiterjoin%
\pgfsetlinewidth{0.803000pt}%
\definecolor{currentstroke}{rgb}{0.000000,0.000000,0.000000}%
\pgfsetstrokecolor{currentstroke}%
\pgfsetdash{}{0pt}%
\pgfpathmoveto{\pgfqpoint{0.800000in}{4.224000in}}%
\pgfpathlineto{\pgfqpoint{5.760000in}{4.224000in}}%
\pgfusepath{stroke}%
\end{pgfscope}%
\begin{pgfscope}%
\pgfsetbuttcap%
\pgfsetmiterjoin%
\definecolor{currentfill}{rgb}{1.000000,1.000000,1.000000}%
\pgfsetfillcolor{currentfill}%
\pgfsetfillopacity{0.800000}%
\pgfsetlinewidth{1.003750pt}%
\definecolor{currentstroke}{rgb}{0.800000,0.800000,0.800000}%
\pgfsetstrokecolor{currentstroke}%
\pgfsetstrokeopacity{0.800000}%
\pgfsetdash{}{0pt}%
\pgfpathmoveto{\pgfqpoint{4.248289in}{3.209490in}}%
\pgfpathlineto{\pgfqpoint{5.623889in}{3.209490in}}%
\pgfpathquadraticcurveto{\pgfqpoint{5.662778in}{3.209490in}}{\pgfqpoint{5.662778in}{3.248379in}}%
\pgfpathlineto{\pgfqpoint{5.662778in}{4.087889in}}%
\pgfpathquadraticcurveto{\pgfqpoint{5.662778in}{4.126778in}}{\pgfqpoint{5.623889in}{4.126778in}}%
\pgfpathlineto{\pgfqpoint{4.248289in}{4.126778in}}%
\pgfpathquadraticcurveto{\pgfqpoint{4.209401in}{4.126778in}}{\pgfqpoint{4.209401in}{4.087889in}}%
\pgfpathlineto{\pgfqpoint{4.209401in}{3.248379in}}%
\pgfpathquadraticcurveto{\pgfqpoint{4.209401in}{3.209490in}}{\pgfqpoint{4.248289in}{3.209490in}}%
\pgfpathclose%
\pgfusepath{stroke,fill}%
\end{pgfscope}%
\begin{pgfscope}%
\definecolor{textcolor}{rgb}{0.000000,0.000000,0.000000}%
\pgfsetstrokecolor{textcolor}%
\pgfsetfillcolor{textcolor}%
\pgftext[x=4.556552in,y=3.901268in,left,base]{\color{textcolor}\rmfamily\fontsize{14.000000}{16.800000}\selectfont Method}%
\end{pgfscope}%
\begin{pgfscope}%
\pgfsetrectcap%
\pgfsetroundjoin%
\pgfsetlinewidth{1.505625pt}%
\definecolor{currentstroke}{rgb}{0.121569,0.466667,0.705882}%
\pgfsetstrokecolor{currentstroke}%
\pgfsetdash{}{0pt}%
\pgfpathmoveto{\pgfqpoint{4.287178in}{3.683923in}}%
\pgfpathlineto{\pgfqpoint{4.676067in}{3.683923in}}%
\pgfusepath{stroke}%
\end{pgfscope}%
\begin{pgfscope}%
\definecolor{textcolor}{rgb}{0.000000,0.000000,0.000000}%
\pgfsetstrokecolor{textcolor}%
\pgfsetfillcolor{textcolor}%
\pgftext[x=4.831623in,y=3.615868in,left,base]{\color{textcolor}\rmfamily\fontsize{14.000000}{16.800000}\selectfont Numpy}%
\end{pgfscope}%
\begin{pgfscope}%
\pgfsetrectcap%
\pgfsetroundjoin%
\pgfsetlinewidth{1.505625pt}%
\definecolor{currentstroke}{rgb}{1.000000,0.498039,0.054902}%
\pgfsetstrokecolor{currentstroke}%
\pgfsetdash{}{0pt}%
\pgfpathmoveto{\pgfqpoint{4.287178in}{3.395769in}}%
\pgfpathlineto{\pgfqpoint{4.676067in}{3.395769in}}%
\pgfusepath{stroke}%
\end{pgfscope}%
\begin{pgfscope}%
\definecolor{textcolor}{rgb}{0.000000,0.000000,0.000000}%
\pgfsetstrokecolor{textcolor}%
\pgfsetfillcolor{textcolor}%
\pgftext[x=4.831623in,y=3.327714in,left,base]{\color{textcolor}\rmfamily\fontsize{14.000000}{16.800000}\selectfont Custom}%
\end{pgfscope}%
\end{pgfpicture}%
\makeatother%
\endgroup%
} & \resizebox{0.47\linewidth}{!}{%% Creator: Matplotlib, PGF backend
%%
%% To include the figure in your LaTeX document, write
%%   \input{<filename>.pgf}
%%
%% Make sure the required packages are loaded in your preamble
%%   \usepackage{pgf}
%%
%% Figures using additional raster images can only be included by \input if
%% they are in the same directory as the main LaTeX file. For loading figures
%% from other directories you can use the `import` package
%%   \usepackage{import}
%% and then include the figures with
%%   \import{<path to file>}{<filename>.pgf}
%%
%% Matplotlib used the following preamble
%%   \usepackage{fontspec}
%%   \setmainfont{DejaVuSerif.ttf}[Path=C:/Users/garaya/AppData/Local/Continuum/anaconda3/lib/site-packages/matplotlib/mpl-data/fonts/ttf/]
%%   \setsansfont{DejaVuSans.ttf}[Path=C:/Users/garaya/AppData/Local/Continuum/anaconda3/lib/site-packages/matplotlib/mpl-data/fonts/ttf/]
%%   \setmonofont{DejaVuSansMono.ttf}[Path=C:/Users/garaya/AppData/Local/Continuum/anaconda3/lib/site-packages/matplotlib/mpl-data/fonts/ttf/]
%%
\begingroup%
\makeatletter%
\begin{pgfpicture}%
\pgfpathrectangle{\pgfpointorigin}{\pgfqpoint{6.400000in}{4.800000in}}%
\pgfusepath{use as bounding box, clip}%
\begin{pgfscope}%
\pgfsetbuttcap%
\pgfsetmiterjoin%
\definecolor{currentfill}{rgb}{1.000000,1.000000,1.000000}%
\pgfsetfillcolor{currentfill}%
\pgfsetlinewidth{0.000000pt}%
\definecolor{currentstroke}{rgb}{1.000000,1.000000,1.000000}%
\pgfsetstrokecolor{currentstroke}%
\pgfsetdash{}{0pt}%
\pgfpathmoveto{\pgfqpoint{0.000000in}{0.000000in}}%
\pgfpathlineto{\pgfqpoint{6.400000in}{0.000000in}}%
\pgfpathlineto{\pgfqpoint{6.400000in}{4.800000in}}%
\pgfpathlineto{\pgfqpoint{0.000000in}{4.800000in}}%
\pgfpathclose%
\pgfusepath{fill}%
\end{pgfscope}%
\begin{pgfscope}%
\pgfsetbuttcap%
\pgfsetmiterjoin%
\definecolor{currentfill}{rgb}{1.000000,1.000000,1.000000}%
\pgfsetfillcolor{currentfill}%
\pgfsetlinewidth{0.000000pt}%
\definecolor{currentstroke}{rgb}{0.000000,0.000000,0.000000}%
\pgfsetstrokecolor{currentstroke}%
\pgfsetstrokeopacity{0.000000}%
\pgfsetdash{}{0pt}%
\pgfpathmoveto{\pgfqpoint{0.800000in}{0.528000in}}%
\pgfpathlineto{\pgfqpoint{5.760000in}{0.528000in}}%
\pgfpathlineto{\pgfqpoint{5.760000in}{4.224000in}}%
\pgfpathlineto{\pgfqpoint{0.800000in}{4.224000in}}%
\pgfpathclose%
\pgfusepath{fill}%
\end{pgfscope}%
\begin{pgfscope}%
\pgfsetbuttcap%
\pgfsetroundjoin%
\definecolor{currentfill}{rgb}{0.000000,0.000000,0.000000}%
\pgfsetfillcolor{currentfill}%
\pgfsetlinewidth{0.803000pt}%
\definecolor{currentstroke}{rgb}{0.000000,0.000000,0.000000}%
\pgfsetstrokecolor{currentstroke}%
\pgfsetdash{}{0pt}%
\pgfsys@defobject{currentmarker}{\pgfqpoint{0.000000in}{-0.048611in}}{\pgfqpoint{0.000000in}{0.000000in}}{%
\pgfpathmoveto{\pgfqpoint{0.000000in}{0.000000in}}%
\pgfpathlineto{\pgfqpoint{0.000000in}{-0.048611in}}%
\pgfusepath{stroke,fill}%
}%
\begin{pgfscope}%
\pgfsys@transformshift{1.287291in}{0.528000in}%
\pgfsys@useobject{currentmarker}{}%
\end{pgfscope}%
\end{pgfscope}%
\begin{pgfscope}%
\definecolor{textcolor}{rgb}{0.000000,0.000000,0.000000}%
\pgfsetstrokecolor{textcolor}%
\pgfsetfillcolor{textcolor}%
\pgftext[x=1.287291in,y=0.430778in,,top]{\color{textcolor}\rmfamily\fontsize{14.000000}{16.800000}\selectfont \(\displaystyle 1\)}%
\end{pgfscope}%
\begin{pgfscope}%
\pgfsetbuttcap%
\pgfsetroundjoin%
\definecolor{currentfill}{rgb}{0.000000,0.000000,0.000000}%
\pgfsetfillcolor{currentfill}%
\pgfsetlinewidth{0.803000pt}%
\definecolor{currentstroke}{rgb}{0.000000,0.000000,0.000000}%
\pgfsetstrokecolor{currentstroke}%
\pgfsetdash{}{0pt}%
\pgfsys@defobject{currentmarker}{\pgfqpoint{0.000000in}{-0.048611in}}{\pgfqpoint{0.000000in}{0.000000in}}{%
\pgfpathmoveto{\pgfqpoint{0.000000in}{0.000000in}}%
\pgfpathlineto{\pgfqpoint{0.000000in}{-0.048611in}}%
\pgfusepath{stroke,fill}%
}%
\begin{pgfscope}%
\pgfsys@transformshift{2.149545in}{0.528000in}%
\pgfsys@useobject{currentmarker}{}%
\end{pgfscope}%
\end{pgfscope}%
\begin{pgfscope}%
\definecolor{textcolor}{rgb}{0.000000,0.000000,0.000000}%
\pgfsetstrokecolor{textcolor}%
\pgfsetfillcolor{textcolor}%
\pgftext[x=2.149545in,y=0.430778in,,top]{\color{textcolor}\rmfamily\fontsize{14.000000}{16.800000}\selectfont \(\displaystyle 2\)}%
\end{pgfscope}%
\begin{pgfscope}%
\pgfsetbuttcap%
\pgfsetroundjoin%
\definecolor{currentfill}{rgb}{0.000000,0.000000,0.000000}%
\pgfsetfillcolor{currentfill}%
\pgfsetlinewidth{0.803000pt}%
\definecolor{currentstroke}{rgb}{0.000000,0.000000,0.000000}%
\pgfsetstrokecolor{currentstroke}%
\pgfsetdash{}{0pt}%
\pgfsys@defobject{currentmarker}{\pgfqpoint{0.000000in}{-0.048611in}}{\pgfqpoint{0.000000in}{0.000000in}}{%
\pgfpathmoveto{\pgfqpoint{0.000000in}{0.000000in}}%
\pgfpathlineto{\pgfqpoint{0.000000in}{-0.048611in}}%
\pgfusepath{stroke,fill}%
}%
\begin{pgfscope}%
\pgfsys@transformshift{3.011800in}{0.528000in}%
\pgfsys@useobject{currentmarker}{}%
\end{pgfscope}%
\end{pgfscope}%
\begin{pgfscope}%
\definecolor{textcolor}{rgb}{0.000000,0.000000,0.000000}%
\pgfsetstrokecolor{textcolor}%
\pgfsetfillcolor{textcolor}%
\pgftext[x=3.011800in,y=0.430778in,,top]{\color{textcolor}\rmfamily\fontsize{14.000000}{16.800000}\selectfont \(\displaystyle 3\)}%
\end{pgfscope}%
\begin{pgfscope}%
\pgfsetbuttcap%
\pgfsetroundjoin%
\definecolor{currentfill}{rgb}{0.000000,0.000000,0.000000}%
\pgfsetfillcolor{currentfill}%
\pgfsetlinewidth{0.803000pt}%
\definecolor{currentstroke}{rgb}{0.000000,0.000000,0.000000}%
\pgfsetstrokecolor{currentstroke}%
\pgfsetdash{}{0pt}%
\pgfsys@defobject{currentmarker}{\pgfqpoint{0.000000in}{-0.048611in}}{\pgfqpoint{0.000000in}{0.000000in}}{%
\pgfpathmoveto{\pgfqpoint{0.000000in}{0.000000in}}%
\pgfpathlineto{\pgfqpoint{0.000000in}{-0.048611in}}%
\pgfusepath{stroke,fill}%
}%
\begin{pgfscope}%
\pgfsys@transformshift{3.874054in}{0.528000in}%
\pgfsys@useobject{currentmarker}{}%
\end{pgfscope}%
\end{pgfscope}%
\begin{pgfscope}%
\definecolor{textcolor}{rgb}{0.000000,0.000000,0.000000}%
\pgfsetstrokecolor{textcolor}%
\pgfsetfillcolor{textcolor}%
\pgftext[x=3.874054in,y=0.430778in,,top]{\color{textcolor}\rmfamily\fontsize{14.000000}{16.800000}\selectfont \(\displaystyle 4\)}%
\end{pgfscope}%
\begin{pgfscope}%
\pgfsetbuttcap%
\pgfsetroundjoin%
\definecolor{currentfill}{rgb}{0.000000,0.000000,0.000000}%
\pgfsetfillcolor{currentfill}%
\pgfsetlinewidth{0.803000pt}%
\definecolor{currentstroke}{rgb}{0.000000,0.000000,0.000000}%
\pgfsetstrokecolor{currentstroke}%
\pgfsetdash{}{0pt}%
\pgfsys@defobject{currentmarker}{\pgfqpoint{0.000000in}{-0.048611in}}{\pgfqpoint{0.000000in}{0.000000in}}{%
\pgfpathmoveto{\pgfqpoint{0.000000in}{0.000000in}}%
\pgfpathlineto{\pgfqpoint{0.000000in}{-0.048611in}}%
\pgfusepath{stroke,fill}%
}%
\begin{pgfscope}%
\pgfsys@transformshift{4.736309in}{0.528000in}%
\pgfsys@useobject{currentmarker}{}%
\end{pgfscope}%
\end{pgfscope}%
\begin{pgfscope}%
\definecolor{textcolor}{rgb}{0.000000,0.000000,0.000000}%
\pgfsetstrokecolor{textcolor}%
\pgfsetfillcolor{textcolor}%
\pgftext[x=4.736309in,y=0.430778in,,top]{\color{textcolor}\rmfamily\fontsize{14.000000}{16.800000}\selectfont \(\displaystyle 5\)}%
\end{pgfscope}%
\begin{pgfscope}%
\pgfsetbuttcap%
\pgfsetroundjoin%
\definecolor{currentfill}{rgb}{0.000000,0.000000,0.000000}%
\pgfsetfillcolor{currentfill}%
\pgfsetlinewidth{0.803000pt}%
\definecolor{currentstroke}{rgb}{0.000000,0.000000,0.000000}%
\pgfsetstrokecolor{currentstroke}%
\pgfsetdash{}{0pt}%
\pgfsys@defobject{currentmarker}{\pgfqpoint{0.000000in}{-0.048611in}}{\pgfqpoint{0.000000in}{0.000000in}}{%
\pgfpathmoveto{\pgfqpoint{0.000000in}{0.000000in}}%
\pgfpathlineto{\pgfqpoint{0.000000in}{-0.048611in}}%
\pgfusepath{stroke,fill}%
}%
\begin{pgfscope}%
\pgfsys@transformshift{5.598563in}{0.528000in}%
\pgfsys@useobject{currentmarker}{}%
\end{pgfscope}%
\end{pgfscope}%
\begin{pgfscope}%
\definecolor{textcolor}{rgb}{0.000000,0.000000,0.000000}%
\pgfsetstrokecolor{textcolor}%
\pgfsetfillcolor{textcolor}%
\pgftext[x=5.598563in,y=0.430778in,,top]{\color{textcolor}\rmfamily\fontsize{14.000000}{16.800000}\selectfont \(\displaystyle 6\)}%
\end{pgfscope}%
\begin{pgfscope}%
\definecolor{textcolor}{rgb}{0.000000,0.000000,0.000000}%
\pgfsetstrokecolor{textcolor}%
\pgfsetfillcolor{textcolor}%
\pgftext[x=3.280000in,y=0.187044in,,top]{\color{textcolor}\rmfamily\fontsize{14.000000}{16.800000}\selectfont Speed range \(\displaystyle u\) (m/s)}%
\end{pgfscope}%
\begin{pgfscope}%
\pgfsetbuttcap%
\pgfsetroundjoin%
\definecolor{currentfill}{rgb}{0.000000,0.000000,0.000000}%
\pgfsetfillcolor{currentfill}%
\pgfsetlinewidth{0.803000pt}%
\definecolor{currentstroke}{rgb}{0.000000,0.000000,0.000000}%
\pgfsetstrokecolor{currentstroke}%
\pgfsetdash{}{0pt}%
\pgfsys@defobject{currentmarker}{\pgfqpoint{-0.048611in}{0.000000in}}{\pgfqpoint{0.000000in}{0.000000in}}{%
\pgfpathmoveto{\pgfqpoint{0.000000in}{0.000000in}}%
\pgfpathlineto{\pgfqpoint{-0.048611in}{0.000000in}}%
\pgfusepath{stroke,fill}%
}%
\begin{pgfscope}%
\pgfsys@transformshift{0.800000in}{0.696304in}%
\pgfsys@useobject{currentmarker}{}%
\end{pgfscope}%
\end{pgfscope}%
\begin{pgfscope}%
\definecolor{textcolor}{rgb}{0.000000,0.000000,0.000000}%
\pgfsetstrokecolor{textcolor}%
\pgfsetfillcolor{textcolor}%
\pgftext[x=0.452550in,y=0.622438in,left,base]{\color{textcolor}\rmfamily\fontsize{14.000000}{16.800000}\selectfont \(\displaystyle 0.0\)}%
\end{pgfscope}%
\begin{pgfscope}%
\pgfsetbuttcap%
\pgfsetroundjoin%
\definecolor{currentfill}{rgb}{0.000000,0.000000,0.000000}%
\pgfsetfillcolor{currentfill}%
\pgfsetlinewidth{0.803000pt}%
\definecolor{currentstroke}{rgb}{0.000000,0.000000,0.000000}%
\pgfsetstrokecolor{currentstroke}%
\pgfsetdash{}{0pt}%
\pgfsys@defobject{currentmarker}{\pgfqpoint{-0.048611in}{0.000000in}}{\pgfqpoint{0.000000in}{0.000000in}}{%
\pgfpathmoveto{\pgfqpoint{0.000000in}{0.000000in}}%
\pgfpathlineto{\pgfqpoint{-0.048611in}{0.000000in}}%
\pgfusepath{stroke,fill}%
}%
\begin{pgfscope}%
\pgfsys@transformshift{0.800000in}{1.127464in}%
\pgfsys@useobject{currentmarker}{}%
\end{pgfscope}%
\end{pgfscope}%
\begin{pgfscope}%
\definecolor{textcolor}{rgb}{0.000000,0.000000,0.000000}%
\pgfsetstrokecolor{textcolor}%
\pgfsetfillcolor{textcolor}%
\pgftext[x=0.452550in,y=1.053598in,left,base]{\color{textcolor}\rmfamily\fontsize{14.000000}{16.800000}\selectfont \(\displaystyle 0.2\)}%
\end{pgfscope}%
\begin{pgfscope}%
\pgfsetbuttcap%
\pgfsetroundjoin%
\definecolor{currentfill}{rgb}{0.000000,0.000000,0.000000}%
\pgfsetfillcolor{currentfill}%
\pgfsetlinewidth{0.803000pt}%
\definecolor{currentstroke}{rgb}{0.000000,0.000000,0.000000}%
\pgfsetstrokecolor{currentstroke}%
\pgfsetdash{}{0pt}%
\pgfsys@defobject{currentmarker}{\pgfqpoint{-0.048611in}{0.000000in}}{\pgfqpoint{0.000000in}{0.000000in}}{%
\pgfpathmoveto{\pgfqpoint{0.000000in}{0.000000in}}%
\pgfpathlineto{\pgfqpoint{-0.048611in}{0.000000in}}%
\pgfusepath{stroke,fill}%
}%
\begin{pgfscope}%
\pgfsys@transformshift{0.800000in}{1.558625in}%
\pgfsys@useobject{currentmarker}{}%
\end{pgfscope}%
\end{pgfscope}%
\begin{pgfscope}%
\definecolor{textcolor}{rgb}{0.000000,0.000000,0.000000}%
\pgfsetstrokecolor{textcolor}%
\pgfsetfillcolor{textcolor}%
\pgftext[x=0.452550in,y=1.484759in,left,base]{\color{textcolor}\rmfamily\fontsize{14.000000}{16.800000}\selectfont \(\displaystyle 0.4\)}%
\end{pgfscope}%
\begin{pgfscope}%
\pgfsetbuttcap%
\pgfsetroundjoin%
\definecolor{currentfill}{rgb}{0.000000,0.000000,0.000000}%
\pgfsetfillcolor{currentfill}%
\pgfsetlinewidth{0.803000pt}%
\definecolor{currentstroke}{rgb}{0.000000,0.000000,0.000000}%
\pgfsetstrokecolor{currentstroke}%
\pgfsetdash{}{0pt}%
\pgfsys@defobject{currentmarker}{\pgfqpoint{-0.048611in}{0.000000in}}{\pgfqpoint{0.000000in}{0.000000in}}{%
\pgfpathmoveto{\pgfqpoint{0.000000in}{0.000000in}}%
\pgfpathlineto{\pgfqpoint{-0.048611in}{0.000000in}}%
\pgfusepath{stroke,fill}%
}%
\begin{pgfscope}%
\pgfsys@transformshift{0.800000in}{1.989786in}%
\pgfsys@useobject{currentmarker}{}%
\end{pgfscope}%
\end{pgfscope}%
\begin{pgfscope}%
\definecolor{textcolor}{rgb}{0.000000,0.000000,0.000000}%
\pgfsetstrokecolor{textcolor}%
\pgfsetfillcolor{textcolor}%
\pgftext[x=0.452550in,y=1.915919in,left,base]{\color{textcolor}\rmfamily\fontsize{14.000000}{16.800000}\selectfont \(\displaystyle 0.6\)}%
\end{pgfscope}%
\begin{pgfscope}%
\pgfsetbuttcap%
\pgfsetroundjoin%
\definecolor{currentfill}{rgb}{0.000000,0.000000,0.000000}%
\pgfsetfillcolor{currentfill}%
\pgfsetlinewidth{0.803000pt}%
\definecolor{currentstroke}{rgb}{0.000000,0.000000,0.000000}%
\pgfsetstrokecolor{currentstroke}%
\pgfsetdash{}{0pt}%
\pgfsys@defobject{currentmarker}{\pgfqpoint{-0.048611in}{0.000000in}}{\pgfqpoint{0.000000in}{0.000000in}}{%
\pgfpathmoveto{\pgfqpoint{0.000000in}{0.000000in}}%
\pgfpathlineto{\pgfqpoint{-0.048611in}{0.000000in}}%
\pgfusepath{stroke,fill}%
}%
\begin{pgfscope}%
\pgfsys@transformshift{0.800000in}{2.420946in}%
\pgfsys@useobject{currentmarker}{}%
\end{pgfscope}%
\end{pgfscope}%
\begin{pgfscope}%
\definecolor{textcolor}{rgb}{0.000000,0.000000,0.000000}%
\pgfsetstrokecolor{textcolor}%
\pgfsetfillcolor{textcolor}%
\pgftext[x=0.452550in,y=2.347080in,left,base]{\color{textcolor}\rmfamily\fontsize{14.000000}{16.800000}\selectfont \(\displaystyle 0.8\)}%
\end{pgfscope}%
\begin{pgfscope}%
\pgfsetbuttcap%
\pgfsetroundjoin%
\definecolor{currentfill}{rgb}{0.000000,0.000000,0.000000}%
\pgfsetfillcolor{currentfill}%
\pgfsetlinewidth{0.803000pt}%
\definecolor{currentstroke}{rgb}{0.000000,0.000000,0.000000}%
\pgfsetstrokecolor{currentstroke}%
\pgfsetdash{}{0pt}%
\pgfsys@defobject{currentmarker}{\pgfqpoint{-0.048611in}{0.000000in}}{\pgfqpoint{0.000000in}{0.000000in}}{%
\pgfpathmoveto{\pgfqpoint{0.000000in}{0.000000in}}%
\pgfpathlineto{\pgfqpoint{-0.048611in}{0.000000in}}%
\pgfusepath{stroke,fill}%
}%
\begin{pgfscope}%
\pgfsys@transformshift{0.800000in}{2.852107in}%
\pgfsys@useobject{currentmarker}{}%
\end{pgfscope}%
\end{pgfscope}%
\begin{pgfscope}%
\definecolor{textcolor}{rgb}{0.000000,0.000000,0.000000}%
\pgfsetstrokecolor{textcolor}%
\pgfsetfillcolor{textcolor}%
\pgftext[x=0.452550in,y=2.778241in,left,base]{\color{textcolor}\rmfamily\fontsize{14.000000}{16.800000}\selectfont \(\displaystyle 1.0\)}%
\end{pgfscope}%
\begin{pgfscope}%
\pgfsetbuttcap%
\pgfsetroundjoin%
\definecolor{currentfill}{rgb}{0.000000,0.000000,0.000000}%
\pgfsetfillcolor{currentfill}%
\pgfsetlinewidth{0.803000pt}%
\definecolor{currentstroke}{rgb}{0.000000,0.000000,0.000000}%
\pgfsetstrokecolor{currentstroke}%
\pgfsetdash{}{0pt}%
\pgfsys@defobject{currentmarker}{\pgfqpoint{-0.048611in}{0.000000in}}{\pgfqpoint{0.000000in}{0.000000in}}{%
\pgfpathmoveto{\pgfqpoint{0.000000in}{0.000000in}}%
\pgfpathlineto{\pgfqpoint{-0.048611in}{0.000000in}}%
\pgfusepath{stroke,fill}%
}%
\begin{pgfscope}%
\pgfsys@transformshift{0.800000in}{3.283267in}%
\pgfsys@useobject{currentmarker}{}%
\end{pgfscope}%
\end{pgfscope}%
\begin{pgfscope}%
\definecolor{textcolor}{rgb}{0.000000,0.000000,0.000000}%
\pgfsetstrokecolor{textcolor}%
\pgfsetfillcolor{textcolor}%
\pgftext[x=0.452550in,y=3.209401in,left,base]{\color{textcolor}\rmfamily\fontsize{14.000000}{16.800000}\selectfont \(\displaystyle 1.2\)}%
\end{pgfscope}%
\begin{pgfscope}%
\pgfsetbuttcap%
\pgfsetroundjoin%
\definecolor{currentfill}{rgb}{0.000000,0.000000,0.000000}%
\pgfsetfillcolor{currentfill}%
\pgfsetlinewidth{0.803000pt}%
\definecolor{currentstroke}{rgb}{0.000000,0.000000,0.000000}%
\pgfsetstrokecolor{currentstroke}%
\pgfsetdash{}{0pt}%
\pgfsys@defobject{currentmarker}{\pgfqpoint{-0.048611in}{0.000000in}}{\pgfqpoint{0.000000in}{0.000000in}}{%
\pgfpathmoveto{\pgfqpoint{0.000000in}{0.000000in}}%
\pgfpathlineto{\pgfqpoint{-0.048611in}{0.000000in}}%
\pgfusepath{stroke,fill}%
}%
\begin{pgfscope}%
\pgfsys@transformshift{0.800000in}{3.714428in}%
\pgfsys@useobject{currentmarker}{}%
\end{pgfscope}%
\end{pgfscope}%
\begin{pgfscope}%
\definecolor{textcolor}{rgb}{0.000000,0.000000,0.000000}%
\pgfsetstrokecolor{textcolor}%
\pgfsetfillcolor{textcolor}%
\pgftext[x=0.452550in,y=3.640562in,left,base]{\color{textcolor}\rmfamily\fontsize{14.000000}{16.800000}\selectfont \(\displaystyle 1.4\)}%
\end{pgfscope}%
\begin{pgfscope}%
\pgfsetbuttcap%
\pgfsetroundjoin%
\definecolor{currentfill}{rgb}{0.000000,0.000000,0.000000}%
\pgfsetfillcolor{currentfill}%
\pgfsetlinewidth{0.803000pt}%
\definecolor{currentstroke}{rgb}{0.000000,0.000000,0.000000}%
\pgfsetstrokecolor{currentstroke}%
\pgfsetdash{}{0pt}%
\pgfsys@defobject{currentmarker}{\pgfqpoint{-0.048611in}{0.000000in}}{\pgfqpoint{0.000000in}{0.000000in}}{%
\pgfpathmoveto{\pgfqpoint{0.000000in}{0.000000in}}%
\pgfpathlineto{\pgfqpoint{-0.048611in}{0.000000in}}%
\pgfusepath{stroke,fill}%
}%
\begin{pgfscope}%
\pgfsys@transformshift{0.800000in}{4.145589in}%
\pgfsys@useobject{currentmarker}{}%
\end{pgfscope}%
\end{pgfscope}%
\begin{pgfscope}%
\definecolor{textcolor}{rgb}{0.000000,0.000000,0.000000}%
\pgfsetstrokecolor{textcolor}%
\pgfsetfillcolor{textcolor}%
\pgftext[x=0.452550in,y=4.071723in,left,base]{\color{textcolor}\rmfamily\fontsize{14.000000}{16.800000}\selectfont \(\displaystyle 1.6\)}%
\end{pgfscope}%
\begin{pgfscope}%
\definecolor{textcolor}{rgb}{0.000000,0.000000,0.000000}%
\pgfsetstrokecolor{textcolor}%
\pgfsetfillcolor{textcolor}%
\pgftext[x=0.396994in,y=2.376000in,,bottom,rotate=90.000000]{\color{textcolor}\rmfamily\fontsize{14.000000}{16.800000}\selectfont PDF (s/m)}%
\end{pgfscope}%
\begin{pgfscope}%
\pgfpathrectangle{\pgfqpoint{0.800000in}{0.528000in}}{\pgfqpoint{4.960000in}{3.696000in}}%
\pgfusepath{clip}%
\pgfsetrectcap%
\pgfsetroundjoin%
\pgfsetlinewidth{1.505625pt}%
\definecolor{currentstroke}{rgb}{0.121569,0.466667,0.705882}%
\pgfsetstrokecolor{currentstroke}%
\pgfsetdash{}{0pt}%
\pgfpathmoveto{\pgfqpoint{1.025455in}{0.697777in}}%
\pgfpathlineto{\pgfqpoint{1.056766in}{0.696763in}}%
\pgfpathlineto{\pgfqpoint{1.070682in}{0.698146in}}%
\pgfpathlineto{\pgfqpoint{1.084599in}{0.701836in}}%
\pgfpathlineto{\pgfqpoint{1.101994in}{0.709758in}}%
\pgfpathlineto{\pgfqpoint{1.115910in}{0.715217in}}%
\pgfpathlineto{\pgfqpoint{1.126348in}{0.716909in}}%
\pgfpathlineto{\pgfqpoint{1.147222in}{0.719000in}}%
\pgfpathlineto{\pgfqpoint{1.161138in}{0.723140in}}%
\pgfpathlineto{\pgfqpoint{1.192450in}{0.735120in}}%
\pgfpathlineto{\pgfqpoint{1.202887in}{0.741661in}}%
\pgfpathlineto{\pgfqpoint{1.213324in}{0.752088in}}%
\pgfpathlineto{\pgfqpoint{1.227241in}{0.771003in}}%
\pgfpathlineto{\pgfqpoint{1.241157in}{0.788773in}}%
\pgfpathlineto{\pgfqpoint{1.251594in}{0.798527in}}%
\pgfpathlineto{\pgfqpoint{1.286385in}{0.826937in}}%
\pgfpathlineto{\pgfqpoint{1.296822in}{0.839844in}}%
\pgfpathlineto{\pgfqpoint{1.307259in}{0.856902in}}%
\pgfpathlineto{\pgfqpoint{1.321175in}{0.885384in}}%
\pgfpathlineto{\pgfqpoint{1.335092in}{0.913426in}}%
\pgfpathlineto{\pgfqpoint{1.345529in}{0.928793in}}%
\pgfpathlineto{\pgfqpoint{1.355966in}{0.938607in}}%
\pgfpathlineto{\pgfqpoint{1.373361in}{0.951524in}}%
\pgfpathlineto{\pgfqpoint{1.380320in}{0.959261in}}%
\pgfpathlineto{\pgfqpoint{1.390757in}{0.974979in}}%
\pgfpathlineto{\pgfqpoint{1.418589in}{1.025522in}}%
\pgfpathlineto{\pgfqpoint{1.435985in}{1.060788in}}%
\pgfpathlineto{\pgfqpoint{1.474254in}{1.143800in}}%
\pgfpathlineto{\pgfqpoint{1.491650in}{1.174955in}}%
\pgfpathlineto{\pgfqpoint{1.509045in}{1.206421in}}%
\pgfpathlineto{\pgfqpoint{1.536878in}{1.261250in}}%
\pgfpathlineto{\pgfqpoint{1.543836in}{1.269722in}}%
\pgfpathlineto{\pgfqpoint{1.550794in}{1.274734in}}%
\pgfpathlineto{\pgfqpoint{1.557752in}{1.277123in}}%
\pgfpathlineto{\pgfqpoint{1.578626in}{1.281699in}}%
\pgfpathlineto{\pgfqpoint{1.589064in}{1.287567in}}%
\pgfpathlineto{\pgfqpoint{1.596022in}{1.293861in}}%
\pgfpathlineto{\pgfqpoint{1.602980in}{1.302583in}}%
\pgfpathlineto{\pgfqpoint{1.613417in}{1.320394in}}%
\pgfpathlineto{\pgfqpoint{1.641250in}{1.375073in}}%
\pgfpathlineto{\pgfqpoint{1.648208in}{1.383423in}}%
\pgfpathlineto{\pgfqpoint{1.658645in}{1.391705in}}%
\pgfpathlineto{\pgfqpoint{1.689957in}{1.412789in}}%
\pgfpathlineto{\pgfqpoint{1.700394in}{1.424419in}}%
\pgfpathlineto{\pgfqpoint{1.717789in}{1.449891in}}%
\pgfpathlineto{\pgfqpoint{1.766496in}{1.519761in}}%
\pgfpathlineto{\pgfqpoint{1.773454in}{1.522746in}}%
\pgfpathlineto{\pgfqpoint{1.780412in}{1.522233in}}%
\pgfpathlineto{\pgfqpoint{1.787371in}{1.521002in}}%
\pgfpathlineto{\pgfqpoint{1.790850in}{1.521283in}}%
\pgfpathlineto{\pgfqpoint{1.794329in}{1.522710in}}%
\pgfpathlineto{\pgfqpoint{1.797808in}{1.525550in}}%
\pgfpathlineto{\pgfqpoint{1.804766in}{1.535733in}}%
\pgfpathlineto{\pgfqpoint{1.811724in}{1.550845in}}%
\pgfpathlineto{\pgfqpoint{1.843036in}{1.625781in}}%
\pgfpathlineto{\pgfqpoint{1.849994in}{1.636794in}}%
\pgfpathlineto{\pgfqpoint{1.856952in}{1.644319in}}%
\pgfpathlineto{\pgfqpoint{1.863910in}{1.648602in}}%
\pgfpathlineto{\pgfqpoint{1.874347in}{1.651025in}}%
\pgfpathlineto{\pgfqpoint{1.898701in}{1.652923in}}%
\pgfpathlineto{\pgfqpoint{1.905659in}{1.651327in}}%
\pgfpathlineto{\pgfqpoint{1.926533in}{1.643779in}}%
\pgfpathlineto{\pgfqpoint{1.933491in}{1.644236in}}%
\pgfpathlineto{\pgfqpoint{1.943929in}{1.648361in}}%
\pgfpathlineto{\pgfqpoint{1.957845in}{1.654850in}}%
\pgfpathlineto{\pgfqpoint{1.989157in}{1.665556in}}%
\pgfpathlineto{\pgfqpoint{2.006552in}{1.675975in}}%
\pgfpathlineto{\pgfqpoint{2.016989in}{1.685258in}}%
\pgfpathlineto{\pgfqpoint{2.051780in}{1.722055in}}%
\pgfpathlineto{\pgfqpoint{2.058738in}{1.726286in}}%
\pgfpathlineto{\pgfqpoint{2.065696in}{1.728475in}}%
\pgfpathlineto{\pgfqpoint{2.072654in}{1.728329in}}%
\pgfpathlineto{\pgfqpoint{2.093529in}{1.724278in}}%
\pgfpathlineto{\pgfqpoint{2.100487in}{1.727388in}}%
\pgfpathlineto{\pgfqpoint{2.107445in}{1.734164in}}%
\pgfpathlineto{\pgfqpoint{2.124840in}{1.755656in}}%
\pgfpathlineto{\pgfqpoint{2.131798in}{1.760195in}}%
\pgfpathlineto{\pgfqpoint{2.138756in}{1.760967in}}%
\pgfpathlineto{\pgfqpoint{2.145715in}{1.758870in}}%
\pgfpathlineto{\pgfqpoint{2.159631in}{1.753365in}}%
\pgfpathlineto{\pgfqpoint{2.166589in}{1.753920in}}%
\pgfpathlineto{\pgfqpoint{2.173547in}{1.758497in}}%
\pgfpathlineto{\pgfqpoint{2.180505in}{1.767094in}}%
\pgfpathlineto{\pgfqpoint{2.197901in}{1.793403in}}%
\pgfpathlineto{\pgfqpoint{2.204859in}{1.799744in}}%
\pgfpathlineto{\pgfqpoint{2.211817in}{1.802612in}}%
\pgfpathlineto{\pgfqpoint{2.229212in}{1.804302in}}%
\pgfpathlineto{\pgfqpoint{2.243129in}{1.805023in}}%
\pgfpathlineto{\pgfqpoint{2.253566in}{1.803132in}}%
\pgfpathlineto{\pgfqpoint{2.274440in}{1.797658in}}%
\pgfpathlineto{\pgfqpoint{2.281398in}{1.799085in}}%
\pgfpathlineto{\pgfqpoint{2.288356in}{1.803017in}}%
\pgfpathlineto{\pgfqpoint{2.309231in}{1.817895in}}%
\pgfpathlineto{\pgfqpoint{2.316189in}{1.819026in}}%
\pgfpathlineto{\pgfqpoint{2.330105in}{1.818105in}}%
\pgfpathlineto{\pgfqpoint{2.337063in}{1.820881in}}%
\pgfpathlineto{\pgfqpoint{2.344022in}{1.827678in}}%
\pgfpathlineto{\pgfqpoint{2.361417in}{1.849856in}}%
\pgfpathlineto{\pgfqpoint{2.368375in}{1.854631in}}%
\pgfpathlineto{\pgfqpoint{2.375333in}{1.856166in}}%
\pgfpathlineto{\pgfqpoint{2.385770in}{1.854614in}}%
\pgfpathlineto{\pgfqpoint{2.399687in}{1.851987in}}%
\pgfpathlineto{\pgfqpoint{2.406645in}{1.852978in}}%
\pgfpathlineto{\pgfqpoint{2.413603in}{1.856335in}}%
\pgfpathlineto{\pgfqpoint{2.437956in}{1.872847in}}%
\pgfpathlineto{\pgfqpoint{2.458831in}{1.880036in}}%
\pgfpathlineto{\pgfqpoint{2.469268in}{1.883270in}}%
\pgfpathlineto{\pgfqpoint{2.476226in}{1.883763in}}%
\pgfpathlineto{\pgfqpoint{2.483184in}{1.882436in}}%
\pgfpathlineto{\pgfqpoint{2.507538in}{1.874798in}}%
\pgfpathlineto{\pgfqpoint{2.524933in}{1.873088in}}%
\pgfpathlineto{\pgfqpoint{2.531891in}{1.874069in}}%
\pgfpathlineto{\pgfqpoint{2.538849in}{1.876965in}}%
\pgfpathlineto{\pgfqpoint{2.556245in}{1.886150in}}%
\pgfpathlineto{\pgfqpoint{2.563203in}{1.887297in}}%
\pgfpathlineto{\pgfqpoint{2.573640in}{1.886274in}}%
\pgfpathlineto{\pgfqpoint{2.587556in}{1.882164in}}%
\pgfpathlineto{\pgfqpoint{2.611910in}{1.873978in}}%
\pgfpathlineto{\pgfqpoint{2.622347in}{1.873540in}}%
\pgfpathlineto{\pgfqpoint{2.632784in}{1.876079in}}%
\pgfpathlineto{\pgfqpoint{2.643221in}{1.878838in}}%
\pgfpathlineto{\pgfqpoint{2.650180in}{1.878905in}}%
\pgfpathlineto{\pgfqpoint{2.657138in}{1.876568in}}%
\pgfpathlineto{\pgfqpoint{2.664096in}{1.871531in}}%
\pgfpathlineto{\pgfqpoint{2.674533in}{1.859908in}}%
\pgfpathlineto{\pgfqpoint{2.695407in}{1.834869in}}%
\pgfpathlineto{\pgfqpoint{2.705845in}{1.826496in}}%
\pgfpathlineto{\pgfqpoint{2.712803in}{1.823186in}}%
\pgfpathlineto{\pgfqpoint{2.719761in}{1.821656in}}%
\pgfpathlineto{\pgfqpoint{2.730198in}{1.821904in}}%
\pgfpathlineto{\pgfqpoint{2.740635in}{1.822035in}}%
\pgfpathlineto{\pgfqpoint{2.747593in}{1.819436in}}%
\pgfpathlineto{\pgfqpoint{2.754552in}{1.813477in}}%
\pgfpathlineto{\pgfqpoint{2.761510in}{1.804262in}}%
\pgfpathlineto{\pgfqpoint{2.775426in}{1.779940in}}%
\pgfpathlineto{\pgfqpoint{2.785863in}{1.762938in}}%
\pgfpathlineto{\pgfqpoint{2.792821in}{1.754958in}}%
\pgfpathlineto{\pgfqpoint{2.799780in}{1.750030in}}%
\pgfpathlineto{\pgfqpoint{2.810217in}{1.746658in}}%
\pgfpathlineto{\pgfqpoint{2.824133in}{1.742928in}}%
\pgfpathlineto{\pgfqpoint{2.834570in}{1.737728in}}%
\pgfpathlineto{\pgfqpoint{2.845007in}{1.730196in}}%
\pgfpathlineto{\pgfqpoint{2.872840in}{1.707062in}}%
\pgfpathlineto{\pgfqpoint{2.886756in}{1.701231in}}%
\pgfpathlineto{\pgfqpoint{2.893714in}{1.697789in}}%
\pgfpathlineto{\pgfqpoint{2.900673in}{1.692210in}}%
\pgfpathlineto{\pgfqpoint{2.911110in}{1.679551in}}%
\pgfpathlineto{\pgfqpoint{2.925026in}{1.662339in}}%
\pgfpathlineto{\pgfqpoint{2.931984in}{1.657441in}}%
\pgfpathlineto{\pgfqpoint{2.938942in}{1.655693in}}%
\pgfpathlineto{\pgfqpoint{2.952859in}{1.655353in}}%
\pgfpathlineto{\pgfqpoint{2.956338in}{1.653767in}}%
\pgfpathlineto{\pgfqpoint{2.959817in}{1.650870in}}%
\pgfpathlineto{\pgfqpoint{2.966775in}{1.640657in}}%
\pgfpathlineto{\pgfqpoint{2.973733in}{1.625417in}}%
\pgfpathlineto{\pgfqpoint{2.994607in}{1.575988in}}%
\pgfpathlineto{\pgfqpoint{3.005045in}{1.558606in}}%
\pgfpathlineto{\pgfqpoint{3.015482in}{1.545953in}}%
\pgfpathlineto{\pgfqpoint{3.025919in}{1.536740in}}%
\pgfpathlineto{\pgfqpoint{3.050272in}{1.517667in}}%
\pgfpathlineto{\pgfqpoint{3.060710in}{1.505782in}}%
\pgfpathlineto{\pgfqpoint{3.078105in}{1.480016in}}%
\pgfpathlineto{\pgfqpoint{3.088542in}{1.466045in}}%
\pgfpathlineto{\pgfqpoint{3.123333in}{1.425714in}}%
\pgfpathlineto{\pgfqpoint{3.137249in}{1.402688in}}%
\pgfpathlineto{\pgfqpoint{3.189435in}{1.306831in}}%
\pgfpathlineto{\pgfqpoint{3.213789in}{1.255191in}}%
\pgfpathlineto{\pgfqpoint{3.227705in}{1.231897in}}%
\pgfpathlineto{\pgfqpoint{3.283370in}{1.144273in}}%
\pgfpathlineto{\pgfqpoint{3.311203in}{1.093451in}}%
\pgfpathlineto{\pgfqpoint{3.318161in}{1.085335in}}%
\pgfpathlineto{\pgfqpoint{3.325119in}{1.080035in}}%
\pgfpathlineto{\pgfqpoint{3.335556in}{1.075944in}}%
\pgfpathlineto{\pgfqpoint{3.352951in}{1.071138in}}%
\pgfpathlineto{\pgfqpoint{3.359910in}{1.067018in}}%
\pgfpathlineto{\pgfqpoint{3.366868in}{1.060619in}}%
\pgfpathlineto{\pgfqpoint{3.377305in}{1.047242in}}%
\pgfpathlineto{\pgfqpoint{3.398179in}{1.018543in}}%
\pgfpathlineto{\pgfqpoint{3.415575in}{1.000439in}}%
\pgfpathlineto{\pgfqpoint{3.429491in}{0.984472in}}%
\pgfpathlineto{\pgfqpoint{3.457323in}{0.950117in}}%
\pgfpathlineto{\pgfqpoint{3.474719in}{0.933178in}}%
\pgfpathlineto{\pgfqpoint{3.499072in}{0.909993in}}%
\pgfpathlineto{\pgfqpoint{3.512989in}{0.892430in}}%
\pgfpathlineto{\pgfqpoint{3.533863in}{0.865083in}}%
\pgfpathlineto{\pgfqpoint{3.544300in}{0.855231in}}%
\pgfpathlineto{\pgfqpoint{3.554737in}{0.849160in}}%
\pgfpathlineto{\pgfqpoint{3.572133in}{0.843053in}}%
\pgfpathlineto{\pgfqpoint{3.593007in}{0.835749in}}%
\pgfpathlineto{\pgfqpoint{3.603444in}{0.829676in}}%
\pgfpathlineto{\pgfqpoint{3.617361in}{0.818749in}}%
\pgfpathlineto{\pgfqpoint{3.634756in}{0.804881in}}%
\pgfpathlineto{\pgfqpoint{3.645193in}{0.799062in}}%
\pgfpathlineto{\pgfqpoint{3.662589in}{0.792660in}}%
\pgfpathlineto{\pgfqpoint{3.679984in}{0.785520in}}%
\pgfpathlineto{\pgfqpoint{3.697379in}{0.775552in}}%
\pgfpathlineto{\pgfqpoint{3.718254in}{0.763692in}}%
\pgfpathlineto{\pgfqpoint{3.728691in}{0.760223in}}%
\pgfpathlineto{\pgfqpoint{3.746086in}{0.757916in}}%
\pgfpathlineto{\pgfqpoint{3.756523in}{0.755339in}}%
\pgfpathlineto{\pgfqpoint{3.766961in}{0.749961in}}%
\pgfpathlineto{\pgfqpoint{3.787835in}{0.737058in}}%
\pgfpathlineto{\pgfqpoint{3.798272in}{0.733421in}}%
\pgfpathlineto{\pgfqpoint{3.815668in}{0.730825in}}%
\pgfpathlineto{\pgfqpoint{3.836542in}{0.726917in}}%
\pgfpathlineto{\pgfqpoint{3.860895in}{0.722675in}}%
\pgfpathlineto{\pgfqpoint{3.906123in}{0.717400in}}%
\pgfpathlineto{\pgfqpoint{3.951351in}{0.708227in}}%
\pgfpathlineto{\pgfqpoint{3.979184in}{0.707789in}}%
\pgfpathlineto{\pgfqpoint{4.013974in}{0.706420in}}%
\pgfpathlineto{\pgfqpoint{4.097472in}{0.702013in}}%
\pgfpathlineto{\pgfqpoint{4.121826in}{0.702110in}}%
\pgfpathlineto{\pgfqpoint{4.215760in}{0.697204in}}%
\pgfpathlineto{\pgfqpoint{4.497565in}{0.696433in}}%
\pgfpathlineto{\pgfqpoint{4.497565in}{0.696433in}}%
\pgfusepath{stroke}%
\end{pgfscope}%
\begin{pgfscope}%
\pgfpathrectangle{\pgfqpoint{0.800000in}{0.528000in}}{\pgfqpoint{4.960000in}{3.696000in}}%
\pgfusepath{clip}%
\pgfsetrectcap%
\pgfsetroundjoin%
\pgfsetlinewidth{1.505625pt}%
\definecolor{currentstroke}{rgb}{1.000000,0.498039,0.054902}%
\pgfsetstrokecolor{currentstroke}%
\pgfsetdash{}{0pt}%
\pgfpathmoveto{\pgfqpoint{2.398000in}{0.703153in}}%
\pgfpathlineto{\pgfqpoint{2.413262in}{0.697715in}}%
\pgfpathlineto{\pgfqpoint{2.425471in}{0.696085in}}%
\pgfpathlineto{\pgfqpoint{2.449889in}{0.696515in}}%
\pgfpathlineto{\pgfqpoint{2.495673in}{0.696448in}}%
\pgfpathlineto{\pgfqpoint{2.590293in}{0.696974in}}%
\pgfpathlineto{\pgfqpoint{2.642181in}{0.698332in}}%
\pgfpathlineto{\pgfqpoint{2.684913in}{0.699556in}}%
\pgfpathlineto{\pgfqpoint{2.727644in}{0.699728in}}%
\pgfpathlineto{\pgfqpoint{2.761219in}{0.699357in}}%
\pgfpathlineto{\pgfqpoint{2.782585in}{0.702045in}}%
\pgfpathlineto{\pgfqpoint{2.816160in}{0.706172in}}%
\pgfpathlineto{\pgfqpoint{2.852787in}{0.710427in}}%
\pgfpathlineto{\pgfqpoint{2.941303in}{0.725534in}}%
\pgfpathlineto{\pgfqpoint{2.965721in}{0.732418in}}%
\pgfpathlineto{\pgfqpoint{2.977930in}{0.738556in}}%
\pgfpathlineto{\pgfqpoint{2.999296in}{0.751204in}}%
\pgfpathlineto{\pgfqpoint{3.014557in}{0.756124in}}%
\pgfpathlineto{\pgfqpoint{3.032870in}{0.762339in}}%
\pgfpathlineto{\pgfqpoint{3.048132in}{0.770304in}}%
\pgfpathlineto{\pgfqpoint{3.063393in}{0.778567in}}%
\pgfpathlineto{\pgfqpoint{3.075602in}{0.782418in}}%
\pgfpathlineto{\pgfqpoint{3.090863in}{0.787120in}}%
\pgfpathlineto{\pgfqpoint{3.103072in}{0.793931in}}%
\pgfpathlineto{\pgfqpoint{3.115282in}{0.804129in}}%
\pgfpathlineto{\pgfqpoint{3.127491in}{0.818015in}}%
\pgfpathlineto{\pgfqpoint{3.154961in}{0.851631in}}%
\pgfpathlineto{\pgfqpoint{3.176327in}{0.874576in}}%
\pgfpathlineto{\pgfqpoint{3.209902in}{0.917541in}}%
\pgfpathlineto{\pgfqpoint{3.222111in}{0.928679in}}%
\pgfpathlineto{\pgfqpoint{3.234320in}{0.936687in}}%
\pgfpathlineto{\pgfqpoint{3.246529in}{0.941455in}}%
\pgfpathlineto{\pgfqpoint{3.258738in}{0.945467in}}%
\pgfpathlineto{\pgfqpoint{3.264842in}{0.949405in}}%
\pgfpathlineto{\pgfqpoint{3.270947in}{0.956171in}}%
\pgfpathlineto{\pgfqpoint{3.277051in}{0.966442in}}%
\pgfpathlineto{\pgfqpoint{3.286208in}{0.988057in}}%
\pgfpathlineto{\pgfqpoint{3.313678in}{1.061221in}}%
\pgfpathlineto{\pgfqpoint{3.319783in}{1.070934in}}%
\pgfpathlineto{\pgfqpoint{3.328940in}{1.080470in}}%
\pgfpathlineto{\pgfqpoint{3.344201in}{1.093137in}}%
\pgfpathlineto{\pgfqpoint{3.350306in}{1.100903in}}%
\pgfpathlineto{\pgfqpoint{3.359462in}{1.117399in}}%
\pgfpathlineto{\pgfqpoint{3.399142in}{1.205998in}}%
\pgfpathlineto{\pgfqpoint{3.411351in}{1.239908in}}%
\pgfpathlineto{\pgfqpoint{3.429664in}{1.293209in}}%
\pgfpathlineto{\pgfqpoint{3.441873in}{1.321206in}}%
\pgfpathlineto{\pgfqpoint{3.463239in}{1.367148in}}%
\pgfpathlineto{\pgfqpoint{3.475448in}{1.400500in}}%
\pgfpathlineto{\pgfqpoint{3.502919in}{1.479146in}}%
\pgfpathlineto{\pgfqpoint{3.515128in}{1.505553in}}%
\pgfpathlineto{\pgfqpoint{3.533441in}{1.543670in}}%
\pgfpathlineto{\pgfqpoint{3.548703in}{1.581060in}}%
\pgfpathlineto{\pgfqpoint{3.557859in}{1.608841in}}%
\pgfpathlineto{\pgfqpoint{3.573121in}{1.665434in}}%
\pgfpathlineto{\pgfqpoint{3.594486in}{1.744199in}}%
\pgfpathlineto{\pgfqpoint{3.606696in}{1.781793in}}%
\pgfpathlineto{\pgfqpoint{3.621957in}{1.819737in}}%
\pgfpathlineto{\pgfqpoint{3.643323in}{1.865278in}}%
\pgfpathlineto{\pgfqpoint{3.689107in}{1.965440in}}%
\pgfpathlineto{\pgfqpoint{3.695211in}{1.987079in}}%
\pgfpathlineto{\pgfqpoint{3.704368in}{2.028236in}}%
\pgfpathlineto{\pgfqpoint{3.719629in}{2.098131in}}%
\pgfpathlineto{\pgfqpoint{3.725734in}{2.117794in}}%
\pgfpathlineto{\pgfqpoint{3.731838in}{2.131854in}}%
\pgfpathlineto{\pgfqpoint{3.753204in}{2.173057in}}%
\pgfpathlineto{\pgfqpoint{3.762361in}{2.201202in}}%
\pgfpathlineto{\pgfqpoint{3.780674in}{2.262095in}}%
\pgfpathlineto{\pgfqpoint{3.792883in}{2.292827in}}%
\pgfpathlineto{\pgfqpoint{3.808145in}{2.331885in}}%
\pgfpathlineto{\pgfqpoint{3.820354in}{2.362911in}}%
\pgfpathlineto{\pgfqpoint{3.826458in}{2.374484in}}%
\pgfpathlineto{\pgfqpoint{3.832563in}{2.382451in}}%
\pgfpathlineto{\pgfqpoint{3.853929in}{2.403899in}}%
\pgfpathlineto{\pgfqpoint{3.860033in}{2.417944in}}%
\pgfpathlineto{\pgfqpoint{3.866138in}{2.438566in}}%
\pgfpathlineto{\pgfqpoint{3.887504in}{2.523201in}}%
\pgfpathlineto{\pgfqpoint{3.893608in}{2.535001in}}%
\pgfpathlineto{\pgfqpoint{3.896660in}{2.537527in}}%
\pgfpathlineto{\pgfqpoint{3.899713in}{2.538115in}}%
\pgfpathlineto{\pgfqpoint{3.905817in}{2.535782in}}%
\pgfpathlineto{\pgfqpoint{3.911922in}{2.533573in}}%
\pgfpathlineto{\pgfqpoint{3.914974in}{2.534070in}}%
\pgfpathlineto{\pgfqpoint{3.918026in}{2.536093in}}%
\pgfpathlineto{\pgfqpoint{3.924131in}{2.544830in}}%
\pgfpathlineto{\pgfqpoint{3.930235in}{2.558548in}}%
\pgfpathlineto{\pgfqpoint{3.954653in}{2.621463in}}%
\pgfpathlineto{\pgfqpoint{3.963810in}{2.637191in}}%
\pgfpathlineto{\pgfqpoint{3.969915in}{2.643963in}}%
\pgfpathlineto{\pgfqpoint{3.976019in}{2.646508in}}%
\pgfpathlineto{\pgfqpoint{3.979071in}{2.645958in}}%
\pgfpathlineto{\pgfqpoint{3.985176in}{2.641679in}}%
\pgfpathlineto{\pgfqpoint{3.994333in}{2.632654in}}%
\pgfpathlineto{\pgfqpoint{4.000437in}{2.630375in}}%
\pgfpathlineto{\pgfqpoint{4.003489in}{2.631294in}}%
\pgfpathlineto{\pgfqpoint{4.009594in}{2.637051in}}%
\pgfpathlineto{\pgfqpoint{4.018751in}{2.651888in}}%
\pgfpathlineto{\pgfqpoint{4.049273in}{2.707406in}}%
\pgfpathlineto{\pgfqpoint{4.055378in}{2.714056in}}%
\pgfpathlineto{\pgfqpoint{4.061482in}{2.716968in}}%
\pgfpathlineto{\pgfqpoint{4.067587in}{2.716210in}}%
\pgfpathlineto{\pgfqpoint{4.076744in}{2.710709in}}%
\pgfpathlineto{\pgfqpoint{4.092005in}{2.700412in}}%
\pgfpathlineto{\pgfqpoint{4.110319in}{2.692040in}}%
\pgfpathlineto{\pgfqpoint{4.116423in}{2.685447in}}%
\pgfpathlineto{\pgfqpoint{4.125580in}{2.669742in}}%
\pgfpathlineto{\pgfqpoint{4.137789in}{2.648037in}}%
\pgfpathlineto{\pgfqpoint{4.143893in}{2.640754in}}%
\pgfpathlineto{\pgfqpoint{4.149998in}{2.636858in}}%
\pgfpathlineto{\pgfqpoint{4.156102in}{2.635765in}}%
\pgfpathlineto{\pgfqpoint{4.165259in}{2.634871in}}%
\pgfpathlineto{\pgfqpoint{4.171364in}{2.630653in}}%
\pgfpathlineto{\pgfqpoint{4.177468in}{2.621148in}}%
\pgfpathlineto{\pgfqpoint{4.186625in}{2.598808in}}%
\pgfpathlineto{\pgfqpoint{4.198834in}{2.568412in}}%
\pgfpathlineto{\pgfqpoint{4.217148in}{2.530439in}}%
\pgfpathlineto{\pgfqpoint{4.226304in}{2.503078in}}%
\pgfpathlineto{\pgfqpoint{4.244618in}{2.445333in}}%
\pgfpathlineto{\pgfqpoint{4.256827in}{2.413501in}}%
\pgfpathlineto{\pgfqpoint{4.281245in}{2.357433in}}%
\pgfpathlineto{\pgfqpoint{4.311768in}{2.262567in}}%
\pgfpathlineto{\pgfqpoint{4.320925in}{2.245935in}}%
\pgfpathlineto{\pgfqpoint{4.333134in}{2.225246in}}%
\pgfpathlineto{\pgfqpoint{4.339238in}{2.210412in}}%
\pgfpathlineto{\pgfqpoint{4.345343in}{2.190478in}}%
\pgfpathlineto{\pgfqpoint{4.357552in}{2.139621in}}%
\pgfpathlineto{\pgfqpoint{4.369761in}{2.089420in}}%
\pgfpathlineto{\pgfqpoint{4.385022in}{2.038723in}}%
\pgfpathlineto{\pgfqpoint{4.397231in}{1.994724in}}%
\pgfpathlineto{\pgfqpoint{4.443015in}{1.812981in}}%
\pgfpathlineto{\pgfqpoint{4.528478in}{1.509066in}}%
\pgfpathlineto{\pgfqpoint{4.546792in}{1.429704in}}%
\pgfpathlineto{\pgfqpoint{4.562053in}{1.366019in}}%
\pgfpathlineto{\pgfqpoint{4.571210in}{1.335869in}}%
\pgfpathlineto{\pgfqpoint{4.580367in}{1.312319in}}%
\pgfpathlineto{\pgfqpoint{4.592576in}{1.288230in}}%
\pgfpathlineto{\pgfqpoint{4.616994in}{1.244790in}}%
\pgfpathlineto{\pgfqpoint{4.629203in}{1.216888in}}%
\pgfpathlineto{\pgfqpoint{4.644464in}{1.175219in}}%
\pgfpathlineto{\pgfqpoint{4.687196in}{1.053847in}}%
\pgfpathlineto{\pgfqpoint{4.696353in}{1.037391in}}%
\pgfpathlineto{\pgfqpoint{4.708562in}{1.022177in}}%
\pgfpathlineto{\pgfqpoint{4.717718in}{1.010347in}}%
\pgfpathlineto{\pgfqpoint{4.726875in}{0.994538in}}%
\pgfpathlineto{\pgfqpoint{4.736032in}{0.974219in}}%
\pgfpathlineto{\pgfqpoint{4.754346in}{0.926274in}}%
\pgfpathlineto{\pgfqpoint{4.766555in}{0.895583in}}%
\pgfpathlineto{\pgfqpoint{4.775711in}{0.878329in}}%
\pgfpathlineto{\pgfqpoint{4.784868in}{0.867048in}}%
\pgfpathlineto{\pgfqpoint{4.809286in}{0.843041in}}%
\pgfpathlineto{\pgfqpoint{4.830652in}{0.815276in}}%
\pgfpathlineto{\pgfqpoint{4.839809in}{0.807294in}}%
\pgfpathlineto{\pgfqpoint{4.861175in}{0.793280in}}%
\pgfpathlineto{\pgfqpoint{4.891697in}{0.772850in}}%
\pgfpathlineto{\pgfqpoint{4.946638in}{0.734857in}}%
\pgfpathlineto{\pgfqpoint{4.961899in}{0.730236in}}%
\pgfpathlineto{\pgfqpoint{4.986317in}{0.723201in}}%
\pgfpathlineto{\pgfqpoint{5.001579in}{0.716522in}}%
\pgfpathlineto{\pgfqpoint{5.016840in}{0.709838in}}%
\pgfpathlineto{\pgfqpoint{5.029049in}{0.707350in}}%
\pgfpathlineto{\pgfqpoint{5.090094in}{0.702416in}}%
\pgfpathlineto{\pgfqpoint{5.114512in}{0.700861in}}%
\pgfpathlineto{\pgfqpoint{5.282387in}{0.696520in}}%
\pgfpathlineto{\pgfqpoint{5.444157in}{0.698142in}}%
\pgfpathlineto{\pgfqpoint{5.444157in}{0.698142in}}%
\pgfusepath{stroke}%
\end{pgfscope}%
\begin{pgfscope}%
\pgfpathrectangle{\pgfqpoint{0.800000in}{0.528000in}}{\pgfqpoint{4.960000in}{3.696000in}}%
\pgfusepath{clip}%
\pgfsetrectcap%
\pgfsetroundjoin%
\pgfsetlinewidth{1.505625pt}%
\definecolor{currentstroke}{rgb}{0.172549,0.627451,0.172549}%
\pgfsetstrokecolor{currentstroke}%
\pgfsetdash{}{0pt}%
\pgfpathmoveto{\pgfqpoint{3.496045in}{0.698606in}}%
\pgfpathlineto{\pgfqpoint{3.510343in}{0.696490in}}%
\pgfpathlineto{\pgfqpoint{3.532812in}{0.696670in}}%
\pgfpathlineto{\pgfqpoint{3.626770in}{0.698113in}}%
\pgfpathlineto{\pgfqpoint{3.663537in}{0.698508in}}%
\pgfpathlineto{\pgfqpoint{3.730942in}{0.702609in}}%
\pgfpathlineto{\pgfqpoint{3.777922in}{0.705902in}}%
\pgfpathlineto{\pgfqpoint{3.806518in}{0.711736in}}%
\pgfpathlineto{\pgfqpoint{3.833072in}{0.712799in}}%
\pgfpathlineto{\pgfqpoint{3.857583in}{0.719167in}}%
\pgfpathlineto{\pgfqpoint{3.875966in}{0.719645in}}%
\pgfpathlineto{\pgfqpoint{3.916818in}{0.727609in}}%
\pgfpathlineto{\pgfqpoint{3.931116in}{0.735770in}}%
\pgfpathlineto{\pgfqpoint{3.941329in}{0.741749in}}%
\pgfpathlineto{\pgfqpoint{3.947456in}{0.742882in}}%
\pgfpathlineto{\pgfqpoint{3.963797in}{0.742423in}}%
\pgfpathlineto{\pgfqpoint{3.969925in}{0.745900in}}%
\pgfpathlineto{\pgfqpoint{3.982180in}{0.753819in}}%
\pgfpathlineto{\pgfqpoint{3.990351in}{0.756235in}}%
\pgfpathlineto{\pgfqpoint{3.998521in}{0.758609in}}%
\pgfpathlineto{\pgfqpoint{4.004649in}{0.762755in}}%
\pgfpathlineto{\pgfqpoint{4.014862in}{0.773593in}}%
\pgfpathlineto{\pgfqpoint{4.023032in}{0.781562in}}%
\pgfpathlineto{\pgfqpoint{4.031202in}{0.786316in}}%
\pgfpathlineto{\pgfqpoint{4.043458in}{0.792345in}}%
\pgfpathlineto{\pgfqpoint{4.051628in}{0.799346in}}%
\pgfpathlineto{\pgfqpoint{4.076139in}{0.823499in}}%
\pgfpathlineto{\pgfqpoint{4.088395in}{0.833563in}}%
\pgfpathlineto{\pgfqpoint{4.096565in}{0.843940in}}%
\pgfpathlineto{\pgfqpoint{4.112906in}{0.868230in}}%
\pgfpathlineto{\pgfqpoint{4.119034in}{0.872406in}}%
\pgfpathlineto{\pgfqpoint{4.127204in}{0.874372in}}%
\pgfpathlineto{\pgfqpoint{4.135374in}{0.876219in}}%
\pgfpathlineto{\pgfqpoint{4.139459in}{0.878782in}}%
\pgfpathlineto{\pgfqpoint{4.143545in}{0.883688in}}%
\pgfpathlineto{\pgfqpoint{4.147630in}{0.891528in}}%
\pgfpathlineto{\pgfqpoint{4.153758in}{0.908347in}}%
\pgfpathlineto{\pgfqpoint{4.166013in}{0.945040in}}%
\pgfpathlineto{\pgfqpoint{4.172141in}{0.956121in}}%
\pgfpathlineto{\pgfqpoint{4.178269in}{0.962023in}}%
\pgfpathlineto{\pgfqpoint{4.188482in}{0.970017in}}%
\pgfpathlineto{\pgfqpoint{4.194609in}{0.977361in}}%
\pgfpathlineto{\pgfqpoint{4.204822in}{0.994311in}}%
\pgfpathlineto{\pgfqpoint{4.217078in}{1.015956in}}%
\pgfpathlineto{\pgfqpoint{4.223205in}{1.032765in}}%
\pgfpathlineto{\pgfqpoint{4.231376in}{1.063897in}}%
\pgfpathlineto{\pgfqpoint{4.243631in}{1.109624in}}%
\pgfpathlineto{\pgfqpoint{4.253844in}{1.136675in}}%
\pgfpathlineto{\pgfqpoint{4.264057in}{1.165130in}}%
\pgfpathlineto{\pgfqpoint{4.276313in}{1.200707in}}%
\pgfpathlineto{\pgfqpoint{4.298781in}{1.252757in}}%
\pgfpathlineto{\pgfqpoint{4.306952in}{1.284632in}}%
\pgfpathlineto{\pgfqpoint{4.317164in}{1.324798in}}%
\pgfpathlineto{\pgfqpoint{4.347803in}{1.424138in}}%
\pgfpathlineto{\pgfqpoint{4.360059in}{1.473338in}}%
\pgfpathlineto{\pgfqpoint{4.376399in}{1.538506in}}%
\pgfpathlineto{\pgfqpoint{4.388655in}{1.579292in}}%
\pgfpathlineto{\pgfqpoint{4.398868in}{1.613110in}}%
\pgfpathlineto{\pgfqpoint{4.404996in}{1.642322in}}%
\pgfpathlineto{\pgfqpoint{4.411123in}{1.684354in}}%
\pgfpathlineto{\pgfqpoint{4.433592in}{1.864715in}}%
\pgfpathlineto{\pgfqpoint{4.437677in}{1.881076in}}%
\pgfpathlineto{\pgfqpoint{4.441762in}{1.889883in}}%
\pgfpathlineto{\pgfqpoint{4.451975in}{1.901455in}}%
\pgfpathlineto{\pgfqpoint{4.456060in}{1.913994in}}%
\pgfpathlineto{\pgfqpoint{4.460145in}{1.935020in}}%
\pgfpathlineto{\pgfqpoint{4.468316in}{1.994986in}}%
\pgfpathlineto{\pgfqpoint{4.476486in}{2.052795in}}%
\pgfpathlineto{\pgfqpoint{4.484656in}{2.092724in}}%
\pgfpathlineto{\pgfqpoint{4.492827in}{2.132609in}}%
\pgfpathlineto{\pgfqpoint{4.500997in}{2.191706in}}%
\pgfpathlineto{\pgfqpoint{4.515295in}{2.299601in}}%
\pgfpathlineto{\pgfqpoint{4.543891in}{2.480088in}}%
\pgfpathlineto{\pgfqpoint{4.568402in}{2.588985in}}%
\pgfpathlineto{\pgfqpoint{4.578615in}{2.661125in}}%
\pgfpathlineto{\pgfqpoint{4.625595in}{3.033505in}}%
\pgfpathlineto{\pgfqpoint{4.631723in}{3.083166in}}%
\pgfpathlineto{\pgfqpoint{4.635808in}{3.105864in}}%
\pgfpathlineto{\pgfqpoint{4.639893in}{3.119667in}}%
\pgfpathlineto{\pgfqpoint{4.648063in}{3.139086in}}%
\pgfpathlineto{\pgfqpoint{4.652148in}{3.156872in}}%
\pgfpathlineto{\pgfqpoint{4.658276in}{3.200687in}}%
\pgfpathlineto{\pgfqpoint{4.682787in}{3.410329in}}%
\pgfpathlineto{\pgfqpoint{4.690958in}{3.458049in}}%
\pgfpathlineto{\pgfqpoint{4.699128in}{3.490976in}}%
\pgfpathlineto{\pgfqpoint{4.711383in}{3.533668in}}%
\pgfpathlineto{\pgfqpoint{4.717511in}{3.568958in}}%
\pgfpathlineto{\pgfqpoint{4.735894in}{3.697257in}}%
\pgfpathlineto{\pgfqpoint{4.739980in}{3.712334in}}%
\pgfpathlineto{\pgfqpoint{4.744065in}{3.721304in}}%
\pgfpathlineto{\pgfqpoint{4.750193in}{3.728612in}}%
\pgfpathlineto{\pgfqpoint{4.754278in}{3.733866in}}%
\pgfpathlineto{\pgfqpoint{4.758363in}{3.742964in}}%
\pgfpathlineto{\pgfqpoint{4.762448in}{3.757480in}}%
\pgfpathlineto{\pgfqpoint{4.768576in}{3.788077in}}%
\pgfpathlineto{\pgfqpoint{4.799215in}{3.959361in}}%
\pgfpathlineto{\pgfqpoint{4.817598in}{4.032965in}}%
\pgfpathlineto{\pgfqpoint{4.821683in}{4.043251in}}%
\pgfpathlineto{\pgfqpoint{4.825768in}{4.048983in}}%
\pgfpathlineto{\pgfqpoint{4.831896in}{4.052583in}}%
\pgfpathlineto{\pgfqpoint{4.842109in}{4.056000in}}%
\pgfpathlineto{\pgfqpoint{4.846194in}{4.053484in}}%
\pgfpathlineto{\pgfqpoint{4.850279in}{4.045703in}}%
\pgfpathlineto{\pgfqpoint{4.854364in}{4.031765in}}%
\pgfpathlineto{\pgfqpoint{4.860492in}{4.000761in}}%
\pgfpathlineto{\pgfqpoint{4.878875in}{3.898799in}}%
\pgfpathlineto{\pgfqpoint{4.889088in}{3.858630in}}%
\pgfpathlineto{\pgfqpoint{4.895216in}{3.841949in}}%
\pgfpathlineto{\pgfqpoint{4.913599in}{3.799412in}}%
\pgfpathlineto{\pgfqpoint{4.919727in}{3.770592in}}%
\pgfpathlineto{\pgfqpoint{4.925855in}{3.727967in}}%
\pgfpathlineto{\pgfqpoint{4.934025in}{3.653419in}}%
\pgfpathlineto{\pgfqpoint{4.946281in}{3.519781in}}%
\pgfpathlineto{\pgfqpoint{4.968749in}{3.259763in}}%
\pgfpathlineto{\pgfqpoint{4.987133in}{3.084424in}}%
\pgfpathlineto{\pgfqpoint{4.997345in}{2.948324in}}%
\pgfpathlineto{\pgfqpoint{5.007558in}{2.820661in}}%
\pgfpathlineto{\pgfqpoint{5.017771in}{2.730418in}}%
\pgfpathlineto{\pgfqpoint{5.027984in}{2.634929in}}%
\pgfpathlineto{\pgfqpoint{5.038197in}{2.512018in}}%
\pgfpathlineto{\pgfqpoint{5.052495in}{2.303088in}}%
\pgfpathlineto{\pgfqpoint{5.068836in}{2.070478in}}%
\pgfpathlineto{\pgfqpoint{5.077006in}{1.980256in}}%
\pgfpathlineto{\pgfqpoint{5.085177in}{1.914469in}}%
\pgfpathlineto{\pgfqpoint{5.097432in}{1.823972in}}%
\pgfpathlineto{\pgfqpoint{5.111730in}{1.692355in}}%
\pgfpathlineto{\pgfqpoint{5.130114in}{1.521388in}}%
\pgfpathlineto{\pgfqpoint{5.146454in}{1.377541in}}%
\pgfpathlineto{\pgfqpoint{5.160752in}{1.246629in}}%
\pgfpathlineto{\pgfqpoint{5.175050in}{1.145428in}}%
\pgfpathlineto{\pgfqpoint{5.187306in}{1.066120in}}%
\pgfpathlineto{\pgfqpoint{5.199561in}{1.002538in}}%
\pgfpathlineto{\pgfqpoint{5.213860in}{0.937828in}}%
\pgfpathlineto{\pgfqpoint{5.224072in}{0.901942in}}%
\pgfpathlineto{\pgfqpoint{5.256754in}{0.803745in}}%
\pgfpathlineto{\pgfqpoint{5.266967in}{0.781467in}}%
\pgfpathlineto{\pgfqpoint{5.275137in}{0.768667in}}%
\pgfpathlineto{\pgfqpoint{5.285350in}{0.757258in}}%
\pgfpathlineto{\pgfqpoint{5.295563in}{0.749534in}}%
\pgfpathlineto{\pgfqpoint{5.315989in}{0.736876in}}%
\pgfpathlineto{\pgfqpoint{5.342542in}{0.717101in}}%
\pgfpathlineto{\pgfqpoint{5.354798in}{0.711900in}}%
\pgfpathlineto{\pgfqpoint{5.373181in}{0.706721in}}%
\pgfpathlineto{\pgfqpoint{5.383394in}{0.706298in}}%
\pgfpathlineto{\pgfqpoint{5.397692in}{0.706779in}}%
\pgfpathlineto{\pgfqpoint{5.405863in}{0.704615in}}%
\pgfpathlineto{\pgfqpoint{5.420161in}{0.700173in}}%
\pgfpathlineto{\pgfqpoint{5.434459in}{0.698914in}}%
\pgfpathlineto{\pgfqpoint{5.503907in}{0.696458in}}%
\pgfpathlineto{\pgfqpoint{5.534545in}{0.696627in}}%
\pgfpathlineto{\pgfqpoint{5.534545in}{0.696627in}}%
\pgfusepath{stroke}%
\end{pgfscope}%
\begin{pgfscope}%
\pgfsetrectcap%
\pgfsetmiterjoin%
\pgfsetlinewidth{0.803000pt}%
\definecolor{currentstroke}{rgb}{0.000000,0.000000,0.000000}%
\pgfsetstrokecolor{currentstroke}%
\pgfsetdash{}{0pt}%
\pgfpathmoveto{\pgfqpoint{0.800000in}{0.528000in}}%
\pgfpathlineto{\pgfqpoint{0.800000in}{4.224000in}}%
\pgfusepath{stroke}%
\end{pgfscope}%
\begin{pgfscope}%
\pgfsetrectcap%
\pgfsetmiterjoin%
\pgfsetlinewidth{0.803000pt}%
\definecolor{currentstroke}{rgb}{0.000000,0.000000,0.000000}%
\pgfsetstrokecolor{currentstroke}%
\pgfsetdash{}{0pt}%
\pgfpathmoveto{\pgfqpoint{5.760000in}{0.528000in}}%
\pgfpathlineto{\pgfqpoint{5.760000in}{4.224000in}}%
\pgfusepath{stroke}%
\end{pgfscope}%
\begin{pgfscope}%
\pgfsetrectcap%
\pgfsetmiterjoin%
\pgfsetlinewidth{0.803000pt}%
\definecolor{currentstroke}{rgb}{0.000000,0.000000,0.000000}%
\pgfsetstrokecolor{currentstroke}%
\pgfsetdash{}{0pt}%
\pgfpathmoveto{\pgfqpoint{0.800000in}{0.528000in}}%
\pgfpathlineto{\pgfqpoint{5.760000in}{0.528000in}}%
\pgfusepath{stroke}%
\end{pgfscope}%
\begin{pgfscope}%
\pgfsetrectcap%
\pgfsetmiterjoin%
\pgfsetlinewidth{0.803000pt}%
\definecolor{currentstroke}{rgb}{0.000000,0.000000,0.000000}%
\pgfsetstrokecolor{currentstroke}%
\pgfsetdash{}{0pt}%
\pgfpathmoveto{\pgfqpoint{0.800000in}{4.224000in}}%
\pgfpathlineto{\pgfqpoint{5.760000in}{4.224000in}}%
\pgfusepath{stroke}%
\end{pgfscope}%
\begin{pgfscope}%
\pgfsetbuttcap%
\pgfsetmiterjoin%
\definecolor{currentfill}{rgb}{1.000000,1.000000,1.000000}%
\pgfsetfillcolor{currentfill}%
\pgfsetfillopacity{0.800000}%
\pgfsetlinewidth{1.003750pt}%
\definecolor{currentstroke}{rgb}{0.800000,0.800000,0.800000}%
\pgfsetstrokecolor{currentstroke}%
\pgfsetstrokeopacity{0.800000}%
\pgfsetdash{}{0pt}%
\pgfpathmoveto{\pgfqpoint{0.936111in}{2.926843in}}%
\pgfpathlineto{\pgfqpoint{1.991276in}{2.926843in}}%
\pgfpathquadraticcurveto{\pgfqpoint{2.030165in}{2.926843in}}{\pgfqpoint{2.030165in}{2.965732in}}%
\pgfpathlineto{\pgfqpoint{2.030165in}{4.087889in}}%
\pgfpathquadraticcurveto{\pgfqpoint{2.030165in}{4.126778in}}{\pgfqpoint{1.991276in}{4.126778in}}%
\pgfpathlineto{\pgfqpoint{0.936111in}{4.126778in}}%
\pgfpathquadraticcurveto{\pgfqpoint{0.897222in}{4.126778in}}{\pgfqpoint{0.897222in}{4.087889in}}%
\pgfpathlineto{\pgfqpoint{0.897222in}{2.965732in}}%
\pgfpathquadraticcurveto{\pgfqpoint{0.897222in}{2.926843in}}{\pgfqpoint{0.936111in}{2.926843in}}%
\pgfpathclose%
\pgfusepath{stroke,fill}%
\end{pgfscope}%
\begin{pgfscope}%
\definecolor{textcolor}{rgb}{0.000000,0.000000,0.000000}%
\pgfsetstrokecolor{textcolor}%
\pgfsetfillcolor{textcolor}%
\pgftext[x=1.121219in,y=3.901268in,left,base]{\color{textcolor}\rmfamily\fontsize{14.000000}{16.800000}\selectfont \(\displaystyle y\) (mm)}%
\end{pgfscope}%
\begin{pgfscope}%
\pgfsetrectcap%
\pgfsetroundjoin%
\pgfsetlinewidth{1.505625pt}%
\definecolor{currentstroke}{rgb}{0.121569,0.466667,0.705882}%
\pgfsetstrokecolor{currentstroke}%
\pgfsetdash{}{0pt}%
\pgfpathmoveto{\pgfqpoint{0.975000in}{3.683923in}}%
\pgfpathlineto{\pgfqpoint{1.363889in}{3.683923in}}%
\pgfusepath{stroke}%
\end{pgfscope}%
\begin{pgfscope}%
\definecolor{textcolor}{rgb}{0.000000,0.000000,0.000000}%
\pgfsetstrokecolor{textcolor}%
\pgfsetfillcolor{textcolor}%
\pgftext[x=1.519444in,y=3.615868in,left,base]{\color{textcolor}\rmfamily\fontsize{14.000000}{16.800000}\selectfont 0.2}%
\end{pgfscope}%
\begin{pgfscope}%
\pgfsetrectcap%
\pgfsetroundjoin%
\pgfsetlinewidth{1.505625pt}%
\definecolor{currentstroke}{rgb}{1.000000,0.498039,0.054902}%
\pgfsetstrokecolor{currentstroke}%
\pgfsetdash{}{0pt}%
\pgfpathmoveto{\pgfqpoint{0.975000in}{3.398523in}}%
\pgfpathlineto{\pgfqpoint{1.363889in}{3.398523in}}%
\pgfusepath{stroke}%
\end{pgfscope}%
\begin{pgfscope}%
\definecolor{textcolor}{rgb}{0.000000,0.000000,0.000000}%
\pgfsetstrokecolor{textcolor}%
\pgfsetfillcolor{textcolor}%
\pgftext[x=1.519444in,y=3.330467in,left,base]{\color{textcolor}\rmfamily\fontsize{14.000000}{16.800000}\selectfont 5.5}%
\end{pgfscope}%
\begin{pgfscope}%
\pgfsetrectcap%
\pgfsetroundjoin%
\pgfsetlinewidth{1.505625pt}%
\definecolor{currentstroke}{rgb}{0.172549,0.627451,0.172549}%
\pgfsetstrokecolor{currentstroke}%
\pgfsetdash{}{0pt}%
\pgfpathmoveto{\pgfqpoint{0.975000in}{3.113123in}}%
\pgfpathlineto{\pgfqpoint{1.363889in}{3.113123in}}%
\pgfusepath{stroke}%
\end{pgfscope}%
\begin{pgfscope}%
\definecolor{textcolor}{rgb}{0.000000,0.000000,0.000000}%
\pgfsetstrokecolor{textcolor}%
\pgfsetfillcolor{textcolor}%
\pgftext[x=1.519444in,y=3.045067in,left,base]{\color{textcolor}\rmfamily\fontsize{14.000000}{16.800000}\selectfont 20.0}%
\end{pgfscope}%
\end{pgfpicture}%
\makeatother%
\endgroup%
}
    \end{tabular}
    \caption{a. Comparison of custom written PDF function with filtering frequency $f_{LP}^{PDF}/ f_s= \num{5e-2}$ and PDF obtained from histogram function from Numpy library b. PDF displayed for the 3 locations.}

    \label{fig:PDF}
\end{figure}



\section{Discussion}

\section{Error Analysis}

\section{Conclusions}

\newpage
% References
\bibliographystyle{plain}
\bibliography{references}
\end{document}
